\documentclass[oneside]{article}

% Compile with latexmk -pdf, or rerun pdfLaTeX until references stabilize.
% Research revision: 15 October 2026.
% Revised proofs: 23 November 2026.
% Standalone source; reference list below.
% Constructive revision: full graded Segal/prism reconstruction, global
% named-cell normalization, explicit filtered relative cellular completion,
% marked NU universes, internal and relative higher classification,
% independent graded multikernels, balanced tree substitution,
% represented recovery, and consistently enumerated proofs.
\usepackage[utf8]{inputenc}
\usepackage{lmodern}
\usepackage{amsmath,amsthm,amsfonts,amssymb,mathtools,mathrsfs}
\usepackage[T1]{fontenc}
\usepackage{graphicx,array,booktabs,longtable}
\usepackage{microtype}
\usepackage[a4paper,total={6in,8in}]{geometry}
\usepackage{enumitem}
\usepackage{needspace}
% Keep the proof heading with the first numbered step.
\AddToHook{env/proof/before}{\Needspace{5\baselineskip}}
\usepackage{tikz-cd}
\usetikzlibrary{decorations.pathmorphing,arrows.meta}
\tikzcdset{every label/.append style={font=\small}}

\usepackage{xcolor}
\usepackage[colorlinks=true,linkcolor=black,citecolor=black,urlcolor=blue]{hyperref}

\setlength{\emergencystretch}{1.5em}

\theoremstyle{plain}
\newtheorem{theorem}{Theorem}[section]
\newtheorem{corollary}[theorem]{Corollary}
\newtheorem{lemma}[theorem]{Lemma}
\newtheorem{proposition}[theorem]{Proposition}
\theoremstyle{definition}
\newtheorem{definition}[theorem]{Definition}
\newtheorem{example}[theorem]{Example}
\newtheorem{judgement}[theorem]{Judgement}
\newtheorem{construction}[theorem]{Construction}
\newtheorem{problem}[theorem]{Problem}
\theoremstyle{remark}
\newtheorem{remark}[theorem]{Remark}
\newtheorem{notation}[theorem]{Notation}
\newtheorem{convention}[theorem]{Convention}
\newtheorem{observation}[theorem]{Observation}

\hypersetup{
  pdftitle={As Above, So Below: Theorization and the Infinity-Type-Theory Microcosm},
  pdfauthor={Conor Sheridan},
  pdfsubject={Marked lax limits and relative classifiers for microcosm principles}
}
\title{%
  As Above, So Below\\[0.5em]
  \normalfont\large Theorization and the $\infty$-Type-Theory Microcosm
}

\author{%
  Conor Sheridan\\
  \small\href{mailto:siridconor@gmail.com}{\texttt{siridconor@gmail.com}}%
}

\date{\today}

\newcommand{\Alg}{\operatorname{Alg}}
\newcommand{\B}{\mathbf B}
\newcommand{\CCyc}{\mathrm{Cyc}}
\newcommand{\CHom}{\mathbf{Hom}}
\newcommand{\CNU}{\mathcal C_{\mathrm{NU}}}
\newcommand{\CSeg}{\operatorname{Seg}}
\newcommand{\CSet}{\mathbf{Set}}
\newcommand{\Cact}{\mathrm{act}}
\newcommand{\Cat}{\mathbf{Cat}_{\infty}}
\newcommand{\Cel}{\mathrm{el}}
\newcommand{\Cim}{\operatorname{im}}
\newcommand{\Cint}{\mathrm{int}}
\newcommand{\CorepCat}{\operatorname{CorepCat}}
\newcommand{\CorepSpace}{\operatorname{CorepSpace}}
\newcommand{\Dep}{\operatorname{Dep}}
\newcommand{\Env}{\operatorname{Env}}
\newcommand{\Eq}{\operatorname{Eq}}
\newcommand{\Eqv}{\operatorname{Eqv}}
\newcommand{\Fin}{\operatorname{Fin}}
\newcommand{\Fun}{\operatorname{Fun}}
\newcommand{\GInfty}{\mathcal G_{\infty}}
\newcommand{\HasCorep}{\operatorname{HasCorep}}
\newcommand{\Homc}{\mathbf{Hom}}
\newcommand{\Ind}{\operatorname{Ind}}
\newcommand{\Int}{\operatorname{Int}}
\newcommand{\J}{J}
\newcommand{\LT}{\mathcal L_T}
\newcommand{\Lex}{\operatorname{Lex}}
\newcommand{\Map}{\operatorname{Map}}
\newcommand{\Mod}{\operatorname{Mod}}
\newcommand{\Moddem}{\operatorname{Mod}^{\mathrm{dem}}}
\newcommand{\NU}{\mathrm{NU}}
\newcommand{\Nat}{\operatorname{Nat}}
\newcommand{\Nn}[1]{N^{\mathrm R}_{#1}}
\newcommand{\Om}{\boldsymbol\Omega}
\newcommand{\RFib}[1]{\operatorname{RFib}(#1)}
\newcommand{\RNU}{\mathcal R_{\infty}}
\newcommand{\Sect}{\operatorname{Sect}}
\newcommand{\Sps}{\mathcal S}
\newcommand{\TT}{\mathbf{TT}}
\newcommand{\Th}{\operatorname{Th}}
\newcommand{\Thz}{\operatorname{Thz}}
\newcommand{\cK}{\mathcal K}
\newcommand{\coc}{\mathrm{cocart}}
\newcommand{\colim}{\operatorname*{colim}}
\newcommand{\core}{\simeq}
\newcommand{\disc}{\operatorname{disc}}
\newcommand{\eT}{e_T}
\newcommand{\ff}{\mathrm{fl}}
\newcommand{\fib}{\operatorname{fib}}
\newcommand{\id}{\operatorname{id}}
\newcommand{\isEq}{\operatorname{isEquiv}}
\newcommand{\laxlim}{\mathop{\mathrm{laxlim}}}
\newcommand{\minrep}{\mathrm{min}}
\newcommand{\one}{\mathbf 1}
\newcommand{\op}{\mathrm{op}}
\newcommand{\yo}{\mathrm y}

\begin{document}
\maketitle
\begin{abstract}
The microcosm principle asks how a structure can occur inside a categorified
structure of the same kind. For dependent structures, this requires an
account of operations, substitution, and the comparisons that must remain
invertible. Starting from the Nguyen--Uemura framework, we construct
theorization by normalized operation kernels and finite cocartesian
exactness. Represented theories recover left-exact categorical models;
maps recover partially strong lax transformations. We prove functoriality
of the kernels, their marked dependent operations, and their universal
envelopes. A marked comma construction followed by left-exact reflection
produces the walking internal theory. Its higher comparison objects give
 the stratified NU microcosm. A limit-closed marked-cocartesian ambient
 category extends categorical stratification and intrinsic internalization
 to unrepresented operation theories. Categories of elements give a relative
 NU classifier for every bounded marked lax limit. A filtered relative
 cellular resolution, starting with the finite named-cell shadow, recovers
 the entire completed marked classifier and all its extension categories.
 Independent higher composition operations are organized as normalized
 multikernels. Balanced tree substitution constructs their composites;
 corepresented rows recover oplax composition, and substitution-stable
 universal operations recover coherent associative composition. The
 construction retains the full base-degree data and recovers lax maps
 between represented systems.
\end{abstract}

\tableofcontents

\section{Introduction}\label{sec:introduction}

A monoid can be defined in a monoidal category because the ambient category
already supplies places for its multiplication, unit, and equations to
live. The same question becomes less transparent for a dependent
structure. An operation may require a previously chosen type or term;
substitution changes that parameter; and some comparison maps express
structure while others express coherence. Before asking for an internal
instance of such a structure, one needs an ambient language in which all
these dependencies can be stated.

The Nguyen--Uemura framework provides that language. Its universal doctrine
$\RNU$ has
\[
 \Lex(\RNU,\Sps)\simeq\TT,
 \qquad
 \Lex(\RNU,\Cat)\simeq\operatorname{Cat}(\TT),
\]
so passing from spaces to categories gives a categorical environment for
$\infty$-type theories. This answers the first question: what should a
categorified ambient theory be? It leads immediately to a second one.
Must every operation act by a chosen transport functor, or can the space
of operations be retained before such a representative exists?

For a categorical model $K:T\to\Cat$, an operation of degree $f:s\to t$
from $x$ to $y$ is a map $K(f)x\to y$. The formula uses transport only to
represent the functor of possible outputs. Keeping that functor, rather
than requiring its representative, leads to the normalized kernels
\[
 P_f(x,y)=
 \fib_f\bigl(\Map_E(x,y)\longrightarrow\Map_T(px,py)\bigr)
\]
of a category $p:E\to T$. Composition in $E$ supplies substitution of
operations. A structural marking records the directions in which
transport must still be represented. Finite cocartesian exactness retains
the finite-limit content of the original doctrine. The recovery theorem
then shows that this enlargement has the intended boundary: on the
represented locus it returns precisely the left-exact categorical models,
and its maps return the appropriate partially strong lax transformations.

This is the form of theorization studied here. The passage from lax maps
to maps of theorized structures follows the principle developed by
Matsuoka~\cite{Matsuoka2017}. Our construction makes its finite-limit and
dependent content explicit in the NU setting. In particular, the kernels
are functorial, their marked-square comparisons are functorial, and the
graded reconstruction retains the entire operation category and its
transformations.

Once operations are available, an internal theory is a coherent choice
of objects and operations over the doctrine, strong in the marked
directions. It is therefore a marked section. The next question is whether
all these choices can be represented by one universal source. The comma
category records a source object together with transport to a test arity;
localization imposes the selected invertibilities; left-exact reflection
places the result back in the categorical model category. For
$a:D^W\to T$ this constructs
\[
 \Omega_{a,W}
 =L_{\mathrm{lex}}\bigl(t\longmapsto(a\downarrow t)[W_t^{-1}]\bigr),
 \qquad
 \Nat(\Omega_{a,W},K)\simeq\laxlim_D^W(Ka).
\]
The source contains the laxity data, so its classifying maps are ordinary
natural transformations. Their morphisms remain arbitrary modifications.
For the positively marked NU doctrine the result is the walking internal
$\infty$-type theory.

The microcosm also has a relative form: a marked source need not itself
be the NU doctrine. To compare its microcosm with the universal
framework, one must retain the doctrine
over which its operations are interpreted. Categories of elements provide
an enriched fully faithful realization
\[
 \J_T:\Lex(T,\Cat)\longrightarrow
 \operatorname{Cat}(\TT)/\B_T,
 \qquad \B_T=\disc(T_{\minrep}).
\]
Combining this realization with the envelope gives
\begin{equation}\label{eq:intro-relative}
 \laxlim_D^W(Ka)\simeq
 \Homc_{\operatorname{Cat}(\TT)/\B_T}
       (\J_T\Omega_{a,W},\J_TK).
\end{equation}
Free finite-limit completion extends the construction to every bounded
marked diagram. The same envelope admits an explicit cellular computation:
starting with any marked source $D^W\to T^\tau$, finite collections of
simplices of $T$ construct relative pushout stages whose envelope colimit
is $\Omega_{T,\tau}$, including all the extension data missing from the
initial source. The base in the relative formula has mathematical content:
an absolute replacement cannot have the same universal property in general.

The construction can then be iterated. To allow independent composition operations, we retain the
categories of lower operations and assign a multikernel to each input
string and output. A substitution product balances over intermediate
operations. Its associative algebras provide coherent grafting, and an
explicit tree formula computes composites of every finite height.
Representing the rows recovers composition functors and their oplax
comparisons. When universal operations remain universal after
substitution, those comparisons become equivalences. This gives a
construction and a recovery test for the new composition data.

We also construct a limit-closed category of marked-cocartesian objects;
its complete categorical objects admit unrepresented operation kernels
in every external degree and an intrinsic higher internalization functor.
The walking source also explains comparisons between internal theories.
Its categorical copowers represent chains of morphisms, and free complete
category objects represent further categorical directions. This gives the
stratified microcosm, with the original structural marking unchanged.

The paper follows two successive passages from theorization to a
microcosm. Section~\ref{sec:nu-foundations} establishes the NU doctrine
and its dependent arities. Section~\ref{sec:theorization} constructs
operation theories, their kernels, recovery, and functoriality.
Section~\ref{sec:microcosm} constructs internal theories and their
walking classifier. Section~\ref{sec:relative} explains the need for a
relative base and proves universal classification, including the
cellular completion of marked sources. Section~\ref{sec:higher-composition}
then theorizes higher composition through graded multikernels and develops
the unrepresented higher categorical operation structures.
Section~\ref{sec:stratified-microcosm} constructs the iterated NU
microcosm, its higher comparison objects, and its stratified relative
classification.

\subsection{Conventions}\label{subsec:conventions}

Categories, limits, slices, and functor categories are $\infty$-categorical
unless an ordinary construction is specified. Ordinary internal categories and complete Segal objects are distinguished
whenever both occur.
Smallness is relative to a fixed hierarchy of universes. Each source is
bounded in one universe; presheaves, localizations, and classifiers are
formed in a sufficiently large subsequent universe. ``Bounded small''
always refers to this convention.

Write $\Sps$ for spaces and $\Cat$ for $\infty$-categories. The symbol
$\Map$ denotes a mapping space, whereas $\Nat(A,K)$ denotes the full
$\infty$-category of ordinary natural transformations and modifications.
Its objects have invertible naturality comparisons; its morphisms need not
be invertible. The underlying mapping space is its core. We use the complete
Segal nerve $N_nC=\Fun([n],C)^\core$.

A \textbf{marking} $W$ contains all equivalences, is invariant under equivalence
of arrows, and is closed under composition. No two-out-of-three condition is
included. The sharp marking marks every arrow. Our lax-cone convention for
$C:D\to\Cat$ is
\[
 \ell_u:C(u)x_d\longrightarrow x_{d'},\qquad u:d\to d'.
\]
The marked lax limit retains cones for which $\ell_u$ is invertible when
$u\in W$, and retains every comparison between such cones. All displayed
commutative diagrams include their specified coherent homotopies.

For orientation, take $D=[1]$ and a diagram $F:C_0\to C_1$.
A lax cone is $(x_0,x_1,\alpha:F(x_0)\to x_1)$. Its morphisms are
pairs $(u_0,u_1)$ with the coherent square
\[
\begin{tikzcd}[column sep=large]
 F(x_0)\ar[r,"\alpha"]\ar[d,"F(u_0)"']&x_1\ar[d,"u_1"]\\
 F(x'_0)\ar[r,"\alpha'"']&x'_1.
\end{tikzcd}
\]
Marking $0\to1$ requires $\alpha$ to be invertible. Evaluation at $0$
then identifies the marked-section category with $C_0$, with inverse
$x\mapsto(x,F(x),\id)$. For the constant diagram at $[1]$ with identity transition, the
unmarked-arrow section category has three objects in a chain
(the arrows $00,01,11$); marking the transition gives $[1]$.
Thus even for a fixed diagram the choice of marking changes the answer.



\begin{center}
\begin{tabular}{@{}ll@{}}
\toprule
Notation & Role \\
\midrule
$T_{\minrep}$ & Equivalences are the representable arrows \\
$D^W$ & Native arities with their structural marking \\
$T^\sharp$ & Every arrow is marked \\
$\langle a(W)\rangle_T$ & Marking generated by the displayed images \\
$\Omega_{a,W}$ & Left-exact marked envelope \\
$\B_T=\disc(T_{\minrep})$ & Relative base for the NU realization \\
\bottomrule
\end{tabular}
\end{center}
All naturality assertions concern the displayed functors and full hom
categories. When a construction is defined by a universal property, its
choices and coherence are those of the associated contractible space of
universal data, as made precise in Lemma~\ref{lem:universal-choices}.

The general classification takes its marked source, indexing doctrine
and target diagram as inputs. For the universal NU source,
Proposition~\ref{prop:nu-cells} constructs the doctrine and
Definition~\ref{def:nu-marking} constructs its marking. Its operation
squares, coherence and simplicial resolution are derived below.
Representability and stability of independently given operations are
recognition conditions. Compact generation is an input only to the
general semantic recognition statement; Proposition~\ref{prop:stratified-compact-doctrines}
derives it at the successive NU stages.

\section{The Nguyen--Uemura framework}
\label{sec:nu-foundations}

A dependent operation cannot be described by its number of inputs alone.
Its later inputs may be types or terms over its earlier ones. Finite
limits express the contexts in which those inputs can be assembled, and
dependent products express the corresponding families of operations.
We therefore begin with the finite-limit and representable-map structure
of an $\infty$-type theory.

The theory--model equivalence of Nguyen and Uemura~\cite{NguyenUemura2025}
then identifies the space-valued models of one universal doctrine with
theories themselves. Applying complete category objects gives the
categorical ambient structures needed for a microcosm. The comprehension
fiber calculation between these two steps explains exactly where the
dependent arities enter: an operation is a point in a fiber over its
previously chosen parameters. These are the same fibers that the
normalized kernels will retain in Section~\ref{sec:theorization}.

\subsection{\texorpdfstring{$\infty$-type theories and their models}{Infinity-type theories and their models}}
\label{subsec:infty-type-theory}

\begin{definition}
\label{def:infty-type-theory}
    An $\infty$-type theory is a small finitely complete $\infty$-category $T$ with a replete class of \textbf{representable} arrows containing identities, closed under composition and pullback, and consisting of exponentiable arrows. Exponentiability of $u:y\to z$ means that the pullback functor has a right adjoint: $\Sigma_u\dashv u^*\dashv u_*: T_{/y}\rightleftarrows T_{/z}.$ Here $\Sigma_u$ is postcomposition, and $u_*$ is the dependent product. A theory morphism $H:T\to T'$ preserves finite limits, representable arrows, and dependent products along them. Write $\TT$ for their category.
\end{definition}

The class of representable arrows is an input for a general theory.
On the same category of finite sets, equivalences alone and all maps
both satisfy the definition: for the latter, dependent product along
$f$ takes the product of each finite family over the fibers of $f$.
These two classes already differ on $\varnothing\to 1$.

Dependent-product preservation means invertibility of the canonical mate
\begin{equation}
\label{eq:pushforward-mate}
     H_z u_*\longrightarrow(Hu)_*H_y.
\end{equation}
It is obtained by applying the unit of $(Hu)^*\dashv(Hu)_*$, the inverse of the pullback comparison $H_yu^*\simeq(Hu)^*H_z$, and the counit of $u^*\dashv u_*$.  The pullback comparison and its right mate have the respective shapes
\[
    \begin{tikzcd}[column sep=large]
        T_{/z}\ar[r,"u^*"]\ar[d,"H_z"']&
        T_{/y}\ar[d,"H_y"]\\
        T'_{/Hz}\ar[r,"(Hu)^*"']&T'_{/Hy}
    \end{tikzcd}
    \qquad
    \begin{tikzcd}[column sep=large]
        T_{/y}\ar[r,"u_*"]\ar[d,"H_y"']&
        T_{/z}\ar[d,"H_z"]\\
        T'_{/Hy}\ar[r,"(Hu)_*"']&T'_{/Hz}.
    \end{tikzcd}
\]
The first square is canonically commutative for a left-exact $H$. For a theory morphism the mate makes the second square commutative; the direction of the mate is \eqref{eq:pushforward-mate}.


\begin{lemma}[Universal choices]
\label{lem:universal-choices}
    For fixed input, the spaces of terminal objects, limiting cones, presheaf representations, and right adjoints with adjunction structure are empty or contractible. An iterated finite construction with nonempty contractible extension spaces has a contractible space of coherent realizations.
\end{lemma}
\begin{proof}
    \leavevmode
    \begin{enumerate}
        \item \textbf{Reduce universal data to terminal objects.} Limits and representations are terminal objects in the corresponding cone or representation categories. Any two terminal objects have a contractible space of equivalences compatible with their universal maps. Adjoints with unit and counit are likewise universal representations of the appropriate mapping-space families.
    
        \item \textbf{Totalize the successive extension spaces.} For the last assertion, write each successive space as the dependent total space of its extension family over the preceding space. A map with contractible homotopy fibers is an equivalence; induction therefore identifies every stage with the initial point.\qedhere
    \end{enumerate}
\end{proof}

The representable arrows are structure on a theory, and need not be all exponentiable arrows. Every finitely complete $T$ has a minimal structure $T_{\minrep}$ with just equivalences representable. An $\infty$-topos has the maximal structure, with all arrows representable. The category $\TT$ is presentable and admits free constructions on marked categories~\cite[\S3.1]{NguyenUemura2025}. In particular, the free theory $\GInfty$ on a representable arrow $\partial:E\to U$ satisfies
\[
    \Map_{\TT}(\GInfty,T)\simeq(\text{representable arrows of }T)^\core.
\]

\begin{definition}
\label{def:theory-model}
    A \textbf{$T$-theory} is a left exact functor $T\to\Sps$, and $\Th(T)=\Lex(T,\Sps).$ For a small category $\mathcal C$, write $\RFib{\mathcal C}\simeq\Fun(\mathcal C^{\op},\Sps)$ for its category of right fibrations. A map of right fibrations is representable when every pullback along a representable presheaf is representable. A \textbf{$T$-model} consists of a small category $\mathcal C$, called a \textbf{category of contexts}, with terminal object and a theory morphism $M:T\to\RFib{\mathcal C}$.
\end{definition}

If $f:B\to A$ is a representable map of right fibrations and $a:\yo\Gamma\to A$, its pullback is $\yo\{a\}_f$ for an object of $\mathcal C$, through a contractible space of representations. Equivalently, the functor between total categories $f$ has a right adjoint; its counit constructs this \textbf{comprehension object} and the square
\[
    \begin{tikzcd}
        \yo\{a\}_f \ar[r]\ar[d] & B\ar[d,"f"]\\
        \yo\Gamma\ar[r,"a"']&A.
    \end{tikzcd}
\]
To see the equivalence, the fiber of the presheaf pullback at $\Delta$ is the space of arrows over a map $\Delta\to\Gamma$ into the object $a$. Representing that presheaf is precisely the universal-arrow property for the right adjoint at $a$. Its naturality supplies substitution and all higher coherences. This is the representability criterion of~\cite[\S3.2]{NguyenUemura2025}.

For a model $M$ and a representable $u:y\to z$, an element $a\in M(z)_\Gamma$ therefore gives a \textbf{context extension $\{a\}_u$}. The model's finite-limit and pushforward preservation supply the pairing and substitution operations. A model morphism has a base functor preserving the terminal context and compatible maps of right fibrations; the Beck--Chevalley comparisons for representable arrows are invertible. In particular it preserves comprehension. We use the category of models and these morphisms from~\cite[\S3.3]{NguyenUemura2025}.

\begin{definition}
\label{def:internal-model-democratic}
    The \textbf{internal language of a model} is $L(M)(x)=M(x)_{1_{\mathcal C}}$. A context is \textbf{contextual} if it is obtained, up to equivalence, by a finite tower of comprehension from the terminal context. A model is \textbf{democratic} if every context is contextual.
\end{definition}

\begin{theorem}[Nguyen--Uemura]\label{thm:democracy}
    Internal language induces an equivalence $L:\Moddem(T)\simeq\Th(T).$ Moreover $\Th(T)\simeq\operatorname{Ind}(T^{\op})$: every $T$-theory is a filtered colimit of representable theories $T(x,-)$.
\end{theorem}
\begin{proof}[Recall of the construction]
    \leavevmode
    \begin{enumerate}
        \item \textbf{Construct the democratic model.} For a representable theory $T(x,-)$, the slice theory $T_{/x}$ freely adjoins a global element of $x$. Its contextual Yoneda model supplies the corresponding democratic model. The general syntactic model is obtained by the left adjoint to internal language. The unit is an equivalence, first on representables and then on their filtered colimits. On democratic models, the comprehension fiber formula below proves that internal language is conservative by induction on contextual towers. The counit is consequently an equivalence there. These constructions and their adjunction are proved in~\cite[\S4, Thm.~4.14]{NguyenUemura2025}.
        
        \item \textbf{Identify the filtered completion.} The final assertion is the Ind characterization proved in Theorem~\ref{thm:Th-filtered-colimit-generation}.\qedhere
    \end{enumerate}
\end{proof}

\subsection{Polynomial arities and comprehension}
\label{subsec:polynomials}

The dependence of an input on a previous parameter should be visible in
the arity itself. A representable map provides the family of possible
inputs, and its polynomial functor packages the corresponding dependent
operation. The fiber formula below identifies its semantic value with the
ordinary comprehension fiber; it will later let us read a dependent
operation directly from a normalized kernel.

For an exponentiable $u:y\to z$ define
\[
    P_u=z_!u_*y^*:T\longrightarrow T,
    \qquad \varpi_{u,x}:P_ux\longrightarrow z.
\]
Here $y^*x=(y\times x\to y)$ and $z_!$ forgets the map to $z$. The counit gives evaluation $(P_ux)\times_z y\to x$. The construction is functorial in $x$ and is preserved by theory morphisms through their canonical mates.

\begin{lemma}[Comprehension fiber formula]
\label{lem:comprehension-fiber}
    For a model $M$, a representable $u:y\to z$, a context $\Gamma$, and $a\in M(z)_\Gamma$, there is a natural equivalence $M(x)_{\{a\}_u}\simeq \fib_a\bigl(M(P_ux)_\Gamma\to M(z)_\Gamma\bigr).$ It is compatible with model morphisms and substitution.
\end{lemma}
\begin{proof}
    \leavevmode
    \begin{enumerate}
        \item \textbf{Compute the comprehension fiber by adjunction.} Put $\mathcal P=\RFib{\mathcal C}$ and $f=M(u)$. Regard $a$ as the slice object $\yo\Gamma\to M(z)$. Counit comprehension represents $f^*a$ by $\yo\{a\}_u$. There are natural equivalences
        \begin{align*}
             M(x)_{\{a\}_u}
             &\simeq\Map_{\mathcal P}(\yo\{a\}_u,M(x))\\
             &\simeq\Map_{\mathcal P_{/M(y)}}(f^*a,M(y)\times M(x))\\
             &\simeq\Map_{\mathcal P_{/M(z)}}(a,f_*M(y^*x))\\
             &\simeq\Map_{\mathcal P_{/M(z)}}(a,M(u_*y^*x)).
        \end{align*}
        The second line uses the fixed map $f^*a\to M(y)$; the third uses the dependent-product adjunction and finite-limit preservation; the fourth is the model pushforward mate. Yoneda identifies the last line with the stated fiber.
    
        \item \textbf{Verify naturality and substitution.} A model morphism's Beck--Chevalley comparison identifies the comprehension counits, and the mate identity identifies the two displayed adjunctions. This proves compatibility with model morphisms. Pullback of the same diagram proves compatibility with substitution.\qedhere
    \end{enumerate}
\end{proof}

The formula also explains contextual induction: an equivalence on closed realizations is an equivalence at the terminal context; its value at an extension is a fiber of values at the preceding context. Iterating proves equivalence on every contextual fiber.

\begin{lemma}[Polynomial base change]
\label{lem:polynomial-bc}
    For a pullback square of representable arrows
    \[
        \begin{tikzcd}
            y'\ar[r,"a"]\ar[d,"u'"']&y\ar[d,"u"]\\
            z'\ar[r,"b"']&z
        \end{tikzcd}
    \]
    there is a canonical equivalence $b^*u_*\simeq u'_*a^*$. It induces cartesian base-change comparisons for polynomial objects and their evaluation maps. These comparisons are coherent under pasting.
\end{lemma}
\begin{proof}
    \leavevmode
    \begin{enumerate}
        \item \textbf{Construct the Beck--Chevalley mate.} For $A\to z'$ and $B\to y$, the adjunctions and the pullback square give
        \begin{align*}
            \Map_{/z'}(A,b^*u_*B)
            &\simeq\Map_{/z}(\Sigma_bA,u_*B)\\
            &\simeq\Map_{/y}(u^*\Sigma_bA,B)\\
            &\simeq\Map_{/y}(\Sigma_a u'^*A,B)\\
            &\simeq\Map_{/y'}(u'^*A,a^*B)\\
            &\simeq\Map_{/z'}(A,u'_*a^*B).
        \end{align*}
        Yoneda gives the stated mate.

        \item \textbf{Identify polynomial evaluation and its coherence.} For polynomial objects take $B=y\times x\to y$. The resulting equivalence identifies $z'\times_zP_ux$ with $P_{u'}x$, including evaluation, since evaluation is the counit used in the same calculation. A map $x'\to x$ then gives the corresponding comparison with changed polynomial argument. Pasting uses composition of these natural mapping-space equivalences, so all higher coherences are inherited from the adjunctions.\qedhere
    \end{enumerate}
\end{proof}

\subsection{The universal theory and categorification}
\label{subsec:universal-theory}

The preceding constructions apply to an arbitrary theory. To make theories
themselves the objects of a microcosm, we now use the universal NU
doctrine. The passage to categories requires no new object-level syntax:
complete category objects can be formed degreewise in its space-valued
models.

The theory $\RNU$ is the NU universal doctrine for $\infty$-type theories ~\cite[Def.~5.27 and \S5]{NguyenUemura2025}. It has a univalent left exact universe $\partial:E\to U$, a monomorphism $r:R\to U$ selecting the representable types, terminal and dependent-sum structures on its pullback $\partial_R:E_R\to R$, and the mixed dependent-product structure. Section~\ref{subsec:nu-doctrine} gives the operation squares explicitly. The selected class of a classified theory is recovered from $R$, so the classification concerns theories with their representable-map structure.

\begin{theorem}[NU classification and its categorical form]
\label{thm:nu-categorical}
    There are natural equivalences
    \[
        \Th(\RNU)\simeq\TT,
        \qquad
        \Lex(\RNU,\Cat)\simeq\operatorname{Cat}(\TT).
    \]
    More generally, $\Lex(T,\Cat)\simeq\operatorname{Cat}(\Th(T))$ for every small finitely complete $T$.
\end{theorem}
\begin{proof}
    \leavevmode
    \begin{enumerate}
        \item \textbf{Apply the theory--model equivalences.} The first equivalence is the composite of~\cite[Thms.~4.14 and~5.29]{NguyenUemura2025}.
        
        \item \textbf{Pass to complete Segal objects.} For the general categorical assertion, the complete Segal nerve identifies a category-valued functor with a simplicial space-valued diagram satisfying the Segal and completeness conditions pointwise. Since the nerve preserves limits and detects equivalences, left exactness of that functor is equivalent to left exactness at every simplicial level. Limits in $\Th(T)$ are pointwise. Thus the Segal and completeness conditions are precisely those of a complete category object in $\Th(T)$. This gives the general equivalence and, with the first equivalence, the categorical NU assertion.\qedhere
    \end{enumerate}
\end{proof}

\begin{proposition}[Categorical dependent arity]
\label{prop:polynomial-arity}
    Let $K\in\Lex(T,\Cat)$, let $u:y\to z$ be representable, and let $A\in K(z)$. Define $\Dep^K_{u,x}(A)=K(P_ux)\times_{K(z)}\{A\}.$ If $M_n$ is the democratic model for $K_n=\Fun([n],K(-))^\core$ and $A_n\in K_n(z)$ is the constant $n$-simplex at $A$, then $N_n\Dep^K_{u,x}(A)\simeq M_n(x)_{\{A_n\}_u}.$ The equivalence is natural in $x$, simplicial operators, and ordinary maps of categorical models.
\end{proposition}
\begin{proof}
\leavevmode
    \begin{enumerate}
        \item \textbf{Compute the fiber levelwise.} The nerve preserves the defining fiber, whose $n$th level is $\fib_{A_n}(K_n(P_ux)\to K_n(z))$. Apply Lemma~\ref{lem:comprehension-fiber} at the terminal context in $M_n$. Its naturality in model morphisms gives compatibility with every simplicial operator and with maps of categorical models.\qedhere
    \end{enumerate}
\end{proof}

Thus $K(P_ux)$ is a category of dependent families over $K(z)$. Its fiber is determined by comprehension, rather than by replacing a dependent family with an ordinary external function space.


\subsection{Filtered generation and coherent universal operations}
\label{subsec:filtered-generation-universal-operations}

Two facts will let us pass between this semantic description and explicit
operations. Representable models generate the theory category, so a
construction can be checked on arities. Universal choices, when they
exist, have contractible spaces of realizations, so those checks assemble
over parameters without a separate choice at each object.

\begin{theorem}[Filtered generation]
\label{thm:Th-filtered-colimit-generation}
    For a small finitely complete $T$, there is a natural equivalence $\Th(T)\simeq\operatorname{Ind}(T^{\op})$. The representable theories $h_t=T(t,-)$ are compact and generate under filtered colimits.
\end{theorem}
\begin{proof}
    \leavevmode
    \begin{enumerate}
        \item \textbf{Apply the Ind characterization.} A functor $T\to\Sps$ is a presheaf on $T^{\op}$, and left exactness says that it carries finite colimits in $T^{\op}$ to limits. The Ind characterization identifies exactly these presheaves with $\operatorname{Ind}(T^{\op})$; see~\cite[\S5.3.5]{Lurie2009} and~\cite[Definition~3.13]{NguyenUemura2025}. Evaluation computes maps from a representable and preserves filtered colimits, proving compactness.\qedhere
    \end{enumerate}
\end{proof}

\begin{proposition}[Coherent dependent-product data]\label{prop:coherent-dependent-data}
    Over the space of representable arrows of an $\infty$-type theory, the right adjoints to pullback and their units and counits assemble coherently. For a left exact functor preserving representable arrows, the comparison of dependent products is the canonical mate \eqref{eq:pushforward-mate}. Its invertibility is a property of that functor.
\end{proposition}
\begin{proof}
    \leavevmode
    \begin{enumerate}
        \item \textbf{Assemble the adjunction family.} Form the family of adjunction records over the space of representable arrows. Its fibers are nonempty by exponentiability and contractible by Lemma~\ref{lem:universal-choices}; its projection is therefore an equivalence. An inverse constructs the coherent adjunction family.
    
        \item \textbf{Construct the coherent mates.} The same argument for the pullback comparison, followed by the unit and counit, gives the mate. Naturality of these operations gives all its coherences, and the equivalence test is proposition-valued.\qedhere
    \end{enumerate}
\end{proof}

\begin{lemma}[The complete-nerve comparison test]
\label{lem:complete-nerve-comparison}
    For categorical diagrams $X,K:T\to\Cat$, put $K^{[n]}(t)=\Fun([n],K(t))$. There are natural equivalences
    \[
        N_n\Nat(X,K)\simeq
        \Map_{\Fun(T,\Cat)}(X,K^{[n]}).
    \]
    Consequently a comparison of hom $\infty$-categories is an equivalence whenever its mapping-space comparison, after the corresponding cotensor by every $[n]$, is an equivalence and is natural in $[n]$. The same test applies in a homotopy slice. Left-exact diagrams are closed under these pointwise cotensors.
\end{lemma}
\begin{proof}
    \leavevmode
    \begin{enumerate}
        \item \textbf{Compare complete Segal nerves.} By definition, the left side is $\Fun([n],\Nat(X,K))^\core$. Currying identifies this with the mapping space on the right. Naturality makes the comparisons into a map of complete Segal nerves, and the complete Segal nerve is fully faithful. In a slice the same argument takes the fiber of the relevant mapping diagram over the specified base map. Finally $\Fun([n],-)$ is a right adjoint and therefore preserves finite limits, proving the last claim.\qedhere
    \end{enumerate}
\end{proof}

\section{Theorization: from categorical models to operation theories}\label{sec:theorization}
A categorical model $K:T\to\Cat$ assigns to each operation degree
$f:s\to t$ a transport functor $K(f)$. For $x\in K(s)$, the possible
operations with output $y\in K(t)$ form
$\Map_{K(t)}(K(f)x,y)$. The transport object $K(f)x$ represents this
functor of $y$. Theorization starts by retaining the functor of operations
when such a representative is not available.

A category over $T$ already contains the required degree information and
composition. Its fibers are the categories of objects of each arity, and
its mapping spaces over a fixed arrow are the normalized operation
kernels. Cocartesian arrows identify the operations that do have universal
representatives. This gives a direct way to retain represented transport
in structural directions while allowing general operation kernels
elsewhere, in the spirit of Matsuoka's theorization~\cite{Matsuoka2017}.

The remaining issue is exactness. Requiring arbitrary limits of the total
category would impose conditions on operations that the original
categorical model does not supply. Instead, finite cocartesian exactness
asks for the universal cones whose diagrams and legs express transport.
The recovery theorem proves that this condition is exactly left
exactness on represented objects. We then identify the operation kernels,
their dependent comparisons, their reconstruction, and their functoriality.

Throughout, $T$ is small and finitely complete, $\tau$ is a marking, and
$\mathcal L_T=\Lex(T,\Cat)$. All universal transport data are retained
through their contractible spaces of choices.

\subsection{Finite cocartesian exactness}\label{sec:thz:recovery}

For $p:E\to T$, an arrow $u:x\to y$ over $f:s\to t$ is
$p$-cocartesian when, for every $z\in E$, the map
\begin{equation}\label{thz:eq:cocart}
 \Map_E(y,z)\longrightarrow
 \Map_E(x,z)\times_{\Map_T(s,pz)}\Map_T(t,pz)
\end{equation}
is an equivalence.  This mapping-space condition is invariant under
categorical equivalence and makes sense for an arbitrary functor of
$\infty$-categories.  We use the invariant notion of cocartesian fibration:
a functor is cocartesian when every base arrow admits a cocartesian lift from
every object over its source.  When an explicit quasicategorical
straightening is required, one may replace the functor by an equivalent
categorical fibration over the same base; this replacement is a proof device
and is not retained as structure.\cite{HHLN2023}

\begin{lemma}[Cocartesian calculus]\label{thz:lem:cocart-calculus}
Cocartesian arrows are invariant under equivalence of arrows and are closed
under composition.  If $u$ and $vu$ are cocartesian, then $v$ is
cocartesian.  A cocartesian arrow over an equivalence is an equivalence.
For a fixed source $x$ and degree $f:px\to t$, the space
$\operatorname{CocLift}_p(f,x)$ of cocartesian lifts of $f$ from $x$,
including their target objects, is either empty or contractible.
\end{lemma}
\begin{proof}
\leavevmode
\begin{enumerate}
\item \textbf{Paste the mapping-space pullbacks.} For composable $u:x\to y$ and $v:y\to w$, the mapping-space square for
$vu$ is the pasting of the squares for $u$ and $v$.  Pasting and
cancellation of pullback squares prove closure under composition and the
cancellation assertion.  Equivalence invariance follows by conjugating
\eqref{thz:eq:cocart} by equivalences of its source and target mapping
spaces.
\begin{samepage}
The two squares and their outer rectangle are the comparisons
for $v$, $u$, and $vu$, respectively:
\[
\begin{tikzcd}[column sep=2em,row sep=large]
 \Map_E(w,z)\ar[r,"v^*"]\ar[d]&
 \Map_E(y,z)\ar[r,"u^*"]\ar[d]&\Map_E(x,z)\ar[d]\\
 \Map_T(pw,pz)\ar[r,"p(v)^*"']&
 \Map_T(py,pz)\ar[r,"p(u)^*"']&\Map_T(px,pz).
\end{tikzcd}
\]
\end{samepage}
\item \textbf{Detect equivalences over equivalences.} Let $u:x\to y$ be
$p$-cocartesian with $p(u)$ an equivalence. For every $z\in E$,
precomposition with $p(u)$ is an equivalence on base mapping spaces.
Its pullback identifies the right-hand side of \eqref{thz:eq:cocart}
with $\Map_E(x,z)$. Under this identification the cocartesian comparison
is precomposition with $u$, so $\Map_E(y,z)\to\Map_E(x,z)$ is an
equivalence for every $z$. Yoneda now makes $u$ an equivalence.
Conversely, if $u$ is an equivalence, both horizontal maps in its
mapping-space square are equivalences. The square is therefore a
pullback, so $u$ is $p$-cocartesian.
\item \textbf{Compare cocartesian lifts.} Let $u:x\to y$ and $u':x\to y'$ be two cocartesian lifts of the same
$f:s\to t$.  The cocartesian universal property of $u$, evaluated at the
point $u'$ and over $\id_t$, gives a contractible space of
factorizations
\[
 u'\simeq vu,\qquad v:y\to y',\qquad p(v)=\id_t.
\]
Interchanging $u$ and $u'$ gives a reverse factorization.  Applying the
same contractible-factorization statement to the two composites identifies
them with the respective identities.  Thus the comparison is an
equivalence, and the whole comparison space is contractible.  Therefore
$\operatorname{CocLift}_p(f,x)$ is contractible whenever it is inhabited.\qedhere
\end{enumerate}
\end{proof}

\begin{definition}[Finite cocartesian exactness]\label{thz:def:exact}
The functor $p:E\to T$ is \textbf{finitely cocartesian exact} if, for every
finite simplicial set $J$, every diagram $d:J\to E$ whose edges are
$p$-cocartesian, and every limiting cone
$\lambda:J^{\triangleleft}\to T$ on $pd$:
\begin{enumerate}\renewcommand{\labelenumi}{(\roman{enumi})}
\item there exists a lift $\bar\lambda:J^{\triangleleft}\to E$ extending
      $d$, lying over $\lambda$, whose cone legs are $p$-cocartesian;
\item every such lift is a limiting cone in $E$.
\end{enumerate}
The empty diagram is included.
\end{definition}

Finite cocartesian exactness is an actual restriction on sources.
For example $[1]\to *$ admits cocartesian lifts of every base arrow,
using identities. Its empty diagram has a lifted cone with vertex
$0$; this vertex is not terminal in $[1]$, violating clause (ii).

\begin{lemma}[Contractibility of exact lifted cones]\label{thz:lem:exact-lift-cones}
Fix a finite diagram $d:J\to E$ whose edges are $p$-cocartesian and a
limiting cone $\lambda:J^{\triangleleft}\to T$ on $pd$.  Let
$\operatorname{CocLiftCone}_p(d,\lambda)$ be the space of lifts
$\bar\lambda:J^{\triangleleft}\to E$ extending $d$, lying over $\lambda$,
and having $p$-cocartesian cone legs.  If the two clauses of
Definition~\ref{thz:def:exact} hold for $(d,\lambda)$, then
$\operatorname{CocLiftCone}_p(d,\lambda)$ is contractible.
\end{lemma}
\begin{proof}
\leavevmode
\begin{enumerate}
\item \textbf{Choose a limiting lift.} Clause (i) supplies a point
$\bar\lambda\in\operatorname{CocLiftCone}_p(d,\lambda)$, and clause (ii)
makes it a limiting cone on $d$.  Let $\operatorname{LimCone}_E(d)$ and
$\operatorname{LimCone}_T(pd)$ denote the spaces of limiting cones on the two
diagrams.  Both are contractible: the first is inhabited by $\bar\lambda$ and
the second by $\lambda$, while Lemma~\ref{lem:universal-choices} identifies the
space of candidates for a fixed universal property as empty or contractible.
\item \textbf{Identify the space of limiting cones over the base.} Every $\mu\in\operatorname{LimCone}_E(d)$ is equivalent as a cone to
$\bar\lambda$.  Applying $p$ gives an equivalence
$p\mu\simeq p\bar\lambda=\lambda$, so $p\mu$ is limiting.  Hence $p$ restricts
to a map $p_*:\operatorname{LimCone}_E(d)\longrightarrow \operatorname{LimCone}_T(pd).$ It is a map between contractible spaces and therefore an equivalence.  Define
the fixed-base limiting-cone space by
\[
 \operatorname{LimCone}_E(d;\lambda):=
 \fib_\lambda\!\left(
   \operatorname{LimCone}_E(d)\longrightarrow
   \operatorname{LimCone}_T(pd)
 \right).
\]
It is contractible because it is a homotopy fiber of the equivalence $p_*$.
\item \textbf{Recover the cocartesian condition.} There is a natural inclusion $\operatorname{CocLiftCone}_p(d,\lambda) \longrightarrow \operatorname{LimCone}_E(d;\lambda)$ by clause (ii).  Conversely, a point
$\mu\in\operatorname{LimCone}_E(d;\lambda)$ consists of a limiting cone on
$d$ together with its specified identification with the fixed base cone
$\lambda$.  Since the homotopy fiber
$\operatorname{LimCone}_E(d;\lambda)$ is contractible and contains
$\bar\lambda$, there is a contractible space of equivalences
$\mu\simeq\bar\lambda$ inside that fixed-base homotopy fiber.  In particular
the equivalence is compatible with the chosen identification of both base
cones with $\lambda$.  Corresponding cone legs are therefore equivalent as
arrows of $E$ over the corresponding fixed legs of $\lambda$, and
cocartesianness is invariant under equivalence of arrows by
Lemma~\ref{thz:lem:cocart-calculus}.  Thus every leg of $\mu$ is
$p$-cocartesian.  The two spaces are consequently the same replete subspace
of the fixed-base cone space, and
\[
 \operatorname{CocLiftCone}_p(d,\lambda)
 \simeq \operatorname{LimCone}_E(d;\lambda)
 \simeq *.
\qedhere
\]
\end{enumerate}
\end{proof}

\begin{lemma}[Cocartesian cone fibers]\label{thz:lem:cocart-cone-fiber}
Let $K:T\to\Cat$, let $D:J\to T$ be a diagram indexed by a small
simplicial set, and let $\lambda:J^{\triangleleft}\longrightarrow T$ be a cone on $D$ with apex $\ell$.  Let
$p_K:\int_TK\to T$ be the cocartesian unstraightening.  If
$d:J\to\int_TK$ is a $p_K$-cocartesian diagram over $D$, write $\xi_d\in\lim_{j\in J}K(Dj)$ for the corresponding object under the cocartesian-section equivalence, and
let $\kappa_\lambda:K(\ell)\longrightarrow\lim_{j\in J}K(Dj)$ be the functor induced by cocartesian transport along the legs of $\lambda$.
Then there is a natural equivalence of spaces
\begin{equation}\label{thz:eq:cocart-cone-fiber}
 \operatorname{CocLiftCone}_{p_K}(d,\lambda)
 \simeq
 \fib_{\xi_d}\!\left(
   K(\ell)^{\simeq}
   \longrightarrow
   \left(\lim_{j\in J}K(Dj)\right)^{\simeq}
 \right),
\end{equation}
where the displayed fiber is taken on maximal subgroupoids.  No choice of a
cleavage enters either side.
\end{lemma}
\begin{proof}
\leavevmode
\begin{enumerate}
\item \textbf{Identify the cocartesian section category.} For a diagram of $\infty$-categories indexed by an arbitrary small
simplicial set, its limit is canonically the $\infty$-category of
cocartesian sections of its cocartesian unstraightening.
\cite[Cor.~3.3.3.2, dual form]{Lurie2009}
Applied to $KD:J\to\Cat$, this gives
\[
 \operatorname{Sect}_{\coc}
 \left(D^*\!\int_TK\longrightarrow J\right)
 \simeq
 \lim_{j\in J}K(Dj).
\]
Under this equivalence the given diagram $d$ corresponds to $\xi_d$.
\item \textbf{Compute the cone fiber.} An object $a\in K(\ell)$ determines the cocartesian transport diagram
$\kappa_\lambda(a)$ over $D$.  To extend the fixed diagram $d$ to a cone
over the specified base cone $\lambda$ with cocartesian cone legs is exactly
to give such an $a$ together with an equivalence of cocartesian sections $\kappa_\lambda(a)\simeq\xi_d.$ The $\infty$-groupoid of such pairs is the homotopy fiber of the induced
map on maximal subgroupoids at $\xi_d$, which proves
\eqref{thz:eq:cocart-cone-fiber}.  Naturality in $D$, $d$, and $\lambda$
is the naturality of pullback, unstraightening, and the section--limit
equivalence.\qedhere
\end{enumerate}
\end{proof}

\begin{lemma}[Exactness detection]\label{thz:lem:exact}
For $K:T\to\Cat$, let
$p_K:\int_TK\to T$ be its cocartesian unstraightening.  Then $p_K$ is
finitely cocartesian exact if and only if $K$ preserves finite limits.
\end{lemma}
\begin{proof}
\leavevmode
\begin{enumerate}
\item \textbf{Construct lifts from left exactness.} Let $D:J\to T$ be finite, let $\lambda$ be a limiting cone with apex
$\ell$, and write $\kappa_\lambda:K(\ell)\longrightarrow\lim_{j\in J}K(Dj).$
Suppose first that $K$ preserves this limit, so that $\kappa_\lambda$ is an
equivalence.  Let $d:J\to\int_TK$ be any cocartesian diagram over $D$, and
let $\xi_d$ be its corresponding object of $\lim_JKD$.  By
Lemma~\ref{thz:lem:cocart-cone-fiber},
\[
 \operatorname{CocLiftCone}_{p_K}(d,\lambda)
 \simeq
 \fib_{\xi_d}\!\left(
   K(\ell)^{\simeq}
   \longrightarrow
   \left(\lim_JKD\right)^{\simeq}
 \right)
 \simeq *.
\]
Thus a cocartesian lift cone exists, and its space of possible realizations
is contractible.
\item \textbf{Verify the limiting property of every lift.} Let $\bar\lambda$ be any such lift, with apex $a\in K(\ell)$.  For
$z\in K(r)$, its mapping-space comparison
\[
 \Map_{\int K}(z,a)\longrightarrow
 \lim_j\Map_{\int K}(z,dj)
\]
lies over the limiting-cone equivalence $\Map_T(r,\ell)\xrightarrow{\simeq} \lim_j\Map_T(r,Dj).$ The homotopy fiber over $h:r\to\ell$ is $\Map_{K(\ell)}(K(h)z,a)\longrightarrow \Map_{\lim_jK(Dj)} \bigl(\kappa_\lambda K(h)z,\xi_d\bigr).$ Because $\kappa_\lambda$ is an equivalence and
$\kappa_\lambda(a)\simeq\xi_d$, this fiber map is an equivalence.  A map
of spaces over an equivalent base is an equivalence when every homotopy
fiber map is one.  Hence the total mapping-space comparison is an
equivalence, so $\bar\lambda$ is limiting.  This proves both clauses of
finite cocartesian exactness.
\begin{samepage}
The mapping-space test lies in the commutative square
\[
\begin{tikzcd}[column sep=large,row sep=large]
 \Map_{\int K}(z,a)\ar[r]\ar[d]&
 \displaystyle\lim_j\Map_{\int K}(z,dj)\ar[d]\\
 \Map_T(r,\ell)\ar[r,"\simeq"']&
 \displaystyle\lim_j\Map_T(r,Dj).
\end{tikzcd}
\]
The fiber calculation above identifies the top comparison over each
point of the bottom row.
\end{samepage}
\item \textbf{Prove essential surjectivity of the comparison.} Conversely, assume that $p_K$ is finitely cocartesian exact.  For any
$\xi\in\lim_JKD$, the cocartesian-section equivalence supplies a
$p_K$-cocartesian diagram $d_\xi$ over $D$.  Clause (i) gives a
cocartesian lift cone over $\lambda$, and
Lemma~\ref{thz:lem:cocart-cone-fiber} therefore makes the fiber of
$\kappa_\lambda$ over $\xi$ inhabited.  Thus $\kappa_\lambda$ is
essentially surjective.
\item \textbf{Prove full faithfulness of the comparison.} For $a,b\in K(\ell)$, let $d_b$ be the cocartesian transport diagram
corresponding to $\kappa_\lambda(b)$.  Its canonical cocartesian cone has
apex $b$, and clause (ii) makes that cone limiting in $\int_TK$.  Test its
universal property against $a$ and take the homotopy fiber over the base cone
determined by $\id_\ell$.  This gives $\Map_{K(\ell)}(a,b)\simeq \Map_{\lim_jK(Dj)} (\kappa_\lambda a,\kappa_\lambda b).$ Hence $\kappa_\lambda$ is fully faithful.  It is therefore an equivalence.
Varying the finite diagram and limiting cone proves that $K$ preserves
finite limits.\qedhere
\end{enumerate}
\end{proof}

\subsection{Theorized objects and marked transport}

Exactness controls the finite-limit behavior of transport. We can now
specify which degrees must admit that transport. Keeping this requirement
as a marking allows the same definition to express fully represented
models and theories whose other degrees carry only operation kernels.

\begin{definition}[Marked theorization]\label{thz:def:thz}
An object of $\Thz_\tau(T)$ is a bounded small functor $p:E\to T$
which is finitely cocartesian exact and for which, for every
$f:s\to t$ in $\tau$ and every $x\in E_s$, the space
$\operatorname{CocLift}_p(f,x)$ is inhabited.  A morphism is a functor
over $T$ preserving $p$-cocartesian arrows whose degree lies in $\tau$.
All natural transformations over $T$ between such morphisms are retained.
A morphism is \textbf{strong} if it preserves every cocartesian arrow.  An
object is \textbf{represented} if it admits a cocartesian lift of every
arrow of $T$ from every object over its source.
\end{definition}

The exactness and marked-lift clauses are properties of $p$, not choices of
cleavages or universal cones.  Lemmas~\ref{thz:lem:cocart-calculus} and
\ref{thz:lem:exact-lift-cones} give contractibly unique universal
representatives whenever a semantic construction requires them.

More explicitly, each hom $\infty$-category is the full subcategory of
$\Fun_{/T}(E,E')$ on the functors satisfying the displayed preservation
property.  Identities and composites preserve those arrows, while vertical
transformations and their higher cells are inherited from the ambient slice.
Thus these homs form an $(\infty,2)$-category without an additional
comparison package.  Lemma~\ref{thz:lem:exact} gives an object
\[
 \Theta K=\left(\int_TK\longrightarrow T\right),
 \qquad K\in\mathcal L_T.
\]

Write $T^\tau$ for the marked $\infty$-category.  The marking also
constructs a wide $\infty$-subcategory $T_\tau\subseteq T$ as follows.  Let
$\operatorname{Ar}_\tau(T)\subseteq\Fun([1],T)^{\simeq}$ be the replete
subspace of marked arrows, and define a simplicial subspace of the complete
Segal nerve of $T$ by
\[
 N(T_\tau)_n:=
 \left\{\sigma\in\Fun([n],T)^{\simeq}\ \middle|\
 \sigma(i\to i+1)\in\operatorname{Ar}_\tau(T)
 \text{ for every }0\leq i<n\right\}.
\]
Because a Berman marking contains all equivalences, is replete, and is closed
under composition, every simplicial operator restricts to this subspace, its
Segal maps are the restrictions of those of $T$, and its completeness map is
the restriction of the completeness map of $T$.  Thus $N(T_\tau)$ is a
complete Segal space and defines the desired wide $\infty$-subcategory.  This
construction contains no choice of representatives.  If $T^\tau$ is bounded,
the same degreewise pullback construction is a bounded categorical record.  For
$p:E\to T$ put $E_\tau:=E\times_TT_\tau$.

\begin{lemma}[Marked straightening]\label{thz:lem:marked-straightening}
For $p\in\Thz_\tau(T)$, the restriction
\begin{equation}\label{thz:eq:marked-straightening}
 p_\tau:E_\tau\longrightarrow T_\tau
\end{equation}
is a cocartesian fibration.  An arrow of $E_\tau$ is $p_\tau$-cocartesian
if and only if it is $p$-cocartesian in $E$.  Consequently straightening
constructs a functor
\[
 C_\tau:T_\tau\longrightarrow\Cat,
 \qquad C_\tau(s)=E_s.
\]
For a marked arrow $a:s\to s'$, we write
$a_!:C_s\to C_{s'}$ for the value $C_\tau(a)$.  No choice of a cleavage
is part of this notation.
\end{lemma}
\begin{proof}
\leavevmode
\begin{enumerate}
\item \textbf{Construct lifts over marked arrows.} By Definition~\ref{thz:def:thz}, every arrow of $T_\tau$ admits a globally
$p$-cocartesian lift from every object over its source.  Such a lift remains
cocartesian after restriction to $E_\tau$, so $p_\tau$ has cocartesian
lifts of all base arrows.
\item \textbf{Compare restricted and global cocartesian arrows.} Conversely, let $v:x\to y$ be $p_\tau$-cocartesian over $f\in\tau$.
Choose only for this argument a globally $p$-cocartesian lift
$u:x\to y'$ over the same $f$.  The restricted universal property of
$u$ and that of $v$ produce mutually inverse comparison arrows between
$y$ and $y'$ over $\id_{t(f)}$.  Hence $v$ differs from $u$ by a
vertical equivalence.  Lemma~\ref{thz:lem:cocart-calculus} makes $v$
globally $p$-cocartesian.  Thus the two cocartesian classes agree.
\item \textbf{Apply straightening coherently.} The invariant cocartesian-fibration criterion and straightening now apply
to \eqref{thz:eq:marked-straightening}.  Contractibility of lift spaces
shows that any presentation by individual universal lifts gives the same
functor $C_\tau$ through a contractible space.\cite{HHLN2023}\qedhere
\end{enumerate}
\end{proof}

\subsection{Operation kernels and represented recovery}

For a fixed degree, source, and target, the operations are already present
as a mapping-space fiber. Their composition is balanced over the category
of intermediate objects. The recognition question is whether these
kernels come from transport functors; the answer requires both
corepresentable rows and the compatibility of their balanced composites.

For arbitrary $p:E\to T$, put $C_s=E_s$ and define the normalized
operation kernel
\begin{equation}\label{thz:eq:kernel}
 P_f(x,y)=\fib_f\bigl(
   \Map_E(x,y)\longrightarrow\Map_T(s,t)
 \bigr),
 \qquad f:s\to t.
\end{equation}
The kernel is the pullback
\[
\begin{tikzcd}[column sep=large]
 P_f(x,y)\ar[r]\ar[d]&\Map_E(x,y)\ar[d,"p"]\\
 \{f\}\ar[r]&\Map_T(s,t).
\end{tikzcd}
\]
Composition in $E$ gives balanced cells
\[
 \mu_{g,f}:P_g\odot P_f\longrightarrow P_{gf},\qquad
 (P_g\odot P_f)(x,z)=
 \int^{y\in C_t}P_f(x,y)\times P_g(y,z).
\]
The unit is normalized:
$P_{\id_s}\simeq\Map_{C_s}$.  The complete Segal diagram of
$E\to T$ supplies the whole coherent family of unit, composition, face,
degeneracy, and higher comparison maps.  Thus the balanced cells are
constructed from the given category over $T$; they are not appended after
binary composition has been selected.

\begin{lemma}[Functorial assembly of corepresenters]\label{thz:lem:corep-assembly}
Let $P:C^{\op}\times D\to\Sps$ be a profunctor.  If every row
$P(c,-)$ is corepresentable, then the representing objects assemble
canonically into a functor $F:C\longrightarrow D$ and a natural equivalence $P(c,d)\simeq\Map_D(Fc,d).$ The $\infty$-groupoid of such factorizations, including the displayed
natural equivalence, is contractible.
\end{lemma}
\begin{proof}
\leavevmode
\begin{enumerate}
\item \textbf{Assemble the corepresenters through Yoneda.} Curry $P$ to
$C^{\op}\to\Fun(D,\Sps)$.  By the rowwise hypothesis its image
lies in the replete essential image of covariant Yoneda
\[
 y_D^{\mathrm{cov}}:D^{\op}\longrightarrow\Fun(D,\Sps),
 \qquad d\longmapsto\Map_D(d,-).
\]
Restricted Yoneda is an equivalence onto this essential image.  Passing to
the corresponding functor categories therefore gives a contractible fiber
of lifts of the curried row functor.  Taking opposites constructs $F$ and
the stated natural equivalence.  All functorial coherence is inherited from this factorization through Yoneda.\qedhere
\end{enumerate}
\end{proof}

For arbitrary $p$, the normalized kernels and their coherent composition
are the profunctors of unital lax straightening. See
\cite[Cor.~6.2 and \S8, eq.~(2)]{Blom2024}, including the identification of
the binary comparison with composition in $E$.
Iteration gives the bars in Section~\ref{sec:higher-composition}.

\begin{theorem}[Base-independent kernel recognition]\label{thz:thm:kernel-recognition}
Let $B$ be an arbitrary $\infty$-category and let $p:E\to B$ be an arbitrary
functor.  For $s\in B$ put $C_s=E_s$, and for $f:s\to t$ define $P_f(x,y)=\fib_f\bigl( \Map_E(x,y)\longrightarrow\Map_B(s,t) \bigr).$ Composition in $E$ gives the balanced maps $\mu_{g,f}:P_g\odot P_f\longrightarrow P_{gf}.$ The following are equivalent.
\begin{enumerate}\renewcommand{\labelenumi}{(\roman{enumi})}
\item for every $f:s\to t$ and every $x\in C_s$, there exists a
      $p$-cocartesian arrow over $f$ with source $x$;
\item every $P_f(x,-)$ is corepresentable and every $\mu_{g,f}$ is an
      equivalence.
\end{enumerate}
When these conditions hold, the representing objects assemble through a
contractible space into transport functors $f_!:C_s\to C_t$.  The invariant
cocartesian-fibration criterion and straightening construct their identity,
composition, and all higher transport coherences, and these identify with the
composition diagram induced by $E$.  Equivalently, $p$ is the cocartesian
unstraightening of the resulting coherent functor
$B\to\Cat$.
\end{theorem}
\begin{proof}
\leavevmode
\begin{enumerate}
\item \textbf{Obtain corepresentable kernels from transport.} Assume (i).  The invariant cocartesian-fibration criterion identifies $p$ with
a cocartesian fibration.  Straightening gives transport functors
$f_!:C_s\to C_t$.  For a cocartesian lift $x\to f_!x$, its mapping-space
universal property gives, naturally in $y$, $P_f(x,y)\simeq\Map_{C_t}(f_!x,y).$ For composable $f:s\to t$ and $g:t\to r$, co-Yoneda followed by
cocartesian factorization gives
\[
 (P_g\odot P_f)(x,z)
 \simeq P_g(f_!x,z)
 \simeq P_{gf}(x,z).
\]
Both maps are induced by composition in $E$, so their composite is precisely
$\mu_{g,f}$.  Thus (ii) holds.
\item \textbf{Construct universal cocartesian arrows from corepresenters.} Conversely suppose (ii).  Lemma~\ref{thz:lem:corep-assembly}, applied to
$P_f:C_s^{\op}\times C_t\to\Sps$, constructs a functor
$f_!:C_s\to C_t$ and a natural equivalence $P_f(x,y)\simeq\Map_{C_t}(f_!x,y).$ Let $u_{f,x}:x\longrightarrow f_!x$ be the arrow over $f$ corresponding to $\id_{f_!x}$.  For $z\in C_r$
and every $g:t\to r$, co-Yoneda and invertibility of $\mu_{g,f}$ give
\begin{equation}\label{thz:eq:cocart-fibers}
 P_g(f_!x,z)\xrightarrow{\simeq}P_{gf}(x,z).
\end{equation}
The two sides are the fibers over $g$ of, respectively, $\Map_E(f_!x,z)\longrightarrow\Map_B(t,r)$ and
\[
 \Map_E(x,z)\times_{\Map_B(s,r)}\Map_B(t,r)
 \longrightarrow\Map_B(t,r).
\]
The maps \eqref{thz:eq:cocart-fibers} are exactly the fibers of the
comparison induced by $u_{f,x}$.  Since these are the homotopy fibers over every point
$g\in\Map_B(t,r)$, the fiberwise equivalence criterion for maps of spaces
over a fixed base makes the total comparison an equivalence.  Hence
$u_{f,x}$ is $p$-cocartesian, proving (i).
Equivalently, composition with $u_{f,x}$ makes the square
\[
\begin{tikzcd}[column sep=large]
 \Map_E(f_!x,z)\ar[r,"u_{f,x}^*"]\ar[d]&
 \Map_E(x,z)\ar[d]\\
 \Map_B(t,r)\ar[r,"f^*"']&\Map_B(s,r)
\end{tikzcd}
\]
a pullback. The kernel calculation verifies this universal property
fiber by fiber, uniformly in the next base arrow $g:t\to r$.
\item \textbf{Straighten the resulting fibration.} We have now constructed a $p$-cocartesian lift over every base arrow from
every source object.  The invariant cocartesian-fibration criterion therefore
applies to $p$, and straightening constructs a coherent functor $C:B\longrightarrow\Cat$ together with an equivalence over $B$ between $p:E\to B$ and its cocartesian
unstraightening.
\item \textbf{Identify the recovered transport and composition.} For every $f:s\to t$, the straightened transport functor $C(f)$ corepresents
the same rows as the functor $f_!$ already constructed above:
\[
 P_f(x,y)\simeq\Map_{C_t}(C(f)x,y)
 \simeq\Map_{C_t}(f_!x,y).
\]
Lemma~\ref{thz:lem:corep-assembly} therefore identifies
$C(f)\simeq f_!$ through a contractible space of representing
factorizations.  Under these identifications the represented image of
$\mu_{g,f}:P_g\odot P_f\to P_{gf}$ has direction
$C(gf)\to C(g)C(f)$: it acts by precomposition on
$\Map(C(g)C(f)x,-)$. At $\id_{C(g)C(f)x}$ its value classifies the
composite of the two chosen cocartesian lifts through the lift of
$gf$. The resulting vertical comparison is invertible since both
lifts are cocartesian. Its inverse is the straightening compositor
$C(g)C(f)\to C(gf)$, by the same factorization universal property.
This is the variance of corepresentables
\cite[Thm.~8.3.3.13]{Kerodon}. The straightening unit corresponds
to $P_{\id_s}\simeq\Map_{C_s}$.

Identity, associativity, and every higher
transport coherence are therefore the coherences of the single straightened
functor $C$, rather than separately supplied comparison data.  No finite-limit
structure on $B$ is used.\qedhere
\end{enumerate}
\end{proof}

\begin{theorem}[Representation and morphism recovery]\label{thz:thm:recovery}
For $p\in\Thz_\tau(T)$, representedness is equivalently the kernel
condition of Theorem~\ref{thz:thm:kernel-recognition}.  On this locus,
finite cocartesian exactness identifies $p$ with the cocartesian
unstraightening of a unique left exact functor $T\to\Cat$ through a
contractible space.  Consequently, on underlying $\infty$-categories, $\Thz_\tau(T)^{\mathrm{represented,strong}} \simeq\mathcal L_T,$ where the superscript denotes represented objects and strong morphisms.
Moreover, for $K,L\in\mathcal L_T$,
\begin{equation}\label{thz:eq:lax-recovery}
 \operatorname{Hom}_{\Thz_\tau(T)}(\Theta K,\Theta L)
 \simeq
 \Nat_{\mathrm{lax}}^{\tau\text{-strong}}(K,L),
\end{equation}
with lax comparison convention
$L(f)H_s\Rightarrow H_tK(f)$.
\end{theorem}
\begin{proof}
\leavevmode
\begin{enumerate}
\item \textbf{Recover the left-exact categorical model.} Theorem~\ref{thz:thm:kernel-recognition} constructs the coherent functor
$T\to\Cat$ on the represented locus.  Since $p$ is finitely
cocartesian exact, Lemma~\ref{thz:lem:exact} makes this functor left exact.
Conversely the cocartesian unstraightening of a left exact functor is
finitely cocartesian exact by the same lemma and is represented by
construction.  Contractibility of lift and representer spaces gives the
stated uniqueness.
\item \textbf{Recover lax morphisms and their strong locus.} For the hom statement, lax unstraightening identifies a functor
$F:\int K\to\int L$ over $T$ with fiber functors $H_s$ and coherent
comparisons $\alpha_f:L(f)H_s\longrightarrow H_tK(f).$ An arrow over $f$, represented by $v:K(f)x\to y$, is sent to $L(f)H_sx\xrightarrow{\alpha_f}H_tK(f)x \xrightarrow{H_tv}H_ty.$ The functor identities are exactly the coherent lax identities, and the same
correspondence on cylinders gives transformations and all higher cells.  The
functor preserves cocartesian arrows over $f$ precisely when $\alpha_f$ is
invertible.  Preservation over $\tau$ therefore gives
\eqref{thz:eq:lax-recovery}, while preservation over every base arrow gives
ordinary natural transformations.\cite{HHLN2023}\qedhere
\end{enumerate}
\end{proof}

\begin{corollary}[Recovery of categorified $\infty$-type theories]\label{thz:cor:categorification}
There is a natural equivalence $\Thz_{\tau_{\NU}}(\RNU)^{\mathrm{represented,strong}} \simeq \operatorname{Cat}(\TT).$ With all morphisms retained, the represented objects instead have the
$\tau_{\NU}$-strong lax transformation categories of
\eqref{thz:eq:lax-recovery}.
\end{corollary}
\begin{proof}
\leavevmode
\begin{enumerate}
\item \textbf{Apply categorical NU classification.} Combine Theorem~\ref{thz:thm:recovery} with Theorem~\ref{thm:nu-categorical}: $\Lex(\RNU,\Cat)\simeq\operatorname{Cat}(\TT).$ The assertion about hom categories is the second part of the recovery
theorem.\qedhere
\end{enumerate}
\end{proof}

\begin{example}[A nonrepresented object]\label{thz:ex:nonrepresented}
Let $T=\mathrm{FinCard}$ be the skeletal category of finite cardinals,
mark only its equivalences, and let $p:\one\to T$ select the cardinal
$1$.  Every finite diagram in $\one$ maps to a constant diagram at
$1$, whose limit in $\mathrm{FinCard}$ is $1$; its limiting
cone has the unique lift to $\one$.  Hence $p$ is finitely
cocartesian exact.  The only marked arrow with source $1$ is
$\id_1$, so all required marked lifts exist.

For any map $f:1\to2$, however, $C_1=\one$ and $C_2=\varnothing$,
so $P_f(*,-)$ is not corepresentable.  If $g:2\to1$ is the unique map,
then $gf=\id_1$ and
\[
 (P_g\odot P_f)(*,*)=\varnothing,
 \qquad
 P_{\id_1}(*,*)\simeq *,
\]
so $\mu_{g,f}:\varnothing\to *$ is not an equivalence.
\end{example}

\subsection{Structural limits and dependent operations}\label{sec:thz:dependent}

Recovery explains represented transport, but the construction must also
retain the ways inputs are assembled. Discrete cocartesian diagrams give
products; marked limiting diagrams give the structural comparisons used
in substitution. These provide the multi-input and dependent forms of the
same normalized operations.

\begin{proposition}[Structural exactness]\label{thz:prop:structural}
Every $p\in\Thz_\tau(T)$ has the following properties.
\begin{enumerate}\renewcommand{\labelenumi}{(\roman{enumi})}
\item $E$ has finite products, $p$ preserves them, and the product
      projections are $p$-cocartesian.
\item The marked restriction is the functor
      $C_\tau:T_\tau\to\Cat$ of
      Lemma~\ref{thz:lem:marked-straightening}.  If a finite limiting
      cone in $T$ has its diagram arrows and cone legs in $T_\tau$, then
      its induced comparison is an equivalence
 $C_\ell\xrightarrow{\simeq}\lim_jC_{Dj}.$ \item Strong morphisms preserve every cocartesian limiting cone occurring
      in Definition~\ref{thz:def:exact}.  An arbitrary morphism of
      $\Thz_\tau(T)$ preserves those such cones whose legs lie over
      $\tau$.
\item The finite-product structures on $E$ and $T$ have their canonical
      cartesian symmetric monoidal refinements, and $p$ has the canonical
      product-preserving symmetric monoidal refinement between them.
\end{enumerate}
\end{proposition}
\begin{proof}
\leavevmode
\begin{enumerate}
\item \textbf{Construct finite products.} For (i), apply finite cocartesian exactness to finite discrete diagrams,
including the empty one.  The lifted cones are product cones in $E$ over
the product cones in $T$, and their legs are cocartesian.  Hence $E$ has
finite products and $p$ preserves them.
\item \textbf{Recover limits along marked diagrams.} For (ii), Lemma~\ref{thz:lem:marked-straightening} supplies the marked
cocartesian fibration.  A coherent object of $\lim_jC_{Dj}$ is a globally
cocartesian diagram over $D$.  Exactness supplies a lifted cone and hence
essential surjectivity of $C_\ell\to\lim_jC_{Dj}$.  Testing a lifted
limiting cone on $b\in C_\ell$ against $a\in C_\ell$ and taking the fiber
over $\id_\ell$ gives full faithfulness exactly as in
Lemma~\ref{thz:lem:exact}.
\item \textbf{Check preservation of exact cones.} For (iii), the image cone has cocartesian legs whenever those legs are
preserved.  Each cone triangle writes a diagram edge $v$ as
$vu\simeq u'$, with $u$ and $u'$ cone legs.  Lemma~\ref{thz:lem:cocart-calculus}
then makes the image of $v$ cocartesian, even when its degree itself is not
marked.  Target exactness makes the image cone limiting.
\item \textbf{Identify the cartesian monoidal structures.} Finally, an $\infty$-category with finite products has an essentially
unique cartesian symmetric monoidal structure, and finite-product-preserving
functors are exactly the corresponding symmetric monoidal functors.  Apply
this to $E$, $T$, and $p$.\cite[Cor.~2.4.1.9]{LurieHA}\qedhere
\end{enumerate}
\end{proof}

For $x_i\in C_{s_i}$, Proposition~\ref{thz:prop:structural}(iv) supplies
the canonical cartesian symmetric-monoidal comparison $p\!\left(\prod_i x_i\right)\simeq\prod_i s_i.$ Using this constructed equivalence, for $f:\prod_i s_i\to t$ define the
finite-input operation space by
\[
 \operatorname{Mul}_f((x_i);y)=
 \fib_f\left(
  \Map_E\left(\prod_i x_i,y\right)
  \longrightarrow
  \Map_T\left(\prod_i s_i,t\right)
 \right).
\]
For
$u_i\in\operatorname{Mul}_{f_i}((x_{ij})_j;y_i)$ and
$v\in\operatorname{Mul}_g((y_i)_i;z)$, block substitution is induced by
\[
 \prod_{i,j}x_{ij}\simeq
 \prod_i\left(\prod_jx_{ij}\right)
 \xrightarrow{\prod_i u_i}
 \prod_i y_i\xrightarrow{v}z
\]
of degree $g\circ\prod_i f_i$.  Proposition~\ref{thz:prop:structural}(iv)
provides associativity, symmetry, permutation, diagonals, terminal maps,
and all higher product coherence; none of these is a new field of the
theorized object.  Strong morphisms preserve this entire structure.
An arbitrary morphism preserves the instances whose structural product
projections have marked degrees, by part (iii).

For a marked $a:s\to s'$, let
$a_!:C_s\to C_{s'}$ be the functor constructed in
Lemma~\ref{thz:lem:marked-straightening}.  For $A\in C_{s'}$ define $\Dep_p(a;A)=C_s\times_{C_{s'}}\{A\}.$ Thus a point of $\Dep_p(a;A)$ is a pair $(x,\epsilon_x)$ with
$x\in C_s$ and a specified equivalence
$\epsilon_x:a_!x\simeq A$.  If $\alpha:A\to B$ is an arrow of $C_{s'}$ and
$(x,\epsilon_x)\in\Dep_p(a;A)$,
$(y,\epsilon_y)\in\Dep_p(a;B)$, conjugation constructs a canonical
equivalence
\[
 c_{x,y}:\Map_{C_{s'}}(a_!x,a_!y)
 \xrightarrow{\simeq}\Map_{C_{s'}}(A,B),
 \qquad
 v\longmapsto \epsilon_y\,v\,\epsilon_x^{-1}.
\]
Define the normalized dependent kernel by
\[
 Q_\alpha(x,y):=
 \fib_\alpha\!\left(
  \Map_{C_s}(x,y)
  \xrightarrow{a_!}
  \Map_{C_{s'}}(a_!x,a_!y)
  \xrightarrow{c_{x,y}}
  \Map_{C_{s'}}(A,B)
 \right).
\]
This is the normalized operation kernel of the category over $C_{s'}$ given
by $a_!:C_s\to C_{s'}$, pulled back to the specified source and target
fibers.  Composition is therefore the normalized coherent composition of
\eqref{thz:eq:kernel} applied to this arbitrary-base functor, followed by the
canonical conjugation equivalences.  Consequently, if all rows of these
kernels are corepresentable and their composition comparisons are
equivalences, Theorem~\ref{thz:thm:kernel-recognition} constructs the
corresponding cocartesian transport functors. This construction applies
to an arbitrary base $C_{s'}$.

\begin{proposition}[Displayed operations and pasting]\label{thz:prop:displayed}
For a coherently commutative square $bf\simeq ga$ in $T$, with
$a:s\to s'$ and $b:t\to t'$ marked, there is a canonical natural map $\chi:P_f(x,y)\longrightarrow P_g(a_!x,b_!y).$ For $A\in C_{s'}$, $B\in C_{t'}$,
$\alpha\in P_g(A,B)$, and
\[
 (x,\epsilon_x)\in\Dep_p(a;A),\qquad
 (y,\epsilon_y)\in\Dep_p(b;B),
\]
conjugation gives a canonical equivalence $c^g_{x,y}:P_g(a_!x,b_!y)\xrightarrow{\simeq}P_g(A,B).$ The homotopy fibers
\begin{equation}\label{thz:eq:displayed-fiber}
 \widetilde P_\alpha(x,y):=
 \fib_\alpha\bigl(c^g_{x,y}\chi\bigr)
\end{equation}
form a profunctor $\widetilde P_\alpha: \Dep_p(a;A)^{\op}\times\Dep_p(b;B) \longrightarrow\Sps.$ These displayed operation families commute with parameter substitution and
coherent horizontal and vertical pasting.
\end{proposition}
\begin{proof}
\leavevmode
\begin{enumerate}
\item \textbf{Construct the comparison on operation kernels.} The marked cocartesian fibration of
Lemma~\ref{thz:lem:marked-straightening} supplies the transport functors
$a_!$ and $b_!$.  Cocartesian factorization supplies natural equivalences of
complete mapping-space families $\Lambda_{a,x}:= \lambda_{a,x}^*: P_g(a_!x,z)\xrightarrow{\simeq}P_{ga}(x,z)$ and postcomposition with the structural transport over $b$ supplies natural
maps $(\lambda_{b,y})_*: P_f(x,y)\longrightarrow P_{bf}(x,b_!y).$ The space of choices of a universal arrow denoted by $\lambda$ is
contractible by Lemma~\ref{thz:lem:cocart-calculus}; equivalently, these are
the canonical maps supplied by the straightened cocartesian fibration.  Using
the specified coherence $bf\simeq ga$, define
\begin{equation}\label{thz:eq:displayed-composite}
 \chi:=
 P_f(x,y)
 \xrightarrow{(\lambda_{b,y})_*}
 P_{bf}(x,b_!y)
 \xrightarrow{\simeq}
 P_{ga}(x,b_!y)
 \xrightarrow{\Lambda_{a,x}^{-1}}
 P_g(a_!x,b_!y).
\end{equation}
It is characterized by
\begin{equation}\label{thz:eq:displayed-defining}
 \chi(u)\lambda_{a,x}\simeq\lambda_{b,y}u.
\end{equation}
Since every arrow in \eqref{thz:eq:displayed-composite} is a natural map of
complete mapping-space families and $\Lambda_{a,x}^{-1}$ is the inverse of an
equivalence of such families, $\chi$ is natural in all variables and is
independent of any temporary lift representative.
\item \textbf{Pass to the displayed fibers.} Now let $(x,\epsilon_x)$ and $(y,\epsilon_y)$ lie in the displayed fibers.
Pre- and postcomposition with these equivalences gives
$c^g_{x,y}$, natural in the homotopy-pullback variables.  Hence
$c^g_{x,y}\chi$ is a natural map to the fixed operation space $P_g(A,B)$,
and its homotopy fiber over $\alpha$ is the profunctor
\eqref{thz:eq:displayed-fiber}.  Equivalently, if $\alpha_{x,y}:=\epsilon_y^{-1}\alpha\epsilon_x \in P_g(a_!x,b_!y),$ then $\widetilde P_\alpha(x,y)\simeq\fib_{\alpha_{x,y}}(\chi)$.
\item \textbf{Verify horizontal pasting.} For horizontal pasting, apply the equivalence $\Lambda_{a,x}$ to the two
candidate composites.  Ordinary associativity in $E$, naturality of the
specified square coherences, and the defining equation
\eqref{thz:eq:displayed-defining} identify their images.  Since
$\Lambda_{a,x}$ is an equivalence of complete mapping-space families, its
inverse transports this identification, together with every higher simplex,
back to the displayed operation family. The resulting comparison is
balanced in the intermediate fiber variable and therefore descends to
the coend defining the composite profunctor.
\item \textbf{Verify vertical pasting and substitution.} For vertical pasting, the composite of two structural cocartesian arrows is
cocartesian and lies in the same contractible lift space as the structural
transport for the composite marked arrow.  The two pasted formulas are thus
the inverse images, under the same cocartesian-factorization equivalence, of
the same ordinary composite in $E$.  Iterating this argument on the complete
simplicial composition diagram constructs every higher pasting compatibility.
Parameter substitution is pullback of the mapping-space, equivalence, and
homotopy-fiber diagrams already formed, and therefore commutes with the
construction.\qedhere
\end{enumerate}
\end{proof}


\subsection{Graded reconstruction of the operation category}\label{sec:graded-reconstruction}

The kernels would be an incomplete description if they forgot how arrows
of different degrees fit together. We therefore combine them into a
graded profunctor monad and reconstruct its complete string diagram. This
recovers the category over the doctrine, including its transformations,
and justifies treating the normalized operations as the theory itself.

Let $p:E\to T$ be bounded.  Let $C_p\subseteq E$ be the wide subcategory
of arrows with equivalence degree in $T$, and write $j:C_p\to E$ and
$q=pj$.  Set
\begin{equation}\label{thz:eq:aggregate}A_p(x,y):=\Map_E(jx,jy),\qquad
 H_q(x,y):=\Map_T(qx,qy),\qquad
 \gamma_p:A_p\longrightarrow H_q.
\end{equation}
Composition makes these profunctor monads, and $p$ makes $\gamma_p$ a monad
map.  These structures are derived from the two functors: co-Yoneda
identifies $A_p$ with the companion--conjoint monad
$j^\square\odot j_\square$, and similarly identifies $H_q$ with the monad
of $q_\square\dashv q^\square$.  Thus their higher coherences are already
those of the companion adjunctions.


The composition comparison is independent even of representability of all
normalized rows. Over the ordinal $[2]$, take the poset with relations
$x<y<z_1$ and $x<z_0<z_1$ and no others beyond transitivity, with
degrees $0,1,2,2$ for $x,y,z_0,z_1$. The representing transports out
of $x$ for $01,02$ select $y,z_0$, respectively, and that out of $y$
for $12$ selects $z_1$. The identity degrees give ordinary hom rows.
At $(x,z_0)$ the composite through $y$ is empty, whereas the direct
row is a point. This verifies the role of the composition test in
Theorem~\ref{thz:thm:kernel-recognition}.

\begin{definition}[Degree-exact grading]
\label{thz:def:graded-normal}
A $T$-graded profunctor monad is a category $C$, a functor $q:C\to T$
sending every arrow to an equivalence, an associative unital profunctor
monad $A$ on $C$, and a monad map $\gamma:A\to H_q$.  It is
\textbf{degree-exact} when the unit induces equivalences
\begin{equation}\label{thz:eq:degree-exact}
 \Map_C(x,y)\xrightarrow{\simeq}
 A(x,y)\times_{\Map_T(qx,qy)}\operatorname{Eq}_T(qx,qy).
\end{equation}
The monad is an associative algebra object in the profunctor endomorphism
category, so its coherent multiplication diagram is part of that notion.
Completeness of the construction below is a consequence of
\eqref{thz:eq:degree-exact}.
\end{definition}

Form its string simplicial space by
\begin{equation}\label{thz:eq:normal-string}
 \operatorname{Str}(A)_0=C^\core,\qquad
 \operatorname{Str}(A)_n=
 \int_{(x_0,\ldots,x_n)\in(C^\core)^{n+1}}
          \prod_{i=1}^n A(x_{i-1},x_i).
\end{equation}
The integral is the homotopy dependent sum over the displayed space.
Outer faces restrict the endpoint string; inner faces compose through the
balanced multiplication; degeneracies insert the unit.  For every monotone
map, use iterated multiplication on its nonempty consecutive blocks and
the unit on empty blocks.  The associative-algebra coherences make this
simplicial, and equivariance of framed cells under endpoint equivalences
makes all maps descend to the homotopy sums.

Degree exactness is a recognition condition on an independently given
monad. For $T=C=*$, any discrete group $G$ gives a profunctor monad
with its unique grading, whereas \eqref{thz:eq:degree-exact} reads
$*\to G$. For a monad constructed from $p$, the same condition follows
from the definition of $C_p$.

\begin{construction}[Intrinsic graded hom categories]
\label{con:graded-homs}
For degree-exact graded data $\mathcal G=(C,q,A,\gamma)$ and
$\mathcal G'=(C',q',A',\gamma')$, let $X_{\mathcal G}=\operatorname{Str}(A)$
and $X_{\mathcal G'}=\operatorname{Str}(A')$ be their complete string
simplicial spaces, with the degree maps to $N(T)$ obtained from $\gamma$
and $\gamma'$. Define the complete hom nerve by
\begin{equation}\label{eq:graded-hom-nerve}
 N_m\mathbf{GrHom}_T(\mathcal G,\mathcal G')
 :=\Map_{\operatorname{CSS}_{/N(T)}}
 \bigl(X_{\mathcal G}\times N([m]),X_{\mathcal G'}\bigr),
\end{equation}
where the product has degree given by its first factor. This is an
intrinsic formula in the graded string diagrams, not an extra class of
modification witnesses. At $m=0$ its vertices are graded morphisms:
a functor $F:C\to C'$ over $T$ with a coherent degree-preserving
monad map $\alpha:A\to F^*A'$. For a pair $(F,\alpha),(G,\beta)$,
a $1$-prism has components
$\eta_x\in A'_{\id_{qx}}(Fx,Gx)$ and coherent operation-naturality
comparisons $\eta_y\alpha(a)\simeq\beta(a)\eta_x$.
Higher prisms contain the compatibility for all operation strings.
\end{construction}

\begin{lemma}[Comparison of graded maps with Segal strings]
\label{lem:graded-map-comparison}
Let $\mathcal G=(C,q,A,\gamma)$ and
$\mathcal G'=(C',q',A',\gamma')$ be degree-exact graded
profunctor monads. The space of coherent degree-preserving monad
morphisms $(F,\alpha)$ is naturally equivalent to
\[
 \Map_{\operatorname{CSS}_{/N(T)}}
       (\operatorname{Str}(A),\operatorname{Str}(A')).
\]
In particular, a coherent monad map is completely determined by its
functor on equivalence-degree arrows, its map on framed operation
spaces, and its compatibility with the associative unital profunctor
multiplication. The equivalence is natural for simplicial families of
morphisms.
\end{lemma}
\begin{proof}
Write $X=X_{\mathcal G}$ and $X'=X_{\mathcal G'}$. For fixed $C$,
\cite[Cor.~5.19]{Kositsyn2021} identifies the category of profunctor
monads on $C$ with the category of maps $N(C)\to X$ of Segal spaces
inducing an equivalence on object spaces. Under this equivalence the
unit is the map from $N(C)$, and $X$ has the strings
\eqref{thz:eq:normal-string}: the arrow fibers are $A(x,y)$ and the
Segal equivalences give all higher levels. This fixed-base result
includes the balancing over all arrows of $C$.

For $F:C\to C'$, restriction of profunctors is lax monoidal. Its
binary constraint sends the coend over $u\in C$ of
$P'(Fx,Fu)\times Q'(Fu,Fy)$ to the coend over $u'\in C'$ by $u\mapsto Fu$.
The unit is $F$ on homs. Applying $F$ in each intermediate variable
of the iterated bar commutes with every face and degeneracy, giving
the iterated constraints. The fixed-base object for $F^*A'$ is therefore
the Segal space $X'_F$ with
\[
 (X'_F)_n=(C^\core)^{n+1}
        \times_{(C'^\core)^{n+1}}X'_n .
\]
Indeed its unit is induced by $F$ and its arrow fibers and composition
are the restricted ones just computed. A map $X\to X'_F$ under $N(C)$
is a map $X\to X'$ extending $F$ on $N(C)$, with the indicated
homotopies: this is the pullback universal property in each degree.
Taking the homotopy fiber expressing commutation with the maps to
$N(T)$ gives exactly the degree-preserving algebra maps over $F$.
Here $H_q$ corresponds to the same endpoint pullback of $N(T)$, by
the displayed formula applied to $q$.

By degree exactness and Segal, $N(C)\simeq X\times_{N(T)}N(T^\core)$
in every degree, and likewise for $C'$. Thus every map over $N(T)$
recovers $F$ by restriction, with contractible space of compatible
identifications. Moreover every equivalence edge of $X$ belongs to
this subspace, together with its inverse and inverse triangles.
Its equivalence edges are consequently exactly those of $N(C)$,
so $X$ is complete, and likewise $X'$. This proves the mapping-space
equivalence in the stated complete Segal category.

Finally apply the same comparison to $C\times[m]$ with kernel
$A(x,y)\times\Map_{[m]}(i,j)$ on $(x,i),(y,j)$ and degree
given by the first projection. It is degree-exact and its strings
are $X\times N([m])$. The construction commutes with every simplicial
operator and with composition of maps. It therefore retains the
noninvertible transformations as well as all their higher simplices.
\end{proof}

\begin{theorem}[Graded--Segal reconstruction]
\label{thz:thm:graded-segal-reconstruction}
For degree-exact data $\mathcal G=(C,q,A,\gamma)$, the string space is
complete Segal and canonically defines a bounded functor
$p_{\mathcal G}:E_{\mathcal G}\to T$.  Its mapping spaces are $A(x,y)$,
its degree map is $\gamma$, and the unit identifies $C$ with its wide
subcategory of arrows of equivalence degree.  These constructions extend
to mutually inverse equivalences of $(\infty,2)$-categories
\begin{equation}\label{thz:eq:graded-category-equivalence}
 \operatorname{Cat}^{\mathrm{bd}}_{/T}
 \simeq\operatorname{GrMnd}^{\mathrm{de}}_T.
\end{equation}
Consequently the graded operation data of \eqref{thz:eq:aggregate} recover
the original functor $p$, including all its comparisons over $T$.
\end{theorem}
\begin{proof}
\leavevmode
\begin{enumerate}
\item \textbf{Identify the complete Segal space.}
Lemma~\ref{lem:graded-map-comparison} identifies the strings with the
Segal space of the profunctor monad and proves completeness by
restriction to equivalence degrees.
\item \textbf{Reconstruct the projection to the base.} For the target define $N_q(T)_n=(C^\core)^{n+1} \times_{(T^\core)^{n+1}}N(T)_n.$ The Segal maps of $T$ give a simplicial equivalence
$s_q:N_q(T)\simeq\operatorname{Str}(H_q)$.  Its inverse is constructed
over a contractible inverse-candidate family.  Set
$N(E_{\mathcal G})=\operatorname{Str}(A)$ and define
$N(p_{\mathcal G})$ as the composite of $\operatorname{Str}(\gamma)$,
$s_q^{-1}$ and $N_q(T)\to N(T)$.  These constructions preserve boundedness.
The degree-one fiber is $A(x,y)$ and its map is $\gamma$;
\eqref{thz:eq:degree-exact} identifies the equivalence-degree arrows with
$\Map_C(x,y)$.  The object cores agree, so the unit is an equivalence
onto that wide subcategory.  Faces and degeneracies recover the monad
multiplication and unit, including their higher coherences.
\begin{samepage}
The reconstructed projection is displayed by
\[
\begin{tikzcd}[column sep=large,row sep=large]
 N(E_{\mathcal G})=\operatorname{Str}(A)
   \ar[r,"\operatorname{Str}(\gamma)"]\ar[d,"N(p_{\mathcal G})"']&
 \operatorname{Str}(H_q)\ar[d,"s_q^{-1}","\simeq"']\\
 N(T)&N_q(T)\ar[l].
\end{tikzcd}
\]
Here the lower map forgets the lifted object string in $C^\core$.
\end{samepage}
\item \textbf{Recover the original category from its kernels.} Conversely, for $p:E\to T$ one has $C_p^\core=E^\core$ and
$A_p(x,y)=\Map_E(x,y)$.  The Segal maps give
$N(E)\simeq\operatorname{Str}(A_p)$, compatibly with the grading to
$N(T)$.  The definition of $C_p$ proves degree exactness.  This establishes
inverse reconstruction on objects.
\item \textbf{Identify graded functors, modifications, and higher prisms.}
Lemma~\ref{lem:graded-map-comparison} identifies the object space
of graded morphisms with the mapping space of complete strings over
$N(T)$. To recover the entire hom $\infty$-category, use the defining
formula of Construction~\ref{con:graded-homs}: for every $m\geq0$,
\begin{align*}
 N_m\mathbf{GrHom}_T(\mathcal G,\mathcal G')
 &\simeq \Map_{\operatorname{CSS}_{/N(T)}}
       (X_{\mathcal G}\times N([m]),X_{\mathcal G'})\\
 &\simeq \Map_{(\Cat)_{/T}}
       (E_{\mathcal G}\times[m],E_{\mathcal G'})\\
 &\simeq N_m\Fun_{/T}
       (E_{\mathcal G},E_{\mathcal G'}).
\end{align*}
Here the degree map of $E_{\mathcal G}\times[m]$ is its projection
to $E_{\mathcal G}$ followed by $p_{\mathcal G}$. The second
comparison is full faithfulness and cartesian compatibility of the
complete Segal nerve; the third is the tensor--cotensor adjunction in
the slice. These are maps of the entire simplicial $m$-diagram.
Consequently
\[
 \mathbf{GrHom}_T(\mathcal G,\mathcal G')
 \simeq \Fun_{/T}(E_{\mathcal G},E_{\mathcal G'})
\]
as full $\infty$-categories. At $m=1$ one recovers all vertical
natural transformations, not merely invertible ones. Their higher
prisms are the corresponding coherent functors
$E_{\mathcal G}\times[m]\to E_{\mathcal G'}$ over $T$.
Composition of string maps corresponds to composition of these functors,
providing horizontal and vertical composition and interchange
coherences simultaneously. The object-level inverse constructions in
Steps~1--4 therefore extend to mutually inverse equivalences on every
hom $\infty$-category, proving
\eqref{thz:eq:graded-category-equivalence}.\qedhere
\end{enumerate}
\end{proof}

\subsection{Functoriality of theorization}
\label{subsec:theorization-functoriality}

The reconstruction theorem identifies the operations of a theory with
the arrows of its category over $T$. A map of theories should therefore
act on those operations by applying its underlying functor. Making this
action explicit also identifies what happens to dependent operation
comparisons and explains why lax maps are the appropriate maps between
represented theories.

\begin{theorem}[Functorial normalized operations]
\label{thz:thm:kernel-functoriality}
Let $F:(E\xrightarrow pT)\to(E'\xrightarrow{p'}T)$ be a morphism of
$\Thz_\tau(T)$ and put $C_s=E_s$, $C'_s=E'_s$.
Restriction gives $F_s:C_s\to C'_s$, and for every $f:s\to t$ there is
a natural map
\begin{equation}\label{eq:kernel-functoriality}
 \Phi_f:P_f^E(x,y)\longrightarrow P_f^{E'}(F_sx,F_ty).
\end{equation}
These maps preserve units, balanced composition, and every higher
composition comparison. They are functorial on vertical transformations
and all higher cells. They also preserve the marked-square comparison
maps of Proposition~\ref{thz:prop:displayed}, with their dependent fibers.
A finite cocartesian limiting cone whose legs have degree in $\tau$ is
carried to a limiting cone.
\end{theorem}
\begin{proof}
\leavevmode
\begin{enumerate}
\item \textbf{Construct the map on operation kernels.}
Apply $F$ to the mapping-space square over $\Map_T(s,t)$ and take its
homotopy fiber over $f$. This constructs \eqref{eq:kernel-functoriality}.
For $f:s\to t$ and $g:t\to u$, the maps on factors and the change of
coend variable along $F_t$ give
\[
 \int^{y\in C_t}P_f^E(x,y)\times P_g^E(y,z)
 \longrightarrow
 \int^{y'\in C'_t}
 P_f^{E'}(F_sx,y')\times P_g^{E'}(y',F_uz).
\]
Composing in $E'$ agrees with first composing in $E$ and then applying
$\Phi_{gf}$, since $F(vu)=F(v)F(u)$ coherently. Identity arrows give
the unit compatibility. More generally, the map of complete nerves
$N(E)\to N(E')$ commutes with every face and degeneracy, and hence
supplies the compatibility for all composable strings at once.


\item \textbf{Retain transformations.}
For a vertical transformation $\eta:F\Rightarrow G$, naturality gives
$\eta_yF(v)\simeq G(v)\eta_x$ in the fiber of the appropriate degree.
Performing the same construction on functors $E\times[n]\to E'$ over
$T$ records every simplex of the hom category. Thus the construction
retains noninvertible transformations and their higher coherence.


\item \textbf{Transport the dependent operations.}
The displayed-operation map is constructed by composition and universal
factorization through the marked cocartesian lifts. Their images under
$F$ are again such lifts. The contractible space of universal
factorizations identifies the two constructions, including their
pasting coherences. Taking homotopy fibers over the corresponding
parameter data gives the dependent comparison. Finally, the image of
the indicated limiting cone is a cocartesian cone over the same limiting
base cone. Finite cocartesian exactness of $p'$ makes it limiting.
\qedhere
\end{enumerate}
\end{proof}

\begin{proposition}[The action on represented theories]
\label{prop:represented-functoriality-explicit}
Under Theorem~\ref{thz:thm:recovery}, a map $\Theta K\to\Theta L$
corresponds to functors $H_s:K(s)\to L(s)$ and coherent comparisons
\[
 \alpha_f:L(f)H_s\longrightarrow H_tK(f),
 \qquad \alpha_f\text{ invertible for }f\in\tau.
\]
Its action on the represented kernel element $v:K(f)x\to y$ is
\begin{equation}\label{eq:lax-map-kernel-action}
 L(f)H_sx\xrightarrow{\alpha_{f,x}}H_tK(f)x
       \xrightarrow{H_t(v)}H_ty.
\end{equation}
For composable maps $H:K\to L$ and $J:L\to M$, the comparison is
\begin{equation}\label{eq:lax-map-composition}
 \alpha_f^{JH}=J_t(\alpha_f^H)\circ\alpha_f^JH_s.
\end{equation}
Strong maps are exactly those for which every $\alpha_f$ is invertible.
\end{proposition}
\begin{proof}
\leavevmode
\begin{enumerate}
\item \textbf{Recover the map and its lax comparison.}
Cocartesian factorization gives
\[
 P_f^K(x,y)\simeq\Map_{K(t)}(K(f)x,y).
\]
Substituting this equivalence in
\eqref{eq:lax-map-kernel-action} gives the required map on operations.
Conversely, apply the map on kernels to $\id_{K(f)x}$. The result is
$\alpha_{f,x}$. Compatibility with vertical composition gives naturality
in $x$, and compatibility with balanced composition gives the lax
composition law. 
\item \textbf{Retain modifications and marked strongness.}
Applying this argument to cylinders identifies all
modifications, so the two constructions are inverse on full hom
categories. An arrow $(f,v)$ in $\int K$ is cocartesian precisely when
$v$ is invertible. Testing the image of $(f,\id)$ proves the assertions
about the marking and strong maps. Applying the two kernel actions in
succession gives \eqref{eq:lax-map-composition}.
\qedhere
\end{enumerate}
\end{proof}

The general lax-map correspondence is already the principle of
\cite[Theorem~2.14]{Matsuoka2017}. Here its compatibility with finite
cocartesian exactness, the selected NU structural arrows, and dependent
operation fibers has been verified within the present framework. In
particular, functoriality adds no choice of transport or separate family
of coherence witnesses to a theorized object.

\section{\texorpdfstring{The $\infty$-type-theory microcosm}{The infinity-type-theory microcosm}}\label{sec:microcosm}

An operation theory gives the possible operations. An internal theory
chooses them coherently: an object at each arity, an operation over each
arrow, and the prescribed strong comparisons in structural directions.
Theorization expresses this intrinsically as a morphism from
$\id_T:T\to T$ to an operation theory $p:E\to T$. On represented
targets these morphisms are precisely marked sections, with their
ordinary morphisms and higher comparisons.

The next task is representability. A universal internal object should
record a source object and all ways to transport it to a test arity.
Comma categories do exactly this. Inverting the marked comparison
arrows imposes structural strongness, and left-exact reflection restores
the finite-limit requirements of a categorical model. We construct this
envelope for a general marked source before specializing it to the NU
operations. The seven operation squares and their dependent fibers then
supply the concrete content of the resulting microcosm.

\subsection{Intrinsic internalization from theorization}
\label{subsec:unrepresented-internalization}

The definition of an internal theory does not require represented
operations. A general theorized object has objects, operations, and
marked universal transport. A morphism from the identity object of the
doctrine selects these coherently, even when transport in other degrees
has no representative.

\begin{definition}[Intrinsic internal-object category]
\label{def:general-internalization}
For $p\in\Thz_\tau(T)$ define
\[
 \Int_\tau(p):=
 \Homc_{\Thz_\tau(T)}(\id_T,p).
\]
Its objects are sections $s:T\to E$ carrying $\tau$ to
$p$-cocartesian arrows, and its morphisms are all vertical natural
transformations between such sections.
\end{definition}

\begin{proposition}[Functorial internalization and represented recovery]
\label{prop:general-internalization}
Postcomposition constructs an $(\infty,2)$-functor
$\Int_\tau:\Thz_\tau(T)\to\Cat$.
For $K\in\mathcal L_T$ there are natural equivalences
\begin{equation}\label{eq:general-internalization-recovery}
 \Int_\tau(\Theta K)
 \simeq\laxlim_T^\tau K
 \simeq\Nat(\Omega_{T,\tau},K).
\end{equation}
For $T=\RNU$ and $\tau=\tau_{\NU}$, this extends the internal
NU theory construction to unrepresented theorized macrocosms.
\end{proposition}
\begin{proof}
\leavevmode
\begin{enumerate}
\item \textbf{Construct intrinsic internalization.}
The identity $T\to T$ is a finitely cocartesian-exact object with all
arrows cocartesian. A map of theorized objects preserves the marked
cocartesian arrows of a section, so postcomposition is defined. It acts
on vertical transformations and their higher cells by ordinary functor
composition.


\item \textbf{Recover the represented section description.}
In kernels a section consists of $x_s\in E_s$ and
$\xi_f\in P_f(x_s,x_t)$, with units, composition, and the higher
simplices of a functor $T\to E$; $\xi_f$ is cocartesian for
$f\in\tau$. A morphism consists of the vertical arrows compatible
with these operation strings. For $p=\Theta K$, the represented kernel
formula turns $\xi_f$ into $K(f)x_s\to x_t$, and cocartesianness
becomes invertibility. This is the marked lax-cone description on full
categories. The envelope of $(\id_T,\tau)$ gives the second
equivalence in \eqref{eq:general-internalization-recovery}.
\qedhere
\end{enumerate}
\end{proof}

For represented inputs, internalization is the limit-preserving
functor $K\mapsto\Nat(\Omega_{T,\tau},K)$ on $\mathcal L_T$.
The intrinsic definition retains the same marked-section calculation
for unrepresented operation theories. In
Section~\ref{sec:higher-composition} we construct its limit-closed
marked-cocartesian ambient category; Section~\ref{sec:stratified-microcosm}
then applies internalization in successive categorical degrees.

\subsection{Marked sections}

An internal object must choose operations compatibly over all arities.
A section makes exactly those choices. Factoring its arrows into
transport followed by a vertical comparison identifies the usual lax
cone, and makes structural strongness the invertibility of that vertical
factor.

Let $C:D\to\Cat$ and let $p_C:\int_D C\to D$ be its cocartesian
unstraightening. A section assigns $x_d\in C(d)$ and arrows over each
$u:d\to d'$. Cocartesian factorization identifies the latter with arrows
$\ell_u:C(u)x_d\to x_{d'}$. Composition in the total category supplies
their coherent lax composition law.
The factorization over $u$ and the base-change construction are
\[
\begin{tikzcd}[column sep=large]
 (d,x_d)\ar[r,"\mathrm{cocart}"]\ar[dr,"s(u)"']&
 (d',C(u)x_d)\ar[d,"\ell_u"]\\
 &(d',x_{d'})
\end{tikzcd}
\qquad
\begin{tikzcd}[column sep=large]
 D\times_T\int_TK\ar[r]\ar[d]&\int_TK\ar[d]\\
 D\ar[r,"a"']&T.
\end{tikzcd}
\]
The vertical factor in the first triangle lies in the fiber over $d'$;
it is invertible exactly when $s(u)$ is cocartesian.


\begin{theorem}[Marked sections and base change]
\label{prop:marked-sections}\label{thm:marked-base-change}
There is a natural equivalence of $\infty$-categories
\begin{equation}\label{eq:marked-sections}
 \laxlim\nolimits_D^W C\simeq
 \Fun^{W\text{-}\coc}_{/D}\left(D,\int_D C\right).
\end{equation}
The right side is the full subcategory of sections carrying $W$ to
cocartesian arrows. For $a:D^W\to T$ and $K:T\to\Cat$, it gives
\[
 \operatorname{Sect}^{\coc(W)}
 (D\times_T\textstyle\int_TK\to D)
 \simeq\laxlim_D^W(Ka)
 \simeq\operatorname{Nat}^{W\text{-strong lax}}(\one,Ka).
\]
These equivalences retain every comparison between the selected sections.
For any lax cone, the arrows on which its comparison is invertible form
a marking.
\end{theorem}
\begin{proof}
\leavevmode
\begin{enumerate}
\item \textbf{Identify marked cones with sections.} The correspondence just described is the section description of lax limits
\cite{Berman2020,HHLN2023}. Explicitly, an arrow of $\int C$ over $u$
factors through its cocartesian lift with vertical part $\ell_u$. Such an
arrow is cocartesian exactly when $\ell_u$ is invertible. The contractible
space of factorizations and the composition diagram in $\int C$ give
mutually inverse constructions on coherent cones. Applied to families
parametrized by $[n]$, they also identify modifications and all higher
cells. Restricting the objects by the displayed equivalence predicate gives
\eqref{eq:marked-sections} as a full-subcategory equivalence.
\item \textbf{Pass from generators to the full marking.} Identity comparisons are invertible. For composable arrows the comparison
for their composite is the composite of the transported comparisons.
Thus invertibility is closed under composition and equivalence of arrows,
and includes every equivalence of $D$. This proves the final assertion. Naturality of unstraightening gives
$D\times_T\int_TK\simeq\int_DKa$, proving the base-change formula.\qedhere
\end{enumerate}
\end{proof}

The relative section notation is
$\Fun^{W\text{-}\coc}_{/T}(D,\int_TK)$; the section condition is taken
over $D$ through the pullback of the unstraightening.

\subsection{The universal categorified envelope}\label{sec:thz:envelope}

The section description still quantifies over a whole family of choices.
To represent them by one source object, evaluate a proposed section at a
source arity and then transport it to a test arity. The comma category
records both steps. Its marked localization and lex reflection are
therefore forced by the desired section universal property.

We construct the envelope for a slightly larger input.  Let
$q:(D,W)\to T$ be a marked small $\infty$-category over $T$; the marked
arrows $W$ need not be cocartesian for $q$.  For $K:T\to\Cat$ write
\[
 \Fun^{W\text{-}\coc}_{/T}
   (D,\textstyle\int K)
\]
for the full $\infty$-category of functors over $T$ carrying every arrow
of $W$ to a $p_K$-cocartesian arrow.

\begin{lemma}[The marked comma construction]\label{thz:lem:comma}
For $t\in T$ let $C_q(t)=q\downarrow t$.  Its objects are pairs
$(d,h:qd\to t)$.  Let $W_t$ be the Berman marking generated by the
arrows
\begin{equation}\label{thz:eq:Wt}
 (d,h'q(u))\longrightarrow(d',h'),
 \qquad u:d\to d'\text{ in }W,
\end{equation}
together with all equivalences.  The categorical localization $R_{q,W}(t):=C_q(t)[W_t^{-1}]$ exists and has the functor-category universal property; it is formed
at a sufficiently large subsequent level of the fixed universe hierarchy.\cite[\S6.3]{Kerodon}
These localizations assemble functorially into
$R_{q,W}:T\to\Cat$, and for every
$K:T\to\Cat$ there is a natural equivalence of full hom $\infty$-categories
\begin{equation}\label{thz:eq:comma-UP}
 \Nat(R_{q,W},K)
 \simeq
 \Fun^{W\text{-}\coc}_{/T}
   (D,\textstyle\int K).
\end{equation}
\end{lemma}
\begin{proof}
\leavevmode
\begin{enumerate}
\item \textbf{Assemble the marked comma diagram.} Before localizing, postcomposition in the comma variable assembles into a
functor of marked $\infty$-categories
\[
 C_q^\dagger:T\longrightarrow\operatorname{Cat}_\infty^+,
 \qquad
 C_q^\dagger(t)=(q\downarrow t,W_t).
\]
Indeed, for $h:t\to t'$, postcomposition gives
$h_*:q\downarrow t\to q\downarrow t'$, and it sends every generator
\eqref{thz:eq:Wt} of $W_t$ to the corresponding generator of $W_{t'}$.
The comma construction is functorial on all simplices of $T$, so these maps
already form a coherent marked categorical diagram; no objectwise choices or
separately supplied higher comparisons occur.
\item \textbf{Localize functorially.} The functorial localization theorem in
\cite[\S6.3]{Kerodon} constructs a functor $\operatorname{Loc}:\operatorname{Cat}_\infty^+ \longrightarrow\Cat$ together with the localization maps, functorially in the marked source.
Define $R_{q,W}:=\operatorname{Loc}\circ C_q^\dagger.$ Then
$R_{q,W}(t)\simeq(q\downarrow t)[W_t^{-1}]$ with its stated
functor-category universal property, and all identity, composition, and
higher functorial coherences are those of this composite functor.
\item \textbf{Construct the correspondence on objects and arrows.} First omit the localization.  A natural transformation $G:C_q\to K$
defines $F:D\to\int K$ over $T$ by $F(d)=G_{qd}(d,\id_{qd}).$ For $u:d\to d'$, the comma arrow
$(d,q(u))\to(d',\id)$ in $C_q(qd')$ gives the vertical comparison $K(q(u))F(d)\longrightarrow F(d').$ Conversely, from $F:D\to\int K$ over $T$, define $G_t(d,h)=K(h)F(d).$ For a comma arrow represented by $u:d\to d'$ with
$h'q(u)\simeq h$, apply $K(h')$ to the vertical comparison of $F(u)$.
Every $(d,h)$ is obtained from $(d,\id)$ by transport along $h$, and
every comma arrow is obtained from the displayed vertical comparison
followed by that transport.  Hence these two constructions are inverse.
The comma relation and its image are
\[
\begin{tikzcd}[column sep=large]
 qd\ar[r,"q(u)"]\ar[dr,"h"']&qd'\ar[d,"h'"]\\
 &t
\end{tikzcd}
\qquad
\begin{tikzcd}[column sep=large]
 K(h)F(d)\ar[r,"\simeq"]&
 K(h')K(q(u))F(d)\ar[d,"K(h')\ell_u"]\\
 &K(h')F(d').
\end{tikzcd}
\]
\item \textbf{Recover transformations and higher cells.} Applying the same formulas to cylinders $D\times[n]$ identifies all
transformations and higher cells.
\item \textbf{Identify the marking condition.} Under this equivalence the generator \eqref{thz:eq:Wt} is sent to
$K(h')$ applied to the vertical comparison of $F(u)$.  If $F(u)$ is
cocartesian, that comparison is invertible.  Conversely, taking
$t=qd'$ and $h'=\id$ shows that inversion of every such generator
forces the vertical comparison of $F(u)$ to be invertible, hence $F(u)$
is cocartesian.  The functor-category universal property of localization
therefore gives \eqref{thz:eq:comma-UP}.\qedhere
\end{enumerate}
\end{proof}

\begin{lemma}[Explicit left-exact reflector]\label{thz:lem:lex-reflector}
There is a left adjoint $L_{\mathrm{lex}}:\Fun(T,\Cat) \longrightarrow\mathcal L_T$ to the inclusion.  Its universal property is an equivalence of full hom
$\infty$-categories
\begin{equation}\label{thz:eq:enriched-reflector}
 \Nat_{\mathcal L_T}(L_{\mathrm{lex}}X,K)
 \simeq\Nat(X,K).
\end{equation}
\end{lemma}
\begin{proof}
\leavevmode
\begin{enumerate}
\item \textbf{Specify the localization maps.} The category $\Fun(T,\Cat)$ is presentable because $T$ is small and
$\Cat$ is presentable.  For $t\in T$ put
$h_t=\Map_T(t,-)$, regarded as a groupoid-valued object of
$\Fun(T,\Cat)$.  For every finite limiting cone
$\lambda$ with apex $\ell$ and diagram $D:J\to T$, and every $n\geq0$,
form
\begin{equation}\label{thz:eq:lex-maps}
 \operatorname*{colim}_{j\in J^{\op}}
   (h_{Dj}\times[n])
 \longrightarrow h_\ell\times[n].
\end{equation}
Here $\times[n]$ is the $\Cat$-copower by the finite category $[n]$.
Because $T$ is small and the indexing shapes are finite, these maps form a
small collection.  We use all limiting cones rather than choose one cone for
each diagram, so no finite-limit cleavage on $T$ enters the construction.
Localizing the presentable functor category at this collection exists by the
presentable-localization theorem.\cite[\S5.5.4]{Lurie2009}
\item \textbf{Identify their local objects.} For every $t\in T$, copowered Yoneda gives a natural equivalence $\Map_{\Fun(T,\Cat)}(h_t\times[n],K) \simeq \Map_{\Cat}([n],K(t)).$ Therefore mapping \eqref{thz:eq:lex-maps} into $K$ gives precisely $\Map_{\Cat}([n],K(\ell)) \longrightarrow \lim_j\Map_{\Cat}([n],K(Dj)).$ Thus $K$ is local for every map \eqref{thz:eq:lex-maps} exactly when this
map is an equivalence for every finite limiting cone and every $n\geq0$.
These are the levels of the complete Segal nerve of the comparison $K(\ell)\longrightarrow\lim_jK(Dj).$ The complete Segal nerve is fully faithful and limits in $\Cat$ are
computed levelwise in that presentation, so the levelwise equivalences are
equivalent to the displayed comparison being an equivalence in
$\Cat$.  Hence the local objects are exactly the finite-limit-preserving
functors.  This constructs $L_{\mathrm{lex}}$ on objects and arrows.
\begin{samepage}
Copowered Yoneda identifies the localization test with the nerve-level
limit comparison in the square
\[
\begin{tikzcd}[column sep=2em,row sep=large]
 \Map(h_\ell\times[n],K)\ar[r]\ar[d,"\simeq"']&
 \displaystyle\lim_j\Map(h_{Dj}\times[n],K)\ar[d,"\simeq"]\\
 \Map_{\Cat}([n],K(\ell))\ar[r]&
 \displaystyle\lim_j\Map_{\Cat}([n],K(Dj)).
\end{tikzcd}
\]
The upper mapping spaces are taken in $\Fun(T,\Cat)$.
\end{samepage}
\item \textbf{Upgrade the adjunction to full hom categories.} For the enriched statement, apply the mapping-space localization
adjunction to every pointwise cotensor $K^{[n]}$. These targets remain
left exact, and the equivalences are natural in $[n]$.
Lemma~\ref{lem:complete-nerve-comparison} gives
\eqref{thz:eq:enriched-reflector}. Its unit defines the reflected arrow
$L_{\mathrm{lex}}f$ as the transpose of $\eta_Yf$; naturality supplies
composition and all higher comparisons.\qedhere
\end{enumerate}
\end{proof}

\begin{theorem}[Universal categorified envelope]\label{thz:thm:envelope}
Define $\Env_T(q,W):=L_{\mathrm{lex}}R_{q,W}.$ Then for every $K\in\mathcal L_T$ there is a natural equivalence
\begin{equation}\label{thz:eq:envelope-UP}
 \Nat_{\mathcal L_T}(\Env_T(q,W),K)
 \simeq
 \Fun^{W\text{-}\coc}_{/T}
   (D,\textstyle\int K)
\end{equation}
of full hom $\infty$-categories.

For $p\in\Thz_\tau(T)$ let $W_p$ be the class of $p$-cocartesian
arrows whose degrees lie in $\tau$.  It is a Berman marking by
Lemma~\ref{thz:lem:cocart-calculus} and closure of $\tau$ under equivalence
and composition.  Put $\mathsf F_\tau(p)=\Env_T(p,W_p).$ Then there is an enriched adjunction
\begin{equation}\label{thz:eq:free-adjunction}
 \mathsf F_\tau:\Thz_\tau(T)
 \rightleftarrows\mathcal L_T:\Theta,
 \qquad
 \Nat_{\mathcal L_T}(\mathsf F_\tau p,K)
 \simeq
 \operatorname{Hom}_{\Thz_\tau(T)}(p,\Theta K).
\end{equation}
\end{theorem}
\begin{proof}
\leavevmode
\begin{enumerate}
\item \textbf{Compose the universal properties.} Compose Lemma~\ref{thz:lem:comma} with
Lemma~\ref{thz:lem:lex-reflector}.  This gives
\eqref{thz:eq:envelope-UP} directly.
\item \textbf{Construct the universal marked functor.} The identity of $\Env_T(q,W)$ under that equivalence constructs the
universal marked functor $D\longrightarrow\int_T\Env_T(q,W).$ On objects it begins with
$d\mapsto(d,\id_{qd})$ in the comma category, then passes through the
localization and the unit of $L_{\mathrm{lex}}$.  An arrow of $W$ has
invertible vertical comparison after the comma localization, hence maps
to a cocartesian arrow.  Thus the universal map itself is constructed by
the same universal property and is not an extra field.
\item \textbf{Verify source and doctrine naturality.} A marked functor $(D,W)\to(D',W')$ over $T$ induces maps of comma
categories preserving the generated inverted arrows.  Localization and
$L_{\mathrm{lex}}$ therefore give functoriality of $\Env_T$, including
all higher transformations.  For $W=W_p$, the right-hand side of
\eqref{thz:eq:envelope-UP} is exactly the hom $\infty$-category in
Definition~\ref{thz:def:thz}.  The target $\Theta K$ is a theorized object
by Lemma~\ref{thz:lem:exact}.  This proves
\eqref{thz:eq:free-adjunction}.  The unit and counit are the transposes of
the relevant identity morphisms, and the triangle identities are the two
inverse identities of the natural hom-category equivalence.\qedhere
\end{enumerate}
\end{proof}

Let $\operatorname{MarkCat}^{\mathrm{bd}}_{/T}$ denote the $\infty$-category
of bounded small marked $\infty$-categories $q:(D,W)\to T$ and marked
functors over $T$.  For $K\in\mathcal L_T$, let
\[
 U_T^{\coc}(K):=
 \left(\int_TK\longrightarrow T,\operatorname{Coc}(p_K)\right),
\]
where the marking consists of all $p_K$-cocartesian arrows.
Lemma~\ref{thz:lem:cocart-calculus} makes this a Berman marking.

\begin{corollary}[Envelope adjunction]\label{thz:cor:envelope-adjunction}
The envelope construction and fully cocartesian-marked unstraightening form
an adjunction on underlying $\infty$-categories $\Env_T: \operatorname{MarkCat}^{\mathrm{bd}}_{/T} \rightleftarrows \mathcal L_T:U_T^{\coc}.$ More precisely,
\[
 \Map_{\mathcal L_T}(\Env_T(D,W),K)
 \simeq
 \Map_{\operatorname{MarkCat}^{\mathrm{bd}}_{/T}}
 \bigl((D,W),U_T^{\coc}(K)\bigr).
\]
The enriched equivalence of hom $\infty$-categories is
\eqref{thz:eq:envelope-UP}.
\end{corollary}
\begin{proof}
\leavevmode
\begin{enumerate}
\item \textbf{Take cores of the enriched universal property.} The right-hand mapping space is the core of
$\Fun^{W\text{-}\coc}_{/T}(D,\int K)$.  Taking cores in
\eqref{thz:eq:envelope-UP} gives the displayed adjunction.  Naturality in both
variables follows from the functoriality constructed in
Theorem~\ref{thz:thm:envelope}; hence the unit, counit, and triangle identities
are those determined by the same natural hom equivalence.\qedhere
\end{enumerate}
\end{proof}

\begin{remark}[Adjunction versus reflection]
The reflector $L_{\mathrm{lex}}$ is a reflection of diagram categories.
The adjunction $\mathsf F_\tau\dashv\Theta$ has a different hom universal
property.  By \eqref{thz:eq:lax-recovery}, morphisms between represented
theorized objects are $\tau$-strong lax transformations; ordinary natural
transformations are the strong ones.  The inclusion of ordinary
transformations need not be an equivalence.  Thus $\Theta$ need not be
fully faithful and the categorified-envelope adjunction need not be a
reflection.
\end{remark}

\begin{proposition}[Recovery by the fully marked envelope]\label{thz:prop:fully-marked}
Let $K\in\mathcal L_T$ and mark $\int_TK$ by all of its cocartesian
arrows, writing the resulting marking as $W_K^{\mathrm{all}}$.  Then
\[
 R_{p_K,W_K^{\mathrm{all}}}\simeq K,
 \qquad
 \Env_T(p_K,W_K^{\mathrm{all}})\simeq K.
\]
In particular, if $\tau$ marks every arrow of $T$, the counit
$\mathsf F_\tau\Theta K\to K$ is an equivalence.
\end{proposition}
\begin{proof}
\leavevmode
\begin{enumerate}
\item \textbf{Identify the representing diagram by straightening.} For every $L:T\to\Cat$, Lemma~\ref{thz:lem:comma} gives
\[
 \Nat(R_{p_K,W_K^{\mathrm{all}}},L)
 \simeq
 \Fun^{\mathrm{all\ cocart}}_{/T}
   (\textstyle\int K,\int L).
\]
Ordinary straightening identifies the right side with $\Nat(K,L)$ on
the full hom $\infty$-categories.  Naturality and Yoneda give the first
equivalence.  Since $K$ is already left exact, applying
$L_{\mathrm{lex}}$ gives the second.  For the maximal base marking,
$W_{\Theta K}=W_K^{\mathrm{all}}$, and the resulting comparison is the
counit of \eqref{thz:eq:free-adjunction}.\qedhere
\end{enumerate}
\end{proof}


\subsection{Changing the source, marking, and doctrine}
\label{subsec:envelope-functoriality}

An envelope should depend on the marked problem it represents, rather
than on a particular presentation of that problem. The comma construction
makes this dependence explicit. It also distinguishes enlarging a marking,
which imposes further invertibilities, from changing the doctrine in which
the same marked source is interpreted.

\begin{theorem}[Functoriality of the marked envelope]
\label{thz:thm:envelope-functoriality}
Let $b:T\to S$ preserve finite limits. There is an enriched adjunction
\[
 b_!^{\mathrm{lex}}:\mathcal L_T\rightleftarrows\mathcal L_S:b^*,
 \qquad b_!^{\mathrm{lex}}=L_{\mathrm{lex},S}\operatorname{Lan}_b.
\]
For every bounded marked source $q:(D,W)\to T$, there is a canonical
equivalence
\begin{equation}\label{eq:envelope-doctrine-functoriality}
 b_!^{\mathrm{lex}}\Env_T(q,W)\simeq\Env_S(bq,W).
\end{equation}
A marked functor $F:(D,W)\to(D',W')$ with $q'F\simeq bq$ induces
\begin{equation}\label{eq:envelope-square-functoriality}
 b_!^{\mathrm{lex}}\Env_T(q,W)\longrightarrow\Env_S(q',W').
\end{equation}
These constructions respect identities, composition, and higher
transformations. On represented objects, reindexing is given by
$b^*(\Theta_SK)\simeq\Theta_T(Kb)$.
\end{theorem}
\begin{proof}
\leavevmode
\begin{enumerate}
\item \textbf{Construct the change-of-doctrine adjunction.}
Restriction along $b$ preserves left exactness. Compose left Kan
extension with the enriched lex reflector to obtain the adjunction;
its full-hom form also follows by applying the mapping-space adjunction
to the pointwise cotensors $K^{[n]}$.

\item \textbf{Identify the envelope after change of doctrine.}
For $K\in\mathcal L_S$, the envelope universal property gives
\begin{align*}
 \Nat(b_!^{\mathrm{lex}}\Env_T(q,W),K)
 &\simeq\Nat(\Env_T(q,W),Kb)\\
 &\simeq\Fun^{W\text{-}\coc}_{/T}(D,\textstyle\int_TKb)\\
 &\simeq\Fun^{W\text{-}\coc}_{/S}(D,\textstyle\int_SK)\\
 &\simeq\Nat(\Env_S(bq,W),K).
\end{align*}
The middle equivalence uses
$\int_TKb\simeq T\times_S\int_SK$ and the pullback description of
cocartesian transport. Enriched Yoneda proves
\eqref{eq:envelope-doctrine-functoriality}.


\item \textbf{Make the comparison functorial.}
The map associated to $F$ begins on comma categories with
$(d,h:q(d)\to t)\mapsto(Fd,bh:q'(Fd)\to b(t))$.
It preserves every generator of the marked localization. Compose with
localization and the lex-reflection unit, and take the adjoint transpose;
this constructs \eqref{eq:envelope-square-functoriality}. The comma maps
are functorial on entire diagrams, so their induced maps supply all
composition comparisons. The represented reindexing formula is the
unstraightening equivalence just used, and $Kb$ is left exact.
\qedhere
\end{enumerate}
\end{proof}

\begin{corollary}[Enlarging the structural marking]
\label{cor:envelope-marking-functoriality}
If $\tau\subseteq\sigma$ and $p\in\Thz_\sigma(T)$, there is a
canonical map $\mathsf F_\tau(p)\to\mathsf F_\sigma(p)$.
Precomposition identifies maps from the latter envelope into $K$ with
the full subcategory of maps $p\to\Theta K$ preserving the additional
$\sigma$-marked cocartesian arrows. In particular there are coherent
maps $\Omega_{T,\tau}\to\Omega_{T,\sigma}$ for the walking sources
constructed in Section~\ref{sec:thz:walking}.
\end{corollary}
\begin{proof}
\leavevmode
\begin{enumerate}
\item \textbf{Localize at the additional marked arrows.}
The inclusion $W_p^\tau\subseteq W_p^\sigma$ induces the map between
localized comma diagrams. The functor-category universal property of
localization identifies maps out of the second diagram with the full
subcategory inverting its additional generators. Apply enriched lex
reflection. Taking $p=\id_T$ gives the last assertion.
\qedhere
\end{enumerate}
\end{proof}

\subsection{The universal NU operation squares}\label{subsec:nu-doctrine}

We now specialize the envelope to the NU doctrine. Its named operations
must first be located in the universal finite-limit and comprehension
data. The squares below make their input dependencies explicit, so that
their internal interpretations can subsequently be obtained by the same
fiber calculation used for general operations.

We spell out the finite arity data in $\RNU$. This specifies the marking
without a closure convention that would inadvertently make every operation
strong. The universal span is
\[
 E\xrightarrow{\partial}U\xleftarrow{r}R,
 \qquad E_R=R\times_U E,
 \qquad \partial_R:E_R\to R.
\]
The map $\partial$ is representable. Its univalence is expressed by the
equivalence $|\id|:U\to\Eq(\partial)$, where the map
$\epsilon_\partial:\Eq(\partial)\to U\times U$ represents equivalences
between its fibers. The relation that $r$ is a monomorphism is expressed by
the equivalence $\Delta_r:R\to R\times_U R$.

For a representable $u:y\to z$, form
\begin{equation}\label{eq:nu-poly-objects}
 A_u=P_uz\times_z y,
 \qquad
 D_u=A_u\times_z y.
\end{equation}
In the second pullback the map $A_u\to z$ is evaluation, not the
projection to $z$. Denote the projections and evaluation by
\[
 \ell_u:A_u\to P_uz,\quad m_u:A_u\to y,\quad
 \mathrm{ev}_{u,z}:A_u\to z,\quad
 d_u:D_u\to A_u,\quad p_u^{\mathrm{poly}}:D_u\to y.
\]
Thus $u\otimes u:D_u\to P_uz$ is the composite $\ell_ud_u$. The map
$p_u^{\mathrm{poly}}$ is an arity projection and is distinct from either
of the term operations $p,p_R$ below.

The left exact universe operations are the following pullback squares:
\begin{equation}\label{eq:nu-lex-squares}
\begin{gathered}
\begin{tikzcd}[column sep=large]
 1\ar[r,"*"]\ar[d,"\id"']&E\ar[d,"\partial"]&
 D_\partial\ar[r,"p"]\ar[d,"\partial\otimes\partial"']&E\ar[d,"\partial"]\\
 1\ar[r,"\mathsf 1"']&U&
 P_\partial U\ar[r,"\Sigma"']&U
\end{tikzcd}\\[6pt]
\begin{tikzcd}[column sep=large]
 E\ar[r,"\mathrm{refl}"]\ar[d,"\delta"']&E\ar[d,"\partial"]\\
 E\times_U E\ar[r,"\mathsf{Id}"']&U.
\end{tikzcd}
\end{gathered}
\end{equation}
The selected representable types have the terminal and sum structures,
and the doctrine has the mixed dependent-product structure:
\begin{equation}\label{eq:nu-relative-squares}
\begin{gathered}
\begin{tikzcd}[column sep=large]
 1\ar[r,"*_R"]\ar[d,"\id"']&E_R\ar[d,"\partial_R"]&
 D_{\partial_R}\ar[r,"p_R"]\ar[d,"\partial_R\otimes\partial_R"']&E_R\ar[d,"\partial_R"]\\
 1\ar[r,"\mathsf 1_R"']&R&
 P_{\partial_R}R\ar[r,"\Sigma_R"']&R
\end{tikzcd}\\[6pt]
\begin{tikzcd}[column sep=large]
 P_{\partial_R}E\ar[r,"\lambda"]\ar[d,"P_{\partial_R}\partial"']&E\ar[d,"\partial"]\\
 P_{\partial_R}U\ar[r,"\Pi"']&U.
\end{tikzcd}
\end{gathered}
\end{equation}
The seventh operation square is the defining pullback with horizontal
arrows $\bar r:E_R\to E$ and $r:R\to U$ and vertical arrows
$\partial_R,\partial$. The seven squares include their coherent
commutativity homotopies.

\begin{proposition}[Cellular presentation of the doctrine]\label{prop:nu-cells}
Starting with the free theory on a representable $\partial:E\to U$,
freely impose its univalence and the three pullback structures in
\eqref{eq:nu-lex-squares}. Adjoin $r:R\to U$, invert $\Delta_r$, and
freely impose the three pullback structures in
\eqref{eq:nu-relative-squares}. The resulting theory is equivalent to
$\RNU$, compatibly with its universal structures. At either stage the
independent attachments may be performed in any order.
\end{proposition}
\begin{proof}
\leavevmode
\begin{enumerate}
\item \textbf{Represent each free attachment.} For fixed boundary arrows $v:A\to B$ and $u:E\to U$, the space of
pullback-square structures with vertical sides $v,u$ is the space of
horizontal arrows $a:A\to E$, $b:B\to U$ and a homotopy $ua\simeq bv$,
restricted to the locus where $A\to B\times_U E$ is an equivalence.
Freely adjoining arrows and homotopies, and then inverting this comparison,
represents exactly that space. The space of inverse witnesses of an arrow
is contractible when it is invertible and empty otherwise. These free
extensions exist in $\TT$ by its presentation constructions
\cite[Exs.~3.11--3.12]{NguyenUemura2025}.
\item \textbf{Identify the resulting universal doctrine.} In the free representable-arrow theory, the equivalence object and the
polynomial arities already exist. Theory morphisms preserve them, by their
finite-limit and dependent-product preservation. Therefore maps from the
first stage into a theory $V$ are exactly a representable univalent arrow
in $V$ with the three displayed left exact universe structures. After
adjoining $r$, its pullback $\partial_R$ is already representable, so every
remaining arity exists before any of the last four attachments. Maps from
the final theory are exactly the universal structures of
\cite[Def.~5.27]{NguyenUemura2025}. The represented functors agree;
Yoneda gives the asserted structured equivalence.
\item \textbf{Compare the orders of independent attachments.} At each stage take the coproduct of the independent free extensions in
the undercategory of the common initial stage. Mapping this coproduct into
$V$ gives the product of the extension spaces over the fixed initial map.
Iterated pushouts in any order give the same product, hence the same
universal object, through a contractible space of structured equivalences.
For two independent attachments this is the pushout
\[
\begin{tikzcd}[column sep=large]
 T_0\ar[r]\ar[d]&T_1\ar[d]\\
 T_2\ar[r]&T_1\amalg_{T_0}T_2.
\end{tikzcd}
\]
Its mapping space into $V$ is the pullback of the two extension spaces
over $\Map_{\TT}(T_0,V)$; the finite family follows by iteration.\qedhere
\end{enumerate}
\end{proof}

\subsection{Canonical comparisons and the structural marking}

Naming an operation does not yet decide which of its comparisons must be
invertible in an internal object. The marking is generated by the
specific structural comparisons attached to these squares. Describing
those generators positively makes the strong directions available for
transport and pasting without imposing unrelated invertibilities.

Polynomial base change applied to the square $(\bar r,r)$ gives the
canonical equivalence
$P_{\partial_R}U\simeq R\times_U P_\partial U$.
Accordingly the sum-arity comparison is the composite
\[
 \gamma_\Sigma:
 P_{\partial_R}R\xrightarrow{P_{\partial_R}(r)}
 P_{\partial_R}U\xrightarrow{\simeq}
 R\times_U P_\partial U\longrightarrow P_\partial U.
\]
The defining pullbacks of $A_u,D_u$ then give the induced comparison
$\widetilde\gamma_\Sigma:D_{\partial_R}\to D_\partial$ and the
relations
\begin{align}
 \varpi_{\partial,U}\gamma_\Sigma
   &\simeq r\varpi_{\partial_R,R},\label{eq:gamma-base}\\
 (\partial\otimes\partial)\widetilde\gamma_\Sigma
   &\simeq\gamma_\Sigma(\partial_R\otimes\partial_R),
     \label{eq:gamma-domain}\\
 \bar r p^{\mathrm{poly}}_{\partial_R}
   &\simeq p^{\mathrm{poly}}_\partial\widetilde\gamma_\Sigma.
     \label{eq:gamma-eval}
\end{align}
To compute the square in \eqref{eq:gamma-domain}, write
$X=P_{\partial_R}R$ and $P=P_\partial U$. The first comprehension
pulls back as follows:
\[
 X\times_P A_\partial
 \simeq X\times_U E
 \simeq X\times_R E_R=A_{\partial_R}.
\]
The map $X\to U$ in the middle is
$\varpi_{\partial,U}\gamma_\Sigma=r\varpi_{\partial_R,R}$.
Write $\alpha:A_{\partial_R}\to A_\partial$ for the resulting map.
In the polynomial adjunction the evaluation map is the transpose
of the identity on the polynomial. Postcomposition by $r$ in
$P_{\partial_R}(r)$ therefore postcomposes its evaluation by $r$,
and the base-change mate identifies evaluation after the pullback
of $\partial$. These are naturality of the counit and the
Beck--Chevalley identity of Lemma~\ref{lem:polynomial-bc}. Together
they give $\mathrm{ev}_{\partial,U}\alpha
\simeq r\,\mathrm{ev}_{\partial_R,R}$. The second comprehension is
thus
\[
 X\times_P D_\partial
 \simeq A_{\partial_R}\times_U E
 \simeq A_{\partial_R}\times_R E_R=D_{\partial_R},
\]
where the map to $R$ is now evaluation. Under this equivalence the
projection to $D_\partial$ is $\widetilde\gamma_\Sigma$.
Projecting the second pullback to $E$ gives
\eqref{eq:gamma-eval} with its homotopy. This constructs the
cartesian comparison together with the second evaluation boundary
used in the sum cube.

\begin{proposition}[Univalent terminal and sum comparison cubes]
\label{prop:nu-univalent-cubes}
The specified NU operation squares construct coherent comparisons
\begin{equation}\label{eq:nu-operation-comparisons}
 r\mathsf 1_R\simeq\mathsf 1,\qquad
 \bar r *_R\simeq *,\qquad
 r\Sigma_R\simeq\Sigma\gamma_\Sigma,\qquad
 \bar r p_R\simeq p\widetilde\gamma_\Sigma.
\end{equation}
For the terminal and sum operations, the paired classifier and term
comparisons are faces of complete coherently commutative comparison
cubes. They are natural in their entire parameter diagrams and
compatible with substitution and pasting.
\end{proposition}
\begin{proof}
\leavevmode
\begin{enumerate}
\item \textbf{Construct the terminal comparison.}
For the terminal comparison, the pullback squares classifying
$\mathsf 1_R$ and $\mathsf 1$ identify the respective pulled-back
$\partial$-families over $1$ with the terminal family. Their composite
is an equivalence of classified families over $1$; it carries the
selected total-space section $\bar r *_R$ to $*$. Univalence, in
its form $U\xrightarrow{\simeq}\Eq(\partial)$ over $U\times U$,
transports that equivalence to a path
$r\mathsf 1_R\simeq\mathsf 1$ together with the displayed term
boundary.


\item \textbf{Construct the sum comparison.}
For the sum comparison, put $X=P_{\partial_R}R$. The selected sum
square identifies $(r\Sigma_R)^*\partial$ over $X$ with
$D_{\partial_R}$. The ordinary sum square identifies
$(\Sigma\gamma_\Sigma)^*\partial$ with
$X\times_{P_\partial U}D_\partial$. The cartesian comparison
\eqref{eq:gamma-domain} identifies the latter with
$D_{\partial_R}$ over $X$, including its term projection by
\eqref{eq:gamma-eval}. Composing the two cartesian identifications
gives an equivalence of complete classified families
\[
 (r\Sigma_R)^*\partial
 \simeq(\Sigma\gamma_\Sigma)^*\partial
 \quad\text{over }X.
\]
Applying the inverse of the univalence equivalence to this entire
family equivalence gives $r\Sigma_R\simeq\Sigma\gamma_\Sigma$;
its induced comparison on total spaces is
$\bar r p_R\simeq p\widetilde\gamma_\Sigma$. 
\item \textbf{Fill the cube with its four side maps fixed.}
For a family $v:Y\to X$ in $\RNU$, cartesian squares $v\to\partial$
form the space
\[
 \operatorname{Cart}(v,\partial)\simeq
 \fib_v\!\left(\Map_{\RNU}(X,U)\longrightarrow
                    ((\RNU)_{/X})^\core\right),\qquad b\longmapsto b^*\partial .
\]
Indeed, over $b$ a square is a map $v\to b^*\partial$ in the slice,
and it is cartesian precisely when that map is an equivalence.
Univalence identifies paths between codes with equivalences of their
families: this is the fiber at the two codes of the map induced by
$U\simeq\Eq(\partial)$. Thus the displayed map of spaces is
$(-1)$-truncated, so every inhabited space
$\operatorname{Cart}(v,\partial)$ is contractible
\cite[Prop.~5.8(3)]{NguyenUemura2025}.

In the arrow category put $v=\id_1$ in the terminal case and
$v=\partial_R\otimes\partial_R$ in the sum case. There are two
cartesian maps $v\to\partial$: paste the selected operation with
$(\bar r,r)$, or paste the arity comparison with the ordinary
operation. The latter arity comparison is the identity for the terminal
operation and $(\widetilde\gamma_\Sigma,\gamma_\Sigma)$ for the sum.
The path space between the two composites in
$\operatorname{Cart}(v,\partial)$ is contractible. In the arrow
category a square with these four boundary maps is precisely a path
between their composites, up to the contractible choices of
composition. By currying, that square is a cube in $\RNU$. Evaluation
at the two endpoints recovers both terms of the corresponding pair in
\eqref{eq:nu-operation-comparisons}.

Pullback in the parameter sends these cartesian maps and their paths
to the corresponding ones for the reindexed family. Pasting gives
paths with the prescribed composite boundaries. Over each parameter, the required space is the path space between
the two specified points of $\operatorname{Cart}(v,\partial)$.
The preceding calculation makes every such fiber contractible, hence
the projection of their dependent sum is an equivalence. This gives
the coherent choice and all its higher compatibilities, including
those on common faces.\qedhere

\end{enumerate}
\end{proof}



\begin{definition}[Primitive structural arrows]\label{def:nu-generators}
Let $S_{\mathcal R}$ be the following finite list of named arrows in the
canonical construction:
\begin{align*}
 &\partial,\ \partial_R,\ \mathrm{pr}_U^i:U\times U\to U,
 \ \pi_E^i:E\times_U E\to E,\ \delta:E\to E\times_U E,\\
 &\pi_1,\pi_2:R\times_U R\to R,
 \ \Delta_U:U\to U\times U,\ \epsilon_\partial:\Eq(\partial)\to U\times U,\\
 &\varpi_{\partial,U},\ \varpi_{\partial_R,R},
 \ \varpi_{\partial_R,E},\ \varpi_{\partial_R,U},
 \ \partial\otimes\partial,\ \partial_R\otimes\partial_R,
 \ P_{\partial_R}\partial,\\
 &\ell_u,\ m_u,\ \mathrm{ev}_{u,z},\ d_u,\ p_u^{\mathrm{poly}}
       \qquad(u=\partial,\partial_R; i=1,2).
\end{align*}
All names refer to their constructions above. In particular $z=U$ for
$u=\partial$ and $z=R$ for $u=\partial_R$.
\end{definition}

\begin{lemma}[Canonical structural components]
\label{lem:canonical-realization}
Relative to the fixed universal NU structure, the complete named
arity, operation, and relation diagrams (including the comparison
cubes) have a contractible space of realizations compatible with
their stated universal constructions. Each structural name therefore
determines a well-defined component of $\Fun([1],\RNU)^\core$.
\end{lemma}
\begin{proof}
\leavevmode
\begin{enumerate}
\item \textbf{Construct the complete family in dependency order.}
Form the dependent family of complete diagrams by adjoining, in order,
the terminal and finite-limit objects, polynomial and equivalence
objects, the two evaluated polynomial pullbacks, the operation squares,
the Beck--Chevalley comparisons, and the cubes of
Proposition~\ref{prop:nu-univalent-cubes}. For the operation step,
the classifier of pullbacks of $\partial_R$ factors as
\[
 \Map(X,R)\longrightarrow\Map(X,U)
       \longrightarrow((\RNU)_{/X})^\core .
\]
The first map is $(-1)$-truncated by monicity of $r$, and the second
by univalence of $\partial$ \cite[Prop.~5.8(3)]{NguyenUemura2025}.
Their composite is $(-1)$-truncated. Hence for either vertical
classifier the space of cartesian squares with a given input family
is empty or contractible. Proposition~\ref{prop:nu-cells} supplies a
point for each of the six prescribed input families in the universal
doctrine. Their spaces of squares are therefore contractible. The
seventh square is the defining pullback of $\partial$ along $r$.
This derives the named operations in the universal theory.
For comparison, in finite sets with all maps representable,
$\varnothing\to1$ is univalent since its fiber has just the identity
automorphism. It admits no terminal operation square: such a square
would include $1\to\varnothing$. And $2\to1$ is representable with
two fiber automorphisms, so representability alone does not give
univalence. 
\item \textbf{Identify the contractible space of realizations.}
Terminal objects,
representations, adjoints with adjunction structure, and equivalent
finite-limit presentations have contractible spaces of realization,
by Lemma~\ref{lem:universal-choices}; the polynomial comparisons are
canonical mates. Proposition~\ref{prop:nu-univalent-cubes} uses the
inverse of a natural equivalence on the whole classified-family diagram,
so the cube and its boundary restrictions are one coherent output, not
separately supplied fillers. The successive dependent total spaces are
therefore contractible. 
\item \textbf{Extract the invariant arrow components.}
Evaluation at each named arrow gives a single
component in the arrow core, invariant under the resulting coherent
identifications.
\qedhere
\end{enumerate}
\end{proof}



\begin{definition}[NU arity marking]\label{def:nu-marking}
The marking $\tau_{\NU}$ on $\RNU$ is generated by the complete
arrow-core components of the names in $S_{\mathcal R}$. Thus it is the
smallest marking containing these structural arrows.
\end{definition}

This definition uses the generated marking of the specified structural
arrows. Pullback closure of that marking is a separate operation and is
not part of the definition. It should also be distinguished from marking
every pullback projection: the latter marks every arrow, since an arrow
is its own pullback along the identity of its codomain. The selected
dependent arities determine the generators used here.

\subsection{The finite named-cell shadow}\label{subsec:nu-shadow}

The finite shadow is a marked $\infty$-category recording the named NU
operations and their complete relation diagrams before completion as
an $\infty$-type theory. In particular, the shadow itself is not
asserted to possess finite limits, representable maps, or dependent
products. Its realization in $\RNU$ records the universal constructions
of these arities; its marked envelope supplies the left-exact categorical
object used subsequently.

Its named vertices are
\begin{gather*}
 1,E,U,R,E_R,E\times_U E,U\times U,R\times_U R,\Eq(\partial),\\
 P_\partial U,P_{\partial_R}R,P_{\partial_R}E,P_{\partial_R}U,
 A_\partial,D_\partial,A_{\partial_R},D_{\partial_R}.
\end{gather*}
The finite list of named complete cells consists of:
\begin{enumerate}
\item the terminal vertex, the product $U\times U$, the pullback
$E\times_U E$, and both defining pullbacks for each $A_u,D_u$;
\item the projection--diagonal triangles, the triangles
$u\otimes u\simeq\ell_ud_u$, and the naturality triangle
$\varpi_{\partial_R,U}P_{\partial_R}\partial
 \simeq\varpi_{\partial_R,E}$;
\item the seven NU operation squares;
\item the predicate pullback with sides $r,r,\pi_1,\pi_2$;
\item the walking equivalences $|\id|$ and $\Delta_r$, together with
$\epsilon_\partial|\id|\simeq\Delta_U$ and
$\pi_i\Delta_r\simeq\id_R$;
\item the polynomial comparisons
\eqref{eq:gamma-base}--\eqref{eq:gamma-eval} and the complete
comparison cubes of Proposition~\ref{prop:nu-univalent-cubes}.
\end{enumerate}
The terminal cell is a copy of $[0]$ sent to the named vertex $1$.
The occurrence shapes and their generating incidences form a finite
list; each shape retains its entire categorical nerve.

\begin{construction}[Coherent global cell attachment]
\label{con:incidence}
Let $I=I_{\mathcal R}$ be the category of named cells and their named
subfaces from the preceding list. A face carries its oriented occurrence
coordinates, its vertex and arrow labels, and the name of each specified
relation on it. Use one presentation for a shared named face. Arrows are
the actual face inclusions, with composite the nested inclusion. In
particular the two endpoints of a loop give two incidences before gluing.
Independent relation names with the same boundary give distinct cells.
Write $C_c^{\mathrm{occ}}$ for the nerve of the occurrence shape before
any names are identified. The shapes are vertices, intervals, the
indexing shapes of the indicated cones, triangles $[2]$, squares
$[1]^2$, cubes $[1]^3$, and the walking-isomorphism groupoid. In this
notation the nerve presents the corresponding object of $\Cat$.
The walking isomorphism is one shape with its specified generating faces.

Let $G$ be the directed graph with one vertex per vertex name and one edge
per nonidentity generator name. In a cell $c$, let $G_c$ record the
actual occurrences of its named vertices and generating edges, including
a named long edge of a composition triangle. Write
\[
 A_c=F(G_c),\qquad B_c=C_c^{\mathrm{occ}},\qquad H=F(G).
\]
Here $F$ is the free category on a graph, equivalently the categorical
realization of its one-dimensional simplicial set. The first map
$A_c\to B_c$ reads the edges in the complete occurrence shape. The map
$A_c\to H$ reads their global labels, sending identity-labelled occurrences
to identities. A displayed composite used as a path is read as that
path. Face inclusions respect both maps. Put
\begin{equation}\label{eq:coherent-cell-normalization}
 A=\colim_{c\in I} A_c,\qquad B=\colim_{c\in I} B_c,\qquad
 P_\infty=B\amalg_A H .
\end{equation}
All three colimits are in $\Cat$. If $j:H\to P_\infty$ is the pushout
map, give the shadow the generated marking
\begin{equation}\label{eq:finite-shadow}
 P_\infty^\tau=(P_\infty,\langle j(S_{\mathcal R})\rangle).
\end{equation}
Here $S_{\mathcal R}$ in the graph denotes the names of
Definition~\ref{def:nu-generators}. Its image in the doctrine will be
identified in Lemma~\ref{lem:coherent-cell-normalization}.

For every $V$ the full functor category is
\begin{equation}\label{eq:normalized-cell-mapping}
 \Fun(P_\infty,V)\simeq
 \Fun(H,V)\mathop{\times}_{\lim_{I^{\op}}\Fun(A_c,V)}
                    \lim_{I^{\op}}\Fun(B_c,V).
\end{equation}
Indeed, testing this formula against any category $Z$ applies the
mapping-space universal property of the pushout and the two incidence
colimits after taking their product with $Z$. That product preserves
colimits by cartesian closure of $\Cat$
\cite[Ex.~7.7.3.11]{Kerodon}. Yoneda proves the formula. The first
factor assigns each name once. The second gives the full relation
diagrams, with coherent restrictions on shared faces. The pullback
identifies the restrictions on all generator occurrences with the
single assignment. The same calculation with $Z=[n]$ exhibits the
noninvertible transformations and their higher compatibilities.

For example, take one name $x$ and a single edge $a:x\to x$. Then
$A_c\to B_c$ is an equivalence at both the vertex and edge
presentations. Hence $P_\infty\simeq H=B\mathbb N$, the free category
on one endomorphism. This calculation is sensitive to the order of
gluing. Identifying the endpoints inside the edge first and then taking
the incidence colimit would form the homotopy coequalizer of the two
identical maps $*\rightrightarrows B\mathbb N$. Mapping it into $V$
would give an endomorphism of $x$ and a free loop at $x$ in $V^\core$,
as follows by taking the homotopy equalizer of the two evaluations.
Equation~\eqref{eq:coherent-cell-normalization} performs the
name identification once.
\end{construction}

\begin{lemma}[Finite direct realization and marking]
\label{lem:coherent-cell-normalization}
\label{prop:incidence-marking}
The shadow is obtained from $H$ by finitely many relative latching
pushouts. The complete NU operation and relation diagrams determine a
marked realization $q:P_\infty^\tau\to\RNU^{\tau_{\NU}}$ through a
contractible space over those diagrams. Its image marking generates
$\tau_{\NU}$.
\end{lemma}
\begin{proof}
\leavevmode
\begin{enumerate}
\item \textbf{Compute the attaching category.}
The named generating faces form a finite direct category. A proper
inclusion increases face rank, and only the listed faces of the walking
isomorphism are used. Its full nerve is retained in $B_c$. For a lower
ideal $J\subseteq I$, put
$P_J=(\colim_J B_c)\amalg_{\colim_J A_c}H$.
Take $c\notin J$ with every proper face in $J$, and form
$L_cA$ and $L_cB$ by colimits over the nonidentity arrows into $c$.
Distinct embeddings of a common face remain distinct in this latching
index. Put $M_c=A_c\amalg_{L_cA}L_cB$. There is a pushout
\[
 P_{J\cup\{c\}}\simeq P_J\amalg_{M_c}B_c .
\]
To see this, a cone from $A$ on $J\cup\{c\}$ is a cone on $J$ and
a map from $A_c$ agreeing on $L_cA$, and the same holds for $B$.
Substitution of these two pushouts into the definition of
$P_{J\cup\{c\}}$ gives the formula by pasting. Equivalently, once
a functor $P_J\to V$ is fixed, its extension category is
$\Fun(B_c,V)\times_{\Fun(M_c,V)}\{M_c\to P_J\to V\}$.
This adds one filler per named relation with the specified common
faces. The initial stage is $P_\varnothing=H$ and the final stage is
\eqref{eq:coherent-cell-normalization}, independent of the order of
incomparable cells.

For a square with fixed side arrows, write $h,k:x\to y$ for the
two boundary composites. Its space of fillings is
$\operatorname{Path}_{\Map_V(x,y)}(h,k)$. Indeed, split the square
along its diagonal. A triangle on two prescribed consecutive edges
exhibits their composite, with a contractible space of choices.
When the diagonal $d$ is fixed, the two triangles accordingly give
paths $h\simeq d$ and $k\simeq d$. Summing over $d$ and composing
one path with the reverse of the other gives the displayed path
space. Consequently, two distinct named squares on those sides
contribute its product with itself. When two cubes share one named
square, their latching maps instead restrict to the single previously
assigned point of that space. The fiber formula above imposes these
identifications on the complete categories of diagrams.

\item \textbf{Realize the common boundary in the doctrine.}
Send each global generator to its named NU arrow. The vertex, cone
and polynomial cells are obtained by restriction of their terminal,
pullback and evaluation diagrams. The projection--diagonal triangles
are the corresponding cone identities. The operation squares are the
specified universe structures; the polynomial comparisons are the
Beck--Chevalley maps of Lemma~\ref{lem:polynomial-bc}. For a walking
equivalence use its named invertible arrow: restriction from the
walking isomorphism to its generating interval identifies functors out
of the former with the full locus of equivalences in the latter.
Proposition~\ref{prop:nu-univalent-cubes} constructs each comparison
cube in the arrow category with the prescribed operation squares as
faces. On a shared named face take the restriction of the same
diagram, with its specified homotopies. Functoriality of restriction
makes successive face restrictions compatible with composition in
$I$. Reading the generating edges of these diagrams recovers their
global labels with the identifications of
Lemma~\ref{lem:canonical-realization}. This gives a point of the right
side of \eqref{eq:normalized-cell-mapping} for $V=\RNU$, including
the comparison in the middle category.

\item \textbf{Compute the fiber of realization.}
Take cores in \eqref{eq:normalized-cell-mapping} and take the fiber
at that complete record. It is contractible because the displayed
map is an equivalence. This use of contractibility fixes the operation
fields and relation cells themselves; a fiber over boundary arrows
alone is the extension space computed in the first step. Finally
$qj$ sends the marked generator names precisely to the arrows used
in Definition~\ref{def:nu-marking}. Closure under equivalences and
composition gives $\langle q(\tau)\rangle_{\RNU}=\tau_{\NU}$.
\qedhere
\end{enumerate}
\end{proof}

\begin{proposition}[Separation in the finite shadow]
\label{prop:charge}
No primitive operation arrow, nor any of
$r,\bar r,\gamma_\Sigma,\widetilde\gamma_\Sigma$, is forced into the
finite marking $\tau$.
\end{proposition}
\begin{proof}
\leavevmode
\begin{enumerate}
\item \textbf{Assign degrees to the generating arrows.}
Let $B\mathbb N^5$ be the one-object category with endomorphism
monoid $\mathbb N^5$ with basis $\rho,\upsilon,\sigma,\iota,\kappa$.
Give structural arrows and relation equivalences degree zero and give
the other named arrows the following degrees:
\[
\begin{array}{c|c}
 \text{arrows}&\text{degree}\\\hline
 r,\bar r,\gamma_\Sigma,\widetilde\gamma_\Sigma&\rho\\
 \mathsf 1_R,*_R&\upsilon\\
 \mathsf 1,*&\rho+\upsilon\\
 \Sigma,p,\Sigma_R,p_R&\sigma\\
 \mathsf{Id},\mathrm{refl}&\iota\\
 \Pi,\lambda&\kappa.
\end{array}
\]

\item \textbf{Check the attaching maps.}
The table gives a functor on the global free category $H$.
The structural cone edges, projection--diagonal identities and
walking inverses have degree zero. The seven operation squares have
matching path degrees, respectively
$\rho+\upsilon,\sigma,\iota,\upsilon,\sigma,\kappa,\rho$
in the order terminal, sum, identity, selected terminal, selected
sum, mixed product and predicate inclusion. In the predicate
self-pullback each path has degree $\rho$. Both paths of each of
\eqref{eq:gamma-base}--\eqref{eq:gamma-eval} have degree $\rho$.
For the terminal comparison the two horizontal pastings have degree
$\upsilon+\rho=\rho+\upsilon$, with identity comparison on the
source family. For the sum comparison they have degree
$\sigma+\rho=\rho+\sigma$, with the arity and family comparison
edges both of degree $\rho$. The other faces of these cubes are
the operation and arity squares just checked. A diagram in the
nerve of the ordinary category $B\mathbb N^5$ is determined by its
edge assignments respecting composition. Thus all these checks give
compatible functors on the complete occurrences, including the
cubes, whose restrictions to $A_c$ read exactly the global degrees.
Equation~\eqref{eq:normalized-cell-mapping} gives a functor from
$P_\infty$ to $B\mathbb N^5$. Since the only invertible degree is
zero, the generated marking lies over zero. Every indicated operation
or comparison has nonzero degree.\qedhere

\end{enumerate}
\end{proof}

The positive degree grading is a property of this finite shadow, not
of the completed NU doctrine. For example, a terminal map $U\to1$
composed after $\mathsf1:1\to U$ gives an identity, preventing the
same positive grading on $\RNU$. No complementary description of
unmarked arrows of the completed doctrine follows from this grading.



\subsection{The microcosm and its dependent operations}

The arities and structural directions are now fixed. A marked section can
therefore be unpacked into the NU objects, dependent operations, and
comparison terms. The following identification checks that the compact
section formulation retains exactly those data and their morphisms.

\begin{definition}
For $K\in\Lex(\RNU,\Cat)$ define the completed and finite internal-object
categories by
\begin{align}
 \Int^{\NU}_\infty(K)
 &=\Fun^{\tau_{\NU}\text{-}\coc}_{/\RNU}
       (\RNU,\textstyle\int K),\label{eq:nu-internal}\\
 \Int^{\mathrm{cell}}_\infty(K)
 &=\Fun^{\tau\text{-}\coc}_{/\RNU}
       (P_\infty,\textstyle\int K).\label{eq:cell-internal}
\end{align}
Both are full subcategories of their section or relative-functor categories.
\end{definition}

\begin{theorem}[NU microcosm principle]\label{thm:nu-microcosm}
For a categorical $\infty$-type theory $K$, there are natural equivalences
\[
 \Int^{\NU}_\infty(K)\simeq\laxlim_{\RNU}^{\tau_{\NU}}K,
 \qquad
 \Int^{\mathrm{cell}}_\infty(K)\simeq\laxlim_{P_\infty}^{\tau}(Kq).
\]
Restriction to the finite shadow describes an internal object by its named
type, term, representable-type, and dependent-arity objects; its seven
operation cells; and the full coherent named relations and comparison
cubes. Structural comparisons are invertible on the positive list
$S_{\mathcal R}$. Polynomial arities have the comprehension interpretation
of Proposition~\ref{prop:polynomial-arity}. Every comparison between these
internal objects is retained.
\end{theorem}
\begin{proof}
\leavevmode
\begin{enumerate}
\item \textbf{Recover the marked section descriptions.} The two equivalences are Theorem~\ref{prop:marked-sections}. Taking
\eqref{eq:normalized-cell-mapping} over the NU realization identifies
a finite-shadow object with one global generator assignment and a
compatible family of full occurrence diagrams over its restrictions.
The named vertices give the indicated objects. A named arrow $f:s\to t$
gives $\chi_f:K(f)X_s\to X_t$, and the marking makes this comparison
invertible on the structural arrows. The normalized functor diagrams give
the operation, relation, and cube coherences together.
\item \textbf{Interpret the dependent comparisons.} For the dependent interpretation, consider a named square $vh\simeq ku$
with structural vertical sides $u,v$. Its cone coherence is $\chi_v\,K(v)(\chi_h) \simeq \chi_k\,K(k)(\chi_u),$ with the base-square comparison understood. For a pullback square with $h:A\to B$ and $k:C\to D$,
left exactness gives, at $c\in K(C)$,
\[
 K(A)\times_{K(C)}\{c\}\simeq K(B)\times_{K(D)}\{K(k)c\}.
\]
For the arrow comparison put
$\theta:K(v)K(h)X_A\simeq K(k)X_C$, the base-square equivalence
followed by $K(k)\chi_u$. The map on the top mapping space is
\[
 e\longmapsto\chi_v\circ K(v)(e)\circ\theta^{-1}
 \quad\text{in }\Map_{K(D)}(K(k)X_C,X_D).
\]
The specified square makes $\chi_h$, together with its coherence,
a point of the homotopy fiber at $\chi_k$. These are precisely the
displayed fibers of Proposition~\ref{thz:prop:displayed},
with its horizontal and vertical pasting laws.
Apply this first to $(\bar r,r)$ and then to the polynomial evaluation
and comparison squares. Proposition~\ref{prop:polynomial-arity} identifies
their complete Segal fibers with evaluation at comprehension contexts.
The polynomial mates and the comparison cubes are those already
constructed, so pasting introduces no additional preservation fields.
\item \textbf{Identify the sum operation.} For example the strong comparison over
$\varpi_{\partial,U}:P_\partial U\to U$ places $X_{P_\partial U}$
in the fiber $\Dep^K_{\partial,U}(X_U)$. The operation $\Sigma$ gives
$K(\Sigma)X_{P_\partial U}\to X_U$, and the comparison over its term
arrow $p$ lies over that operation by the sum square. The other dependent
operations are read from their displayed squares in exactly the same way.
No such dependent input fiber is declared contractible. Finally, each
step is an equivalence of full functor categories or their full marked
subcategories, so the same description includes every morphism and higher
cell.\qedhere
\end{enumerate}
\end{proof}

\begin{proposition}[Marked descent]\label{prop:marked-descent}
If $(D,W)=\operatorname*{colim}_{i\in I}(D_i,W_i)$ in marked categories over
$T$, then for every $K:T\to\Cat$,
\[
 \Fun^{W\text{-}\coc}_{/T}(D,\textstyle\int K)
 \simeq\lim_{i\in I^{\op}}\Fun^{W_i\text{-}\coc}_{/T}(D_i,\textstyle\int K).
\]
For the finite shadow use the occurrence cells and the global graph
in \eqref{eq:normalized-cell-mapping}. Completed objects are the
limit over the opposite of the marked simplex category of $\RNU$.
\end{proposition}
\begin{proof}
\leavevmode
\begin{enumerate}
\item \textbf{Descend functors and markings along the colimit.} The functor category out of a colimit is the limit of the functor
categories out of its terms. A functor carries the generated marking to
cocartesian arrows exactly when it does so on every local marked arrow,
since cocartesian arrows are closed under composition and equivalence.
The assertion follows by full restriction. Explicitly put
$S(V)=\Fun_{/\RNU}(V,\int K)$ for an occurrence or graph over the
doctrine, and let $S_\tau(H)$ impose cocartesianness on the global
structural generators. Then the finite category of internal objects is
\[
 S_\tau(H)\times_{\lim_{I_{\mathcal R}^{\op}} S(A_c)}
                     \lim_{I_{\mathcal R}^{\op}} S(B_c).
\]
Apply \eqref{eq:normalized-cell-mapping} to $\int K$ and $\RNU$
and take the fiber at $q$. Commutation of limits gives the displayed
pullback before marking, and full restriction on its first factor
imposes precisely the generated marking. This computation retains
all vertical modifications. For the completed resolution each marked
arrow occurs as a marked one-simplex, and
Lemma~\ref{thz:lem:marked-resolution} proves the required density.\qedhere
\end{enumerate}
\end{proof}

\begin{proposition}[Comparison with the completion]\label{prop:shadow-comparison}
Restriction defines
$\Int^{\NU}_\infty(K)\to\Int^{\mathrm{cell}}_\infty(K)$.
For a fixed diagram $K$, its homotopy fiber at a finite object is precisely
the category of its marked extensions along $q$. If $q$ is marked coinitial,
meaning that it induces an equivalence of marked lax limits for every
category-valued diagram on the target, restriction is an equivalence.
For left exact macrocosms alone, the necessary and sufficient criterion is
the envelope comparison of Theorem~\ref{thz:thm:shadow} below.
\end{proposition}
\begin{proof}
\leavevmode
\begin{enumerate}
\item \textbf{Identify restriction fibers and apply the envelope criterion.} The fiber description is the definition of restriction in the section
category. The marked-coinitial assertion follows by applying its universal
property to $K$. The envelope criterion follows from its representing
property and Yoneda, as proved below.\qedhere
\end{enumerate}
\end{proof}

For clarity, adjoining a single source object $p$ and an arrow $p\to z$
to an existing diagram with specified value $X_z$ has extension category
$K(p)\times_{K(z)}K(z)_{/X_z}$. If the new arrow is required to be
strong, this becomes $K(p)\times_{K(z)}\{X_z\}$. These fibers need
not be contractible or inhabited. The relative cellular completion in
Theorem~\ref{thm:relative-cellular-completion} below computes the
entire extension diagram rather than replacing its weak fillers by
universal-choice claims.

\subsection{The walking internal theory}\label{sec:thz:walking}

Having identified the section data, we can package them into a single
walking source. Applying the envelope to the identity of the marked
doctrine makes its universal property the internal-object functor itself.
The resulting representability is what will support the higher
categorical constructions developed later in the paper.

The identity object $\one_{\Thz}=(T\xrightarrow{\id}T)$ belongs to $\Thz_\tau(T)$.  Every arrow is cocartesian, and its lifted
finite cones are exactly the limiting cones in $T$.  It is terminal:
for $p:E\to T$, the unique functor to $\id_T$ over $T$ is $p$ itself,
and it preserves the required marked cocartesian arrows.  The corresponding
hom $\infty$-category is contractible.  The marking used by the envelope
on this terminal object is precisely $\tau$.

\begin{theorem}[Walking internal theory]\label{thz:thm:walking}
Define $\Omega_{T,\tau}:=\Env_T(\id_T,\tau) =L_{\mathrm{lex}} \bigl(t\mapsto (T_{/t})[W_t^{-1}]\bigr).$ Then $\Omega_{T,\tau}$ represents the marked internal-object functor:
\begin{equation}\label{thz:eq:walking-UP}
 \Nat_{\mathcal L_T}(\Omega_{T,\tau},K)
 \simeq
 \operatorname{Hom}_{\Thz_\tau(T)}
   (\one_{\Thz},\Theta K)
 \simeq
 \operatorname*{laxlim}\nolimits^{\tau}_T K.
\end{equation}
The identity of $\Omega_{T,\tau}$ corresponds to a generic marked section $\eta:T\longrightarrow\int_T\Omega_{T,\tau}.$ Every marked internal object in a macrocosm $K$ is the specialization of
this section along its classifying map $\Omega_{T,\tau}\to K$, with the
complete category of classifying maps given by
\eqref{thz:eq:walking-UP}.
\end{theorem}
\begin{proof}
\leavevmode
\begin{enumerate}
\item \textbf{Represent the marked section category.} A functor over $T$ from $\id_T$ to $\int_TK\to T$ is exactly a
section.  The marked-morphism condition says that its edge over every
$f\in\tau$ is cocartesian.  Theorem~\ref{thz:thm:envelope} therefore
gives the first equivalence.  Theorem~\ref{prop:marked-sections} identifies the
marked section category with the partially lax limit.  The identity of the
representing object gives $\eta$ by Yoneda, and naturality gives all
specializations and their morphisms.\qedhere
\end{enumerate}
\end{proof}

\begin{corollary}[NU specialization]\label{thz:cor:NU}
For $(T,\tau)=(\RNU,\tau_{\NU})$, write
$\Omega_{\NU}=\Omega_{T,\tau}$.  Then $\Nat_{\Lex(\RNU,\Cat)}(\Omega_{\NU},K) \simeq \Int^{\NU}_\infty(K).$ Under Theorem~\ref{thm:nu-categorical}, the object $\Omega_{\NU}$ is an object of
$\operatorname{Cat}(\TT)$ carrying the generic internal $\infty$-type
theory represented by this equivalence.
\end{corollary}
\begin{proof}
\leavevmode
\begin{enumerate}
\item \textbf{Specialize the walking envelope to the NU doctrine.} Theorem~\ref{thm:nu-microcosm} identifies the partially lax limit in
\eqref{thz:eq:walking-UP} with $\Int^{\NU}_\infty(K)$ for the
completed NU marking.  Theorem~\ref{thm:nu-categorical} identifies
$\Lex(\RNU,\Cat)$ with $\operatorname{Cat}(\TT)$.\qedhere
\end{enumerate}
\end{proof}

\begin{corollary}[Limits of internal theories]\label{thz:cor:limits}
On ordinary natural transformations of macrocosms,
$\Int^{\NU}_\infty$ preserves every bounded small limit.  More
generally, for fixed $p\in\Thz_\tau(T)$, the functor $K\longmapsto \operatorname{Hom}_{\Thz_\tau(T)}(p,\Theta K)$ preserves bounded small limits.
\end{corollary}
\begin{proof}
\leavevmode
\begin{enumerate}
\item \textbf{Commute enriched hom with limits.} Limits in $\mathcal L_T$ are computed pointwise because limits in
$\Cat$ commute with finite limits.  For fixed $A$, the enriched hom
$\Nat(A,-)$ preserves limits.  On complete Segal level $n$ this is the
preservation of limits by the representable mapping-space functor
after cotensoring by $[n]$.  Apply
Theorems~\ref{thz:thm:envelope} and \ref{thz:thm:walking}.\qedhere
\end{enumerate}
\end{proof}

The generic section has both a record and a semantic description.  As a
record it consists of one fiber object over every object of $T$, one
operation over every arrow, invertibility of the operation comparison on
$\tau$, and the complete coherent composition diagram.  Semantically it
is the image of $\id_{\Omega_{T,\tau}}$ under
\eqref{thz:eq:walking-UP}.  The generic-arrow calculation below identifies its structural terms under specialization.

\subsection{The classifier and its structural terms}

The envelope identifies classifying maps with marked cones. To identify
their mathematical content, we trace a distinguished comma arrow through
localization and reflection. Its image is the generic comparison term;
specialization recovers the original cone comparison with the correct
boundaries.

For $a:D^W\to T$ set
\begin{equation}\label{eq:envelope}
 \Omega_{a,W}=\Env_T(a,W)
 =L_{\mathrm{lex}}\bigl(t\mapsto(a\downarrow t)[W_t^{-1}]\bigr).
\end{equation}
Theorem~\ref{thz:thm:envelope} and
Theorem~\ref{prop:marked-sections} give the following explicit form.

\begin{proposition}[Classifier with explicit inverse]\label{prop:classifier}
There is a natural equivalence $\Nat(\Omega_{a,W},K)\simeq\laxlim_D^W(Ka).$ A cone $(x_d,\ell_u)$ is classified by the reflection transpose of
\begin{equation}\label{eq:classifyformula}
 (a\downarrow t)[W_t^{-1}]\longrightarrow K(t),
 \qquad(d,h)\longmapsto K(h)x_d.
\end{equation}
On the comma arrow $(d,h'a(u))\to(d',h')$, this functor applies
$K(h')$ to $\ell_u$. Conversely, evaluation at $(d,\id_{ad})$ and
at $(d,a(u))\to(d',\id_{ad'})$ recovers the cone and its arrows.
\end{proposition}
\begin{proof}
\leavevmode
\begin{enumerate}
\item \textbf{Compose the explicit comma and reflection correspondences.} Apply the explicit equivalence of full hom categories in
Lemma~\ref{thz:lem:comma}, followed by
Lemma~\ref{thz:lem:lex-reflector}. The displayed formula records that
composite on objects and arrows.\qedhere
\end{enumerate}
\end{proof}

For the smallest marking $\Eq(D)$ write $\Omega_a^-=\Omega_{a,\Eq(D)}.$ The identity of $\Omega_a^-$ classifies a universal cone $s$. For every arrow $u:d\to d'$ of $D$ it has the arrow
\[
 \chi_u:\Omega_a^-(a(u))s_d\longrightarrow s_{d'}
 \quad\text{in }\Omega_a^-(ad').
\]
This arrow is the image of the distinguished comma arrow under the localization and reflection units. In particular it is not specified by an abstract identification of categories of cones.

\begin{lemma}[The term comparison]\label{lem:term}
Let $f:\Omega_a^-\to K$ classify an unmarked lax cone $x$. There is a natural boundary-compatible identification
\begin{equation}\label{eq:termcompare}
 f(\chi_u)\simeq\ell_u(x).
\end{equation}
The source boundary uses the naturality equivalence $f_{ad'}\Omega_a^-(a(u))s_d \simeq K(a(u))f_{ad}s_d.$ Consequently
\begin{equation}\label{eq:strongness}
 \prod_{u\in W}\isEq(f(\chi_u))
 \simeq
 \prod_{u\in W}\isEq(\ell_u(x)).
\end{equation}
The products are over the complete marked-arrow family, with its dependencies on endpoints.
\end{lemma}
\begin{proof}
\leavevmode
\begin{enumerate}
\item \textbf{Evaluate the universal comparison arrow.} The inverse construction in Proposition~\ref{prop:classifier} evaluates the classifying transformation on precisely the comma arrow that defines $\chi_u$. Its image is therefore $\ell_u(x)$, with the displayed boundary comparison supplied by naturality. The unit and counit of the localization and reflection adjunctions give the same identification after passing to $\Omega_a^-$. All these comparisons are natural in the full parameter record $(u,f)$.
\item \textbf{Identify structural strongness.} Both selected-arrow families are now sections of the arrow category of $K(ad')$ with the same boundaries. Their term comparison is the identity/equivalence family of \eqref{eq:termcompare}. Pulling back equivalence membership proves \eqref{eq:strongness}.
\item \textbf{Express the comparison through cocartesian factorization.} Equivalently one may retain the complete cocartesian-factorization
parameter: a cocartesian arrow $c:x_d\to Y$, a vertical arrow
$v:Y\to x_{d'}$, and an identification $vc\simeq x(u)$.
Its universal property gives $e:K(a(u))x_d\simeq Y$ and
$ve\simeq\ell_u(x)$. This identifies the arrow terms and their
boundaries. Forgetting the factorization gives one comparison of the
parameter families; in the reverse direction its inhabited contractible
space supplies existence, without choosing a cleavage.\qedhere
\end{enumerate}
\end{proof}

\begin{proposition}[Marking is localization of the classifier]\label{prop:marklocal}
The canonical map $q_W:\Omega_a^-\longrightarrow\Omega_{a,W}$ has the following enriched universal property: for every $K\in\LT$, precomposition is fully faithful and identifies $\Nat(\Omega_{a,W},K)$ with the full subcategory of $\Nat(\Omega_a^-,K)$ on those $f$ for which every $f(\chi_u)$, $u\in W$, is an equivalence.
\end{proposition}
\begin{proof}
\leavevmode
\begin{enumerate}
\item \textbf{Identify localization with the strong-cone inclusion.} Under Proposition~\ref{prop:classifier}, precomposition is the inclusion of $W$-strong cones among all lax cones. That inclusion is a full subcategory: strongness is a condition on objects, and no extra condition is imposed on comparisons between them. Lemma~\ref{lem:term} identifies its object predicate with the displayed predicate on classifying maps. This proves the assertion, including all higher cells.\qedhere
\end{enumerate}
\end{proof}

Thus the structural directions have been compiled into a universal source object. This is why ordinary natural transformations out of $\Omega_{a,W}$ classify a microcosm with lax operation directions.

\subsection{The two mapping properties of the envelope}

\begin{proposition}[Internal objects are not preserved by the envelope unit]
\label{prop:envelope-internal-object-obstruction}
For a marked finite-limit doctrine there can be
$p\in\Thz_\tau(T)$ with
\[
 \Int_\tau(p)=\varnothing,
 \qquad
 \Int_\tau(\Theta\mathsf F_\tau(p))\simeq\one.
\]
\end{proposition}
\begin{proof}
\leavevmode
\begin{enumerate}
\item \textbf{Construct the marked operation theory.}
Take a small skeleton $T=\mathbf{FinSet}_*$, let
$\tau=\Eq(T)$, and let $p:\one\to T$ select the one-point pointed
set $0$. The object $0$ is both initial and terminal. A finite diagram
in $\one$ maps to a constant diagram at this terminal object; every
limiting cone lifts uniquely and is limiting. The equivalence-marked
lift condition also holds. Thus $p$ is an object of $\Thz_\tau(T)$.


\item \textbf{Compute its represented envelope.}
For every $t$, the comma category $p\downarrow t$ is terminal because
there is exactly one pointed map $0\to t$. Its marking has only the
identity. Hence $R_{p,W_p}$ is the constant terminal diagram, already
left exact, and
\[
 \mathsf F_\tau(p)=\one_T,
 \qquad \Theta\mathsf F_\tau(p)\simeq\id_T.
\]

\item \textbf{Compute the two intrinsic internal-object categories.}
A section of $p$ cannot exist: there is no object over a nonzero
pointed finite set. Therefore $\Int_\tau(p)=\varnothing$.
The identity $T\to T$ has a contractible category of sections over
itself, giving the second assertion.
\qedhere
\end{enumerate}
\end{proof}

The two calculations use different mapping categories. The envelope
represents maps from its source into represented targets. Intrinsic
internalization computes maps from $\id_T$ into the source itself.
Keeping the original source, as in
Definition~\ref{def:general-internalization}, retains that second
mapping category.

\section{Relative microcosms and universal classification}\label{sec:relative}
The preceding microcosm is internal to the NU doctrine itself. A native
theory of operations may instead be indexed by a different finite-limit
category $T$. Its envelope already represents the desired internal
objects in $\Lex(T,\Cat)$; the remaining question is how to place that
entire category inside the NU framework without losing its indexing data.
We first establish why an unrestricted absolute classification cannot
serve this purpose.

For a space-valued model, the category of elements records an element
together with its arity. Applying this construction to every complete
nerve level of a categorical model retains both its objects and its
noninvertible morphisms. The projection to $T$ survives as a relative
base. This provides the fully faithful realization used in the universal
relative classification theorem below.

For a small finitely complete $\infty$-category $T$, put
\[
 \Th(T)=\Lex(T,\Sps),\qquad \LT=\Lex(T,\Cat).
\]
The enrichment of $\LT$ has ordinary natural transformations as objects of $\Nat(A,K)$, and modifications as its morphisms. In particular, the naturality comparisons in an object of this hom category are invertible. Laxity will be encoded in the \textbf{source classifier}.

Write $T_{\minrep}$ for $T$ with equivalences as its representable maps. Let
\[
 \CNU=\operatorname{Cat}(\TT)\simeq\Lex(\RNU,\Cat),
 \qquad \B_T=\disc(T_{\minrep}).
\]
Here $\disc(T_{\minrep})$ is the constant simplicial object on the \textbf{theory} $T_{\minrep}$. Under the NU equivalence it is the space-valued model corresponding to $T_{\minrep}$, regarded as category-valued. It is not a constant functor with value the category $T$.

The notation $\Homc_{\CNU/\B_T}(A,B)$ denotes the full category of internal functors over $\B_T$ and internal natural transformations over that base. The base is fixed by a specified invertible comparison, as in a homotopy slice.
A classified internal object is consequently a triangle
\[
\begin{tikzcd}[column sep=large]
 \Om_{a,W}\ar[rr,"F"]\ar[dr]&&\cK\ar[dl]\\
 &\B_T&
\end{tikzcd}
\]
together with its specified comparison over $\B_T$. Morphisms of
classified objects are internal natural transformations over that same
base.


\subsection{Why the relative base is necessary}

The relative base and source classifier carry information that can
be detected by internal objects. The canonical object of an absolute
NU microcosm gives an obstruction to classifying arbitrary marked
lax limits without retaining that base.

\begin{lemma}[The zero theory]\label{lem:zero}
The terminal category with its minimal representable-map structure is both initial and terminal in $\TT$. Hence, under $\Th(\RNU)\simeq\TT,$ the terminal space-valued model $\one$ is also initial.
\end{lemma}
\begin{proof}
\leavevmode
\begin{enumerate}
\item \textbf{Verify terminality and initiality.} There is a unique theory morphism from any theory to the terminal theory. A left-exact functor from the terminal category to a theory $A$ selects a terminal object of $A$, and the space of terminal objects is contractible. All arrows of the source are equivalences. They remain representable in $A$, and their dependent-product mates are invertible as in Lemma~\ref{lem:minimal}. Thus the functor is a theory morphism, uniquely up to a contractible space. This proves initiality. Equivalences of $\infty$-categories preserve initial and terminal objects, and the terminal object of $\Lex(\RNU,\Sps)$ is the constant point-valued model.\qedhere
\end{enumerate}
\end{proof}

\begin{theorem}[Canonical internal object]\label{thm:pointed}
For every $K\in\Lex(\RNU,\Cat)$ and every Berman marking $\tau$ of $\RNU$, the category $\laxlim_{\RNU}^{\tau}K$ has a canonical initial object,
natural under ordinary natural transformations of $K$. In particular
$\Int^{\NU}_\infty(K)$ has an initial object.
\end{theorem}
\begin{proof}
The representable $h_1=\Map_{\RNU}(1,-)$ is left exact and initial in
$\Lex(\RNU,\Sps)$: by Yoneda its maps to $F$ form $F(1)\simeq *$.
Lemma~\ref{lem:zero} makes the constant point initial there as well.
Hence $h_1\simeq *$, so the terminal object $1$ of $\RNU$ is also initial.

More generally, let $i$ be initial in $T$ and $K:T\to\Cat$.
Write $i_t:i\to t$ for its coherent cone: the projection
$T_{i/}\to T$ is the left fibration of the contractible-valued
functor $\Map_T(i,-)$ and hence an equivalence. Transport defines
the strong section $L(y)_t=K(i_t)y$ for $y\in K(i)$.
On the category of lax sections strong on any marking,
$L$ is left adjoint to evaluation at $i$. Its unit is the identification
$y\simeq K(i_i)y$. Its counit at a section $x$ has component
$\chi_{i_t}(x):K(i_t)x_i\to x_t$. This is a modification: at $f:s\to t$
the section's composition law on $f i_s\simeq i_t$ is exactly its
naturality square. Applying the section to the entire cone gives this
square on every simplex; the same operation on a family of sections
gives naturality in $x$. At $i$ the counit is the identity, and on
$L(y)$ it is the transport identification for the same cone. These
are the two triangle identities. In particular
\[
 \Map(L(y),x)\simeq\Map_{K(i)}(y,x_i).
\]
Now $K(1)\simeq *$ by left exactness. Its object gives the asserted
initial section, and naturality in $K$ follows by transport along the
same $i_t$.
\end{proof}

\begin{theorem}[Obstruction to unrestricted absolute reduction]\label{thm:noabsolute}
There is no construction assigning to every bounded marked diagram $(D,W,C)$ a categorical NU theory $K_C$ such that $\laxlim_D^W C\simeq\Int^{\NU}_\infty(K_C)$ for all such diagrams. The assertion already fails among categorical models of the Lawvere theory of plain objects. The obstruction already occurs for a fixed nonempty category, independently of functoriality in that category.
\end{theorem}
\begin{proof}
\leavevmode
\begin{enumerate}
\item \textbf{Rule out an unrestricted absolute realization.} Let $L=\mathrm{FinSet}^{\op}$ and let
\[
 C^{\bullet}:L\longrightarrow\Cat,\qquad n\longmapsto C^n.
\]
This is a finite-product-preserving categorical model of the Lawvere theory of plain objects. Every arrow belongs to its structural finite-power marking. The category $L$ has initial object $1$, so
\[
 \laxlim\nolimits_L^{\sharp}C^{\bullet}
 \simeq\lim_LC^{\bullet}\simeq C.
\]
Equivalently a strong cone is determined by its component at $1$, and its remaining components are the corresponding tuples. For $C=\varnothing$ the result is empty. Theorem~\ref{thm:pointed} rules out the asserted NU realization, even without any naturality requirement.
\item \textbf{Treat a fixed nonempty carrier.}
Take the discrete category $C=\{0,1\}$. Neither object is initial,
since there is no arrow to the other. Theorem~\ref{thm:pointed}
therefore precludes an equivalence for this individual $C$.\qedhere
\end{enumerate}
\end{proof}


The relative hom category in Theorem~\ref{thm:main} imposes a specified
map to $\B_T$ and a specified source classifier. A zero map need not
commute with those maps to the base. Consequently the canonical-object
argument for an unrestricted NU microcosm does not apply to this slice
hom category. 


\paragraph{An unrepresented theory at the universal marking.}
The point $e:*\to\RNU$ at $1$ belongs to $\Thz_{\tau_{\NU}}(\RNU)$.
To check its marked lifts, use the NU model of finite sets with
\emph{all} maps representable, which exists because finite sets are
locally cartesian closed. At the terminal context both $U$ and $R$
classify finite sets, and $E,E_R$ classify pointed finite sets.
No domain of a generator in Definition~\ref{def:nu-generators} is
equivalent to $1$. Indeed in $U,R,U\times U,R\times_U R$ use sets
of cardinalities one and two (in each factor as required).
In $E,E_R,E\times_U E$ mark a point, twice the same point in the last
case. In $\Eq(\partial)$ use their identity equivalences. In the
polynomial domains $P_u z,P_{\partial_R}U,P_{\partial_R}E$ vary the
outer set in the same way and use singleton fibers, pointed for the
parameter $E$. In $A_u,D_u$ also select the outer point and, where
required, the unique point of the corresponding fiber. These are
elements by the comprehension interpretation of
Proposition~\ref{prop:polynomial-arity}. Projection remembers
the differing outer cardinalities throughout. A lex model sends $1$
to a contractible space, proving the domain assertion. Hence a word
in the structural generators and equivalences starting at $1$ consists
entirely of equivalences. These admit cocartesian lifts from the point.
Finite cocartesian exactness holds because every diagram in $*$ maps
to a diagram constantly terminal and every limiting cone on it is
equivalent to the terminal cone.

Now $\mathsf1:1\to U$ has empty target fiber for $e$. Its outgoing
row from the unique source object cannot be corepresented. Also $e$
admits no section, while $e\downarrow t\simeq\Map_{\RNU}(1,t)\simeq *$
by Theorem~\ref{thm:pointed}. Thus its represented envelope is the constant terminal model.
Enlarging the marking to include $\mathsf1$ would leave finite exactness
intact while destroying marked lift existence. Taking constant external simplicial objects
gives the same separation at every finite categorical stage.

\subsection{Universal relative classification}

\begin{theorem}[Universal relative classification]\label{thm:main}
Let $a:D\to T$ be any functor from a bounded small $\infty$-category to a small finitely complete $\infty$-category, let $W$ be a Berman marking of $D$, and let $K\in\LT$. There are constructed objects
\[
 \Omega_{a,W}=\Env_T(a,W)\in\LT,
 \quad \Om_{a,W}=\J_T\Omega_{a,W}\longrightarrow\B_T,
 \quad \cK=\J_TK\longrightarrow\B_T
\]
and natural equivalences
\begin{align}
 \laxlim\nolimits_D^W(Ka)
 &\simeq \Nat_{\LT}(\Omega_{a,W},K)\label{eq:main1}\\
 &\simeq \Homc_{\CNU/\B_T}(\Om_{a,W},\cK).
 \label{eq:main2}
\end{align}
The functor $\J_T$ is fully faithful on these hom $\infty$-categories. It is constructed from categories of elements, preserves bounded small limits relative to $\B_T$, and admits an explicit inverse on its essential image.

For an arbitrary bounded marked category $(D,W)$ and an arbitrary diagram $C:D\to\Cat$, these conclusions apply canonically with
\[
 T=QD=D^{\ff},\qquad a=i_D:D\to QD,\qquad K=\overline C,
 \qquad \overline C\,i_D\simeq C.
\]
Thus every bounded marked lax-limit microcosm has this relative NU realization. No compact-generation condition or structural marking on the whole completion $QD$ is required.
\end{theorem}

\subsection{\texorpdfstring{Categories of elements as $\infty$-type theories}{Categories of elements as infinity-type theories}}

The principal new ingredient works first with space-valued models. Recall that a covariant functor $F:T\to\Sps$ is classified by a \textbf{left} fibration. We use the covariant convention throughout this section.

\begin{lemma}[Minimal theories and limits]\label{lem:minimal}
The assignment $T\mapsto T_{\minrep}$ is fully faithful from the $\infty$-category of small finitely complete $\infty$-categories and left-exact functors to $\TT$. Its image is closed under bounded small limits, and those limits are created on underlying categories.
\end{lemma}
\begin{proof}
\leavevmode
\begin{enumerate}
\item \textbf{Embed the minimal theories fully faithfully.} Equivalences
are exponentiable and are stable under pullback and composition. Let
$H:T\to T'$ be left exact and let $u:y\to z$ be an equivalence. Then
$Hu$ is an equivalence as well. Pullback along $u$ is an equivalence
of slices with inverse $\Sigma_u$, so its right adjoint satisfies
$u_*\simeq\Sigma_u$. Likewise $(Hu)_*\simeq\Sigma_{Hu}$. Under these
identifications the mate \eqref{eq:pushforward-mate} is the canonical
comparison $H_z\Sigma_u\to\Sigma_{Hu}H_y$ for postcomposition.
At $v:a\to y$ it is the equivalence in the slice over $Hz$ induced
by $H(uv)\simeq(Hu)(Hv)$. Hence $H$ preserves dependent products
along every representable arrow of $T_{\minrep}$, proving full
faithfulness.
\item \textbf{Construct the underlying limits.} For a diagram of finitely complete categories and left-exact functors, the limit of the underlying categories has finite limits computed componentwise. Its projection functors preserve them. A cone of left-exact functors factors through this limit by a left-exact functor. In the corresponding limit of representable-map structures an arrow is distinguished exactly when all its projections are distinguished. When all the structures are minimal, this says exactly that all projections are equivalences, or equivalently that the arrow is an equivalence.
\item \textbf{Test the limit against arbitrary theories.} To check the universal property against an arbitrary theory $S$, a cone
of theory morphisms $S\to T_i$ induces a left exact functor into the
underlying limit. It sends a representable arrow to an equivalence because
every projection does. Its dependent-product mate is invertible after
every projection, and the projections jointly detect equivalences, so it
is a theory morphism. This proves the limit universal property in $\TT$,
including cones whose source is not minimal. The empty diagram gives the
terminal minimal theory.\qedhere
\end{enumerate}
\end{proof}

\begin{construction}[Element realization]\label{con:elements}
For $F\in\Th(T)$ let $p_F:E_F=\int_T F\longrightarrow T$ be its left fibration. Define $\eT(F)=\big((E_F)_{\minrep}\xrightarrow{p_F}T_{\minrep}\big).$ An object of $E_F$ is a pair $(t,x)$ with $x\in F(t)$. A morphism over $u:t\to t'$ consists of the identification $F(u)x\simeq x'$. In complete Segal notation the construction is $\Nn m E_F\simeq \int_{\sigma\in\Nn m T}F(\sigma(0)).$ The first face applies the transport of $F$; the remaining faces and degeneracies are induced by those of $T$. This formula includes all coherence rather than choosing lifts arrow by arrow.
\end{construction}

\begin{lemma}[Exactness of the total category]\label{lem:totallex}
If $F:T\to\Sps$ is left exact, then $E_F$ has finite limits and $p_F$ preserves them. Conversely, if a left fibration $p:E\to T$ has a finitely complete total category and preserves finite limits, its straightening $F$ is left exact.
\end{lemma}
\begin{proof}
\leavevmode
\begin{enumerate}
\item \textbf{Construct a limiting object in the total category.} Let $d:J\to E_F$ be a finite diagram. Its projection has a limit $t=\lim_J p_Fd$. The elements and their coherences in $d$ form a point of $\lim_JF(p_Fd)$. Left exactness identifies this with a point $x\in F(t)$, giving an object $(t,x)$ and a cone over $d$.
\begin{samepage}
For each vertex $j$ the lifted cone has the form
\[
\begin{tikzcd}[column sep=large,row sep=large]
 (t,x)\ar[r]\ar[d,"p_F"']&d(j)\ar[d,"p_F"]\\
 t\ar[r]&p_Fd(j).
\end{tikzcd}
\]
The compatibility of these squares is the point of
$\lim_JF(p_Fd)$ used in the construction.
\end{samepage}
\item \textbf{Verify its mapping-space universal property.} To check its universal property, let $(s,y)\in E_F$. The mapping space to $(t,x)$ is the space of maps $v:s\to t$ equipped with $F(v)y\simeq x$. The corresponding space of cones to $d$ consists of compatible maps $s\to p_Fd(j)$ and compatible identifications with its chosen elements. The limit properties of $t$ and $F(t)$ identify these spaces. This proves the limit property in $E_F$, including the empty diagram, and proves preservation by $p_F$.
\item \textbf{Recover left exactness from limits in the total category.} Conversely fix a finite limiting cone $t\to t_j$ in $T$ and a compatible family of points in $F(t_j)$. The family gives a diagram in $E$. Its limit projects to a limit of the $t_j$ and therefore determines a lift over $t$. Every arrow of a left fibration is cocartesian. The space of cones over this fixed base cone and this diagram identifies with the fiber of $F(t)\longrightarrow\lim_jF(t_j).$ Every such cone is limiting: comparison with a limiting lift gives a map whose projection identifies two limiting cones in $T$, hence is an equivalence. An arrow of a left fibration over an equivalence is an equivalence, so this is an equivalence of cones. The space of limiting lifts over the specified base cone is contractible. Thus every fiber of the displayed comparison is contractible. For the empty diagram, the terminal object of $E$, whose image is terminal in $T$, gives the same argument. Hence the comparison is an equivalence.\qedhere
\end{enumerate}
\end{proof}

\begin{theorem}[Faithful element realization]\label{thm:elementembedding}
Construction~\ref{con:elements} defines a fully faithful functor $\eT:\Th(T)\longrightarrow\TT/T_{\minrep}.$ It preserves bounded small limits. Its essential image consists precisely of theory morphisms $p:E_{\minrep}\to T_{\minrep}$ whose underlying functors are left fibrations. The inverse is the space-valued straightening of $p$.
\end{theorem}
\begin{proof}
\leavevmode
\begin{enumerate}
\item \textbf{Realize natural transformations as theory morphisms.} Lemma~\ref{lem:totallex} constructs the displayed theory morphism. A transformation $F\to G$ induces a functor $E_F\to E_G$ over $T$. It preserves finite limits, since the construction of a limiting object in Lemma~\ref{lem:totallex} is natural in $F$. It is therefore a morphism of minimal theories.
\begin{samepage}
For a transformation $\alpha:F\to G$, this construction gives
\[
\begin{tikzcd}[column sep=large,row sep=large]
 E_F\ar[r,"\int_T\alpha"]\ar[d,"p_F"']&E_G\ar[d,"p_G"]\\
 T\ar[r,"\id_T"']&T.
\end{tikzcd}
\]
Thus the transformation is realized over the fixed base, with limits
carried to the corresponding limits in the target total category.
\end{samepage}
\item \textbf{Prove full faithfulness and identify the essential image.} Conversely every functor over $T$ between these left fibrations automatically preserves their cocartesian arrows, since all their arrows are cocartesian. Straightening identifies its space of choices with the mapping space from $F$ to $G$. The same natural limit calculation shows that every such functor is left exact. Thus passing to theory morphisms discards no maps or homotopies, proving full faithfulness.
\item \textbf{Verify preservation of limits.} Left fibrations and their maps form a full subcategory of $\Cat/T$ closed under limits, with straightening computing those limits pointwise. Left-exact space-valued diagrams are also closed under limits, since limits commute with finite limits. Lemma~\ref{lem:minimal} identifies the resulting relative limits with those in $\TT/T_{\minrep}$. This proves limit preservation. Finally the converse part of Lemma~\ref{lem:totallex} proves the claimed essential-image characterization.\qedhere
\end{enumerate}
\end{proof}

\begin{construction}[Categorical realization]\label{con:J}
For $K\in\LT$ write $K_n(t)=\Fun([n],K(t))^{\core}.$ Each $K_n$ is left exact. Apply $\eT$ degreewise to obtain $(\J_TK)_n=(E_{K_n})_{\minrep}\longrightarrow T_{\minrep}.$ Limit preservation in Theorem~\ref{thm:elementembedding} carries the Segal and completeness conditions to the slice. Thus this is a complete category object in $\TT/T_{\minrep}$, equivalently a category object of $\TT$ over $\B_T$.

More explicitly, before applying the NU theory--model equivalence, its two simplicial directions are the spaces $\int_{\sigma\in\Nn m T}\Fun([n],K(\sigma(0)))^{\core}.$ This is an explicit formula for the entire realization. The $m$ direction gives the underlying theory category; the $n$ direction gives its internal categorical structure.
\end{construction}

\begin{theorem}[Fully faithful categorical NU realization]\label{thm:J}
Construction~\ref{con:J} defines an enriched fully faithful functor $\J_T:\LT\longrightarrow\CNU/\B_T.$ It preserves bounded small limits. Its essential image consists of those category objects over $T_{\minrep}$ whose simplicial levels have minimal representable-map structures and whose levelwise projections to $T$ are left fibrations. In particular,
\begin{equation}\label{eq:Jhom}
 \Nat_{\LT}(A,K)\simeq
 \Homc_{\CNU/\B_T}(\J_TA,\J_TK).
\end{equation}
\end{theorem}
\begin{proof}
\leavevmode
\begin{enumerate}
\item \textbf{Apply element realization levelwise.} Apply the fully faithful, limit-preserving functor $\eT$ to complete category objects. The equivalence $\operatorname{Cat}(\Th(T))\simeq\Lex(T,\Cat)$ is obtained by interchanging the $T$ and complete-Segal directions. This proves full faithfulness on mapping spaces and limit preservation. Degreewise straightening supplies the inverse on the specified image; its fibers recover $K_n(t)$, its transport recovers the $T$ direction, and the simplicial operators recover $K(t)$ and its composition.
\begin{samepage}
For every simplicial degree $n$, the construction commutes with evaluation:
\[
\begin{tikzcd}[column sep=large,row sep=large]
 \Lex(T,\Cat)\ar[r,"\J_T"]\ar[d,"K\mapsto K_n"']&
 \CNU/\B_T\ar[d,"\operatorname{ev}_n"]\\
 \Th(T)\ar[r,"\eT"']&\TT/T_{\minrep}.
\end{tikzcd}
\]
The right vertical functor evaluates the complete category object in
degree $n$.
\end{samepage}
\item \textbf{Recover full hom categories by cotensors.} For enrichment, cotensoring complete category objects by $[r]$ is a
limit construction. Limit preservation by $\eT$ therefore gives
$\J_T(K^{[r]})\simeq(\J_TK)^{[r]}$ in the slice, naturally in $[r]$.
Apply mapping-space full faithfulness to these cotensors and then
Lemma~\ref{lem:complete-nerve-comparison}. This proves
\eqref{eq:Jhom} on full hom $\infty$-categories. Degreewise straightening
gives the same essential-image description for category objects.\qedhere
\end{enumerate}
\end{proof}

This construction uses the underlying finite-limit category of $T$
as its indexing doctrine. It equips both the category of elements
and its base with the minimal representable-map structures.
Lemma~\ref{lem:minimal} then verifies the dependent-product condition
for their projection, giving the stated relative theory morphism.

\subsection{Free finite-limit completion}\label{subsec:free-completion}

An arbitrary native diagram may not arrive with a finite-limit doctrine.
We can construct one without changing the diagram on its original
arities: freely adjoin the finite limits needed for its interpretation.
The extension universal property applies to entire functor categories,
so it also retains natural transformations and modifications.

For a bounded small category $D$, let $\mathcal Q(D)$ be the full
subcategory of $\Fun(D,\Sps)$ generated by the covariant representables
$D(d,-)$ under finite colimits. Put
\[
 QD=\mathcal Q(D)^{\op},\qquad i_D:D\longrightarrow QD.
\]
The opposite Yoneda embedding $i_D$ is fully faithful, $QD$ is small
and finitely complete, and for every finitely complete $X$ restriction
gives
\begin{equation}\label{eq:free-finite-limit}
 \Lex(QD,X)\simeq\Fun(D,X).
\end{equation}
This is the free-completion universal property
\cite[Prop.~5.3.6.2]{Lurie2009}: extend a functor by finite limits on
the opposite finite-colimit expressions. The relations between these
expressions are those of their presheaf universal properties, so the
extension exists and is unique through a contractible space. Applying
the same construction to functor categories proves naturality on
transformations. For $X=\Cat$, applying it also to cotensors by $[n]$
makes \eqref{eq:free-finite-limit} an equivalence on full hom categories.
Thus every $C:D\to\Cat$ has a canonical left exact extension
$\overline C:QD\to\Cat$, with $\overline C i_D\simeq C$.
\subsection{Proof and naturality of universal relative classification}

The two universal properties can now be composed. The envelope classifies
the marked cone, and the categories-of-elements realization identifies
its maps with maps over the fixed NU base. We then check source and
doctrine changes so that the classification can be used with different
presentations of the same native operations.

\begin{proof}[Proof of Theorem~\ref{thm:main}]
\leavevmode
\begin{enumerate}
\item \textbf{Construct the classifier and its relative realization.} Construct $\Omega_{a,W}$ by \eqref{eq:envelope}. Proposition~\ref{prop:classifier} gives \eqref{eq:main1}. Apply Construction~\ref{con:J} to $\Omega_{a,W}$ and $K$; Theorem~\ref{thm:J} gives \eqref{eq:main2}.
\item \textbf{Describe the inverse correspondences.} In one direction a marked cone is sent to \eqref{eq:classifyformula}, then to its reflected classifying map, and finally to its degreewise element realization. In the other direction a relative NU map is straightened degreewise, giving a map $\Omega_{a,W}\to K$, and then evaluated on the universal cone. The straightening equivalence and the two adjunction identities show that these are inverse. Their constructions on $[r]$-families give the inverse identities on complete hom categories.
\item \textbf{Extend an arbitrary bounded diagram freely.} For arbitrary $C:D\to\Cat$, the free finite-limit universal property gives $\Lex(QD,\Cat)\simeq\Fun(D,\Cat)$ as full $\infty$-categories. It constructs $\overline C$ and the equivalence $\overline C i_D\simeq C$. Take $a=i_D$ in the first part. Notice that we keep the marked source $D^W$ in the envelope. At no point do we replace its cone category by all cones on $QD$.\qedhere
\end{enumerate}
\end{proof}

\begin{corollary}[A canonical relative NU classifier for every marked source]\label{cor:allmarked}
For every bounded marked $D^W$, define
\[
 \B_D=\disc((QD)_{\minrep}),\quad
 \Om_{D,W}=\J_{QD}\Env_{QD}(i_D,W),\quad
 \cK_C=\J_{QD}\overline C.
\]
Then $\laxlim_D^W C\simeq \Homc_{\CNU/\B_D}(\Om_{D,W},\cK_C)$ naturally in $C$, with all its ordinary natural transformations and modifications. For fixed $D$, the construction $C\mapsto\cK_C$ is enriched fully faithful.
\end{corollary}
\begin{proof}
\leavevmode
\begin{enumerate}
\item \textbf{Combine free completion with the relative embedding.} Compose the enriched equivalence given by free completion with the enriched fully faithful functor $\J_{QD}$. The representing equivalence is Theorem~\ref{thm:main}.\qedhere
\end{enumerate}
\end{proof}

\begin{proposition}[The relative NU realization of an arbitrary category]
\label{prop:relative-one-generator}
Let $D=\one$, let $T=QD$ be its free finite-limit completion, and
let $i:D\to T$ select its generator $x$. For every small
$\infty$-category $C$ let $\overline C:T\to\Cat$ denote its
left-exact extension. Then the marked envelope for the equivalence
marking is the constant terminal diagram and there is a natural
equivalence of full categories
\begin{equation}\label{eq:relative-one-generator}
 C\simeq\Nat_{\Lex(T,\Cat)}(\one_T,\overline C)
 \simeq\Homc_{\CNU/\B_T}(\B_T,\J_T\overline C).
\end{equation}
This realizes arbitrary morphisms of $C$, not merely its core; it
also applies to the empty category.
\end{proposition}
\begin{proof}
\leavevmode
\begin{enumerate}
\item \textbf{Compute the source envelope.}
The category $\mathcal Q(\one)$ defining $T$ is the finite-colimit
closure of the point in spaces. That point is terminal in
$\mathcal Q(\one)$, so its opposite $x$ is initial in $T$.
Thus $\Map_T(x,t)\simeq *$ for every $t$, and
$(i\downarrow t)$ is terminal. Its marked localization is terminal,
so $\Env_T(i,\Eq(\one))=\one_T$, already left exact.

\item \textbf{Recover the original category and its relative realization.}
Restriction along $i$ induces the enriched equivalence
$\Lex(T,\Cat)\simeq\Fun(\one,\Cat)=\Cat$ from
\eqref{eq:free-finite-limit}. Applied to the full hom from the
terminal diagram this gives
$\Nat(\one_T,\overline C)\simeq C$.
The complete nerve of $\one_T$ is levelwise the constant point-valued
functor on $T$, whose category of elements is $T$ itself. Hence
$\J_T\one_T\simeq\B_T$, and Theorem~\ref{thm:J} gives the second
equivalence. Since all equivalences hold on full hom categories,
noninvertible maps of $C$ are retained.
\qedhere
\end{enumerate}
\end{proof}

\begin{remark}[Finite limits versus finite products]
The free finite-limit doctrine $Q\one$ should not be identified with
the ordinary finite-product Lawvere theory $\mathrm{FinSet}^{\op}$.
Its defining finite colimits are homotopy finite colimits in spaces;
for example the homotopy coequalizer of two identities of the point
is $S^1$, also in $\mathcal Q(\one)$. The relative realization above
uses the genuine finite-limit completion.
\end{remark}

\begin{remark}[What is fixed]
The ambient theory category $\TT$ and the NU doctrine $\RNU$ are
the same for all sources. The marked source determines the relative
base $\B_D$ and the classifier $\Om_{D,W}$. Their maps to the base
retain the native arities in the NU realization.
\end{remark}

\begin{proposition}[Change of indexing doctrine]\label{prop:basechange}
Let $b:T\to S$ preserve finite limits. Pullback of models and pullback of relative NU objects satisfy a natural equivalence
\[
 \J_T(b^*K)\simeq
 \B_T\times_{\B_S}\J_SK,\qquad K\in\mathcal L_S.
\]
The map $\B_T\to\B_S$ is induced by $b_{\minrep}$. These equivalences are coherent under identities and composition.
\end{proposition}
\begin{proof}
\leavevmode
\begin{enumerate}
\item \textbf{Compute doctrine change by pullback.} For space-valued $F:S\to\Sps$, the left fibration of $Fb$ is the pullback $\int_T Fb\simeq T\times_S\int_SF.$ All these categories have minimal representable-map structures. By Lemma~\ref{lem:minimal} this is also the pullback of theories. Apply the formula degreewise to the complete Segal nerve of $K$. The comparisons are the natural base-change equivalences of unstraightening, which already include their higher coherence.
\begin{samepage}
The underlying base-change comparison is the pullback square
\[
\begin{tikzcd}[column sep=large,row sep=large]
 \displaystyle\int_T Fb\ar[r]\ar[d]&
 \displaystyle\int_SF\ar[d]\\
 T\ar[r,"b"']&S.
\end{tikzcd}
\qedhere
\]
\end{samepage}
\end{enumerate}
\end{proof}

\begin{proposition}[Transport of classifiers]\label{prop:classifierchange}
There is an enriched adjunction
\[
 b_!^{\mathrm{lex}}:\LT\rightleftarrows\mathcal L_S:b^*,
 \qquad b_!^{\mathrm{lex}}=L_{\mathrm{lex},S}\operatorname{Lan}_b.
\]
For every $a:D\to T$ and $W$ there is a canonical equivalence
\begin{equation}\label{eq:classbase}
 b_!^{\mathrm{lex}}\Omega_{a,W}\simeq\Omega_{ba,W}.
\end{equation}
If $v:D^W\to E^V$ is marked and $a_D\simeq a_Ev$ over $T$, it induces $c_v:\Omega_{a_D,W}\longrightarrow\Omega_{a_E,V},$ and restriction of cones is precomposition with $c_v$. In particular $c_v$ is an equivalence if and only if restriction is an equivalence for every $K\in\LT$.
\end{proposition}
\begin{proof}
\leavevmode
\begin{enumerate}
\item \textbf{Transport the representing object by adjunction.} Left Kan extension followed by the enriched reflector is left adjoint to restriction. Equivalently its enriched universal property follows from the ordinary adjunction applied also to every cotensor $K^{[r]}$. For $K\in\mathcal L_S$,
\begin{align*}
 \Nat(b_!^{\mathrm{lex}}\Omega_{a,W},K)
 &\simeq\Nat(\Omega_{a,W},b^*K)\\
 &\simeq\laxlim\nolimits_D^W(Kba)\\
 &\simeq\Nat(\Omega_{ba,W},K).
\end{align*}
Enriched Yoneda constructs \eqref{eq:classbase}; its definition by the same natural equivalences proves coherence for composable changes of doctrine.
\item \textbf{Identify source restriction and its equivalence criterion.} Functoriality of comma categories, localization, and reflection constructs $c_v$. Evaluation on the universal cone identifies precomposition with restriction. If restriction is an equivalence for every $K$, take cores and apply the ordinary Yoneda criterion to $c_v$; the converse is immediate.\qedhere
\end{enumerate}
\end{proof}

For a native arity map $a:D\to T$, let $\widehat a:QD\to T$ be its
finite-limit extension. Formula~\eqref{eq:classbase} gives
\[
 \widehat a_!^{\mathrm{lex}}\Env_{QD}(i_D,W)
 \simeq\Env_T(a,W).
\]
Thus change of doctrine carries the classifier from the universal
free completion to the native finite-limit doctrine while retaining
the source marking $W$.


\subsection{Cellular envelopes and finite shadows}\label{sec:thz:cellular}

The walking source uses the completed doctrine, whereas calculations
often begin with a smaller collection of named cells. Colimits of marked
sources pass through the envelope, and its universal property supplies a
test for whether a smaller source classifies the same internal objects.
This is the appropriate comparison between a finite presentation and its
completion.

Corollary~\ref{thz:cor:envelope-adjunction} identifies $\Env_T$ as the
left adjoint from bounded marked categories over $T$ to $\mathcal L_T$;
the right adjoint is cocartesian unstraightening equipped with its full
cocartesian marking.  We now apply that adjunction to marked colimits.

\begin{lemma}[Marked simplicial resolution]\label{thz:lem:marked-resolution}
Take a quasicategorical presentation $\phi:Q\xrightarrow{\simeq}T$ and
give $Q$ the pullback marking $\phi^*\tau$. Explicitly, from a
simplicial-space model of $T$, apply functorial fibrant replacement
in the complete Segal model structure and take its first row.
The derived counit gives $\phi$ by \cite[Thm.~4.11]{JoyalTierney2007}.
Applying the same functors to a map constructs its presentation.  For every simplex
$\sigma:\Delta^n\to Q$, give $\Delta^n$ the pullback marking
$\sigma^*\phi^*\tau$.  Then the canonical cocone induces an equivalence $T^\tau\simeq \operatorname*{colim}_{\sigma:\Delta^n\to Q} (\Delta^n)^{\sigma^*\phi^*\tau}$ in marked $\infty$-categories and in the slice over $T^\tau$.
\end{lemma}
\begin{proof}
\leavevmode
\begin{enumerate}
\item \textbf{Resolve the underlying category by simplices.} Regard $Q$ as a degreewise discrete simplicial space.
In degree $m$ the diagram of representables over $\Delta/Q$ has
homotopy colimit the space of factorizations $[m]\to[n]\to Q$
with changes of factorization. Over each $x\in Q_m$ its indexing
category has initial object $(x,\id_{[m]})$, and there are no maps
between the different fibers. Thus its colimit is the discrete
space $Q_m$. Apply complete-Segal reflection. Each representable
$\Delta[n]$ presents $[n]$, and the discrete simplicial space $Q$
presents its given quasicategory by
\cite[Thm.~4.11]{JoyalTierney2007}. This gives
$Q\simeq\operatorname*{colim}_{\sigma}[n]$ in $\Cat$.  Colimits of marked
$\infty$-categories have the same underlying colimit, and their marking
is the Berman marking generated by the images of the local markings.
\item \textbf{Recover the marking from marked edges.} Every locally marked edge maps to a $\phi^*\tau$-marked edge.  Conversely,
every $\phi^*\tau$-marked edge of $Q$ occurs as a marked $1$-simplex in
the simplex diagram.  Hence the generated marking is exactly
$\phi^*\tau$.  Transporting along the categorical equivalence $\phi$
gives the displayed equivalence over $T^\tau$.  Applying
Proposition~\ref{prop:marked-descent} to this marked colimit gives
the completed NU section formula.\qedhere
\end{enumerate}
\end{proof}

\begin{theorem}[Cellular construction of envelopes]\label{thz:thm:cellular}
For a marked colimit over $T$, $(D,W)\simeq\operatorname*{colim}_i(D_i,W_i),$ there is a canonical equivalence
\begin{equation}\label{thz:eq:env-colim}
 \Env_T(D,W)
 \simeq
 \operatorname*{colim}^{\mathcal L_T}_i
 \Env_T(D_i,W_i).
\end{equation}
In particular, Lemma~\ref{thz:lem:marked-resolution} constructs
$\Omega_{T,\tau}$ as the colimit of the envelopes of its marked simplices.
For the NU doctrine, the finite global named-cell diagram of
Construction~\ref{con:incidence} and its marked colimit construct the
envelope of the finite cellular shadow.
\end{theorem}
\begin{proof}
\leavevmode
\begin{enumerate}
\item \textbf{Use the envelope adjunction.} Corollary~\ref{thz:cor:envelope-adjunction} makes $\Env_T$ a left adjoint on the underlying $\infty$-categories, so it preserves the
displayed colimit.
\item \textbf{Verify the enriched colimit formula.} Equivalently, mapping the right side into $K\in\mathcal L_T$ gives the
limit of the marked functor categories
\[
 \Fun^{W_i\text{-}\coc}_{/T}
   (D_i,\textstyle\int K),
\]
which is the marked functor category from the colimit $(D,W)$.  Yoneda
gives \eqref{thz:eq:env-colim}.  Applying the same calculation to all
$K^{[n]}$ gives the full categorical statement.  The final assertions
apply this result first to Lemma~\ref{thz:lem:marked-resolution}.
For the finite NU source, give $A$ and $B$ the equivalence marking
and $H$ its structural marking over the doctrine. Their marked pushout
has underlying category $B\amalg_A H$ and the marking generated from
$H$. Hence its envelope is
$\Env(B)\amalg_{\Env(A)}\Env(H)$ in $\mathcal L_{\RNU}$, with the
three maps to the doctrine supplied by the realization. The incidence
colimits defining $A$ and $B$ can be evaluated by the same formula.\qedhere
\end{enumerate}
\end{proof}

\begin{theorem}[Finite-shadow criterion]\label{thz:thm:shadow}
For a marked functor $q:D^W\to T^\tau$, define
\[
 \Omega_q=\Env_T(q,W),
 \qquad
 \Omega_T=\Omega_{T,\tau}.
\]
There is a canonical map
$c_q:\Omega_q\to\Omega_T$.  Under the natural equivalences
\[
 \Nat(\Omega_T,K)\simeq\operatorname*{laxlim}^{\tau}K,
 \qquad
 \Nat(\Omega_q,K)\simeq
 \operatorname*{laxlim}^{W}(Kq),
\]
restriction along $q$ is precomposition with $c_q$.
Thus the comparison on internal objects is the horizontal map in
\[
\begin{tikzcd}[column sep=large,row sep=large]
 \Nat(\Omega_T,K)\ar[r,"c_q^*"]\ar[d,"\simeq"']&
 \Nat(\Omega_q,K)\ar[d,"\simeq"]\\
 \laxlim_T^\tau K\ar[r,"q^*"']&\laxlim_D^W(Kq).
\end{tikzcd}
\]
  The following are
equivalent:
\begin{enumerate}\renewcommand{\labelenumi}{(\roman{enumi})}
\item $c_q$ is an equivalence in $\mathcal L_T$;
\item for every left exact $K:T\to\Cat$, restriction is an equivalence $\operatorname*{laxlim}^{\tau}K \xrightarrow{\simeq} \operatorname*{laxlim}^{W}(Kq).$ \end{enumerate}
\end{theorem}
\begin{proof}
\leavevmode
\begin{enumerate}
\item \textbf{Identify restriction with precomposition.} The marked map $q:(D,W)\to(T,\tau)$ over $T$ induces $c_q$ by
functoriality of the envelope.  Naturality of
\eqref{thz:eq:envelope-UP} identifies restriction with precomposition by
$c_q$.  Hence (i) implies (ii).
\item \textbf{Detect equivalence by Yoneda.} Conversely, taking cores in (ii) says
that precomposition with $c_q$ induces an equivalence of mapping spaces
into every object of $\mathcal L_T$.  The $\infty$-categorical Yoneda
criterion makes $c_q$ an equivalence.\qedhere
\end{enumerate}
\end{proof}

\paragraph{The finite NU shadow has sections that do not extend.}
For the map $q:P_\infty^\tau\to\RNU^{\tau_{\NU}}$,
$c_q$ is not an equivalence. We exhibit a left-exact $K$ and a finite
section outside the essential image of restriction.

Put $p=\partial_R$. The selected terminal square gives
$\mathsf1_R^*p\simeq\id_1$. Base change for the polynomial along
$\mathsf1_R:1\to R$ gives
$1\times_R P_p X\simeq X$, naturally in $X$: the base-changed
dependent product is along $\id_1$, with constant parameter $X$.
Let $\iota_X:X\to P_pX$ be the resulting map and $j=\iota_U$.
Naturality gives a cartesian square from $\partial$ to $P_p\partial$,
with base map $j$: in both degrees it is pullback to the same selected
terminal domain. The mixed product square now gives
$(\Pi j)^*\partial\simeq\partial$ over $U$, hence
$\Pi j\simeq\id_U$ by univalence. For a completed section $x$ the
composition law consequently makes
\[
 \chi_\Pi(x)\circ K(\Pi)(\chi_j(x))\simeq \id_{x_U}
\]
after the displayed base identification. In particular $\chi_\Pi(x)$
is a split epimorphism.

Let $A=\mathrm{FinSet}_{\minrep}$ and let $F:\RNU\to\Sps$ be its model.
The cartesian product functor $A\times A\to A$ preserves finite limits
and is a theory morphism by minimality. Its coherent symmetric products
give $F$ a commutative monoid structure. In simplicial degree $n$ put
$K_n=F^n$, with faces and degeneracies of the reduced bar construction.
These are left-exact functors and form a Segal object. For completeness
write $F(t)\simeq\Map_{\TT}(A_t,A)$, where $A_t$ corresponds to the
representable $h_t$. The product on this space is pointwise product of
functors. If a theory functor has an inverse under it, each of its
values is a singleton. The full subgroupoid of singleton-valued theory
functors is contractible, since the full subcategory of singletons is
equivalent to $*$. Conversely all such functors are units. Thus
$K_0(t)\to K_1(t)^{\mathrm{eq}}$ is an equivalence, giving
$K\in\Lex(\RNU,\Cat)$. At $U$, the sole object's endomorphism space is
the groupoid of finite sets under product. At $E$ it is the groupoid
of pointed finite sets, with $K(\partial)$ forgetting the point:
these are the objects over the terminal context in the classifier
and its generic term \cite[Def.~5.2, Thm.~5.29]{NguyenUemura2025}.

Restrict the unique-object strong section of this bar diagram to the
shadow. On the graph replace the unit labels of $\lambda,\Pi$ by a
pointed two-element set and its underlying set, respectively. In the
mixed square both composites are that underlying set with their
canonical unit identifications, so its identity supplies the square
filler. These edges occur only in that square and its subfaces among
the named occurrences (the list in Proposition~\ref{prop:charge}).
Use the same restrictions on those subfaces and the original strong
fillers on every other occurrence. The pullback
\eqref{eq:normalized-cell-mapping} assembles them to a finite marked
section. Its comparison at $\Pi$ has no right inverse: such an inverse
would be a finite set $S$ with $2|S|=1$. Equivalence of sections
conjugates the comparisons by equivalences and preserves splitting.
This proves the claimed failure even up to equivalence. The criterion
of Theorem~\ref{thz:thm:shadow} gives the assertion about $c_q$.

\begin{theorem}[Pointwise marked-coinitiality test]
\label{thm:pointwise-marked-coinitial}
For a marked functor $q:(D,W)\to(T,\tau)$ let
\[
 \vartheta_t:(q\downarrow t)[W_t^{-1}]
 \longrightarrow (T_{/t})[\tau_t^{-1}]
\]
be the map induced by $q$. Restriction
$\laxlim_T^\tau K\to\laxlim_D^W(Kq)$ is an equivalence of full
$\infty$-categories for every functor $K:T\to\Cat$ if and only if
$\vartheta_t$ is an equivalence for every $t\in T$.
\end{theorem}
\begin{proof}
\leavevmode
\begin{enumerate}
\item \textbf{Identify restriction with the comma comparison.}
The marked comma universal property, before left-exact reflection,
identifies the restriction functor with precomposition by the natural
map $R_{q,W}\to R_{\id_T,\tau}$. If this is an equivalence then every
restriction is an equivalence. 
\item \textbf{Detect the equivalence by Yoneda.}
Conversely, the latter property for all
$K\in\Fun(T,\Cat)$ implies, on cores, that precomposition induces
equivalences of mapping spaces into every such $K$. Yoneda then makes
the map of comma diagrams an equivalence. Equivalences of functor
diagrams are detected pointwise, and their components are exactly
$\vartheta_t$. The original enriched comma equivalence retains all
noninvertible modifications.
\qedhere
\end{enumerate}
\end{proof}

For the NU map $q:P_\infty^\tau\to\RNU^{\tau_{\NU}}$, the theorem
tests restriction to the initial finite shadow. To compute the
completed source directly, we next continue that shadow by finite
relative simplex stages. Theorem~\ref{thm:relative-cellular-completion}
constructs their colimit and computes the extension category at
each stage.


\subsection{Relative cellular completion of a marked source}
\label{subsec:relative-cellular-completion}

Starting with the finite named-cell shadow, we adjoin finite collections
of coherent simplices of the doctrine. The same construction works for
any marked source. Its envelope colimit is the completed walking
classifier, and the restriction maps retain each extension category.

\begin{construction}[Finite relative simplex stages]
\label{con:relative-simplex-stages}
Let $q:(D,W)\to(T,\tau)$ be bounded and marked, and choose a
quasicategory $Q\xrightarrow{\simeq}T$ carrying the pulled-back
marking, as in Lemma~\ref{thz:lem:marked-resolution}. Put
$\Sigma_Q=\Delta/Q$. Let $\mathcal F_Q$ be the filtered poset of
finite full subcategories $J\subseteq\Sigma_Q$ (including the empty
one). For $J\in\mathcal F_Q$ form, in marked categories over $T$,
\[
 (P_J,V_J):=\operatorname*{colim}_{\sigma\in J}
        (\Delta^{n_\sigma},W_\sigma),\qquad
 (D_J,W_J):=(D,W)\amalg_{(D,W)\times_T(P_J,V_J)}(P_J,V_J).
\]
Write $U_J$ for the pullback marking on $D\times_TP_J$.
Thus all markings used here are obtained by marked pullbacks,
colimits, and pushouts, not specified as additional fields.
Write $q_J:(D_J,W_J)\to(T,\tau)$ for its canonical map. The stage
$J=\varnothing$ is $(D,W)$, and inclusions $J\subseteq J'$ give
coherent marked maps $D_J\to D_{J'}$ over $T$.
\end{construction}

\begin{lemma}[Filtered exactness for the relative stages]
\label{lem:filtered-marked-exactness}
Filtered colimits of bounded marked $\infty$-categories commute with
finite homotopy limits, at the corresponding universe level. In
particular, for the construction above,
\[
 \operatorname*{colim}_{J\in\mathcal F_Q}
       (D\times_TP_J)
 \simeq D\times_T
       \operatorname*{colim}_{J\in\mathcal F_Q}P_J.
\]
\end{lemma}
\begin{proof}
\leavevmode
\begin{enumerate}
\item \textbf{Compute the filtered colimit levelwise.}
Present $\Cat$ by complete Segal spaces. Filtered colimits of
simplicial spaces are levelwise, and filtered colimits commute with
finite limits of spaces by~\cite[Thm.~9.1.5.1]{Kerodon}.
Consequently the Segal comparisons of a filtered colimit are
again equivalences. In a Segal space an edge $u:x\to y$ is
invertible precisely when there are $v,w:y\to x$ and triangles
with long edges $\id_x,\id_y$ witnessing $vu\simeq\id_x$ and
$uw\simeq\id_y$. This holds because it is the two-sided inverse
test in the homotopy category. The space of these witnesses over
the edge space is formed by finite homotopy limits of degrees
$0,1,2$, using the identity degeneracy for the long edges.
It commutes with filtered colimits. The equivalence locus is the
union of components hit by its projection. Such images commute
with filtered colimits of spaces: a point and a path witnessing
membership are represented at a common stage, and the locus
over each component is either empty or the entire component.
Thus the equivalence loci also commute with this colimit.
Applying this to the completeness maps from degree zero shows
that the colimit is complete.
The resulting levelwise colimit is a complete Segal space and
computes the filtered colimit in $\Cat$. Its finite-limit
comparisons are levelwise the corresponding comparisons in spaces,
and hence are equivalences.


\item \textbf{Retain the marked arrows.}
In marked categories, retain at each simplicial level the replete
subspace of marked edges. Equivalences and finite composites are
detected at a filtered stage; the generated marked-edge family of
the filtered colimit is therefore the filtered colimit of the marked
families. Applying the same finite-limit argument to these families
proves the marked assertion. The displayed pullback comparison is its
specialization to the fixed object $(D,W)$ and the diagrams $P_J$.
\qedhere
\end{enumerate}
\end{proof}

\begin{theorem}[Explicit relative cellular completion]
\label{thm:relative-cellular-completion}
For Construction~\ref{con:relative-simplex-stages}, the canonical
cocones induce equivalences in marked categories over $T$ and in
left-exact categorical models:
\begin{align}
 \operatorname*{colim}_{J\in\mathcal F_Q}(D_J,W_J)
 &\simeq(T,\tau),\label{eq:relative-stage-colim}\\
 \operatorname*{colim}_{J\in\mathcal F_Q}^{\mathcal L_T}
       \Env_T(q_J,W_J)
 &\simeq\Omega_{T,\tau}.\label{eq:relative-env-colim}
\end{align}
For every $K\in\Lex(T,\Cat)$ this gives an equivalence of full
$\infty$-categories
\begin{equation}\label{eq:relative-stage-internal-lim}
 \operatorname*{laxlim}_T^\tau K
 \simeq
 \lim_{J\in\mathcal F_Q^{\op}}
 \Fun^{W_J\text{-}\coc}_{/T}(D_J,\textstyle\int_TK).
\end{equation}
The restrictions and their extension categories are explicit: for a
marked functor $s_J:D_J\to\int_TK$ the category of extensions to
$J'\supseteq J$ is the homotopy fiber
\begin{equation}\label{eq:relative-stage-extension}
 \Fun^{W_{J'}\text{-}\coc}_{/T}
 (D_{J'},\textstyle\int_TK)
 \times_{\Fun^{W_J\text{-}\coc}_{/T}(D_J,\int_TK)}\{s_J\}.
\end{equation}
These constructions are natural in the marked source and in ordinary
maps of left-exact categorical macrocosms.
\end{theorem}
\begin{proof}
\leavevmode
\begin{enumerate}
\item \textbf{Recover the completed marked source.}
Let $\pi:\int_{J\in\mathcal F_Q} J\to\Sigma_Q$ be evaluation.
For $\sigma\in\Sigma_Q$ an object of $\sigma/\pi$ is a triple
$(J,\eta,\alpha:\sigma\to\eta)$ with $\eta\in J$. Enlarge $J$
to the finite full subcategory $J^+$ on its objects and $\sigma$.
This gives an endofunctor $A$ on $\sigma/\pi$ and a natural map
$\id\to A$. The functor
$B(J,\eta,\alpha)=(J^+,\sigma,\id)$ has a natural map to $A$
given by $\alpha$ in $J^+$. The constant functor at
$(\{\sigma\}_{\rm full},\sigma,\id)$ maps naturally to $B$.
This zigzag contracts the nerve of $\sigma/\pi$, so $\pi$ is final.
Iterated colimits over the Grothendieck construction therefore
recover the colimit over $\Sigma_Q$.

Thus the marked simplicial resolution,
Lemma~\ref{thz:lem:marked-resolution}, gives
\[
 \operatorname*{colim}_{J\in\mathcal F_Q}P_J
 \simeq (T,\tau).
\]
By Lemma~\ref{lem:filtered-marked-exactness},
$\operatorname*{colim}_J(D\times_TP_J)\simeq D\times_TT\simeq D$.
As a colimit functor preserves pushouts, the definition of $D_J$
then yields
\begin{align*}
 \operatorname*{colim}_JD_J
 &\simeq
 \left(\operatorname*{colim}_JD\right)
 \amalg_{\operatorname*{colim}_J(D\times_TP_J)}
        \operatorname*{colim}_JP_J\\
 &\simeq D\amalg_DT\simeq T.
\end{align*}
All these equivalences lie in the marked slice over $T$, proving
\eqref{eq:relative-stage-colim}.
\item \textbf{Apply the enriched envelope.}
The envelope is the left adjoint of
Corollary~\ref{thz:cor:envelope-adjunction}; hence it preserves this
marked colimit. Its value on the terminal marked source
$(\id_T,\tau)$ is $\Omega_{T,\tau}$. This proves
\eqref{eq:relative-env-colim}. Mapping the latter equivalence into
$K$, using the full hom-category universal property of the envelope
and the naturality of its cotensor tests, proves
\eqref{eq:relative-stage-internal-lim}.
\item \textbf{Construct the extension categories.}
For $J\subseteq J'$, restriction gives a functor between the
respective marked-section categories. By definition its homotopy
fiber at $s_J$ is \eqref{eq:relative-stage-extension}. The
homotopy-pushout universal property computes each stage category as
\[
 \Fun^{W_J\text{-}\coc}_{/T}(D_J,\textstyle\int K)
 \simeq
 \Fun^{W\text{-}\coc}_{/T}(D,\textstyle\int K)
 \times_{\Fun^{U_J\text{-}\coc}_{/T}(D\times_TP_J,\int K)}
 \Fun^{V_J\text{-}\coc}_{/T}(P_J,\textstyle\int K).
\]
The markings $U_J$ and $V_J$ are exactly those constructed in
Construction~\ref{con:relative-simplex-stages}. Thus a finite selection of target simplices
gives an explicit matching category, with all its morphisms.
The overlap $D\times_TP_J$ can be large, and a transition can identify
old objects: for $T=*$, $D$ two discrete points and $J$ the target
vertex, the pushout is $D\amalg_D *=*$. Proposition~\ref{prop:relative-simplex-attachments}
also presents completion by actual simplicial boundary attachments,
with a finite boundary at each elementary step.\qedhere
\end{enumerate}
\end{proof}

\begin{corollary}[Completed NU classifier from the named-cell shadow]
\label{cor:NU-relative-completion}
Apply Construction~\ref{con:relative-simplex-stages} to
$q:P_\infty^\tau\to\RNU^{\tau_{\NU}}$. There is a filtered
system of explicitly constructed marked sources beginning at
$P_\infty^\tau$ whose colimit is $\RNU^{\tau_{\NU}}$; their
left-exact envelopes have colimit $\Omega_{\NU}$. Consequently the
completed internal NU category of any macrocosm $K$ is the limit of
its marked-section categories over these stages, including all
extension categories and noninvertible modifications.
\end{corollary}
\begin{proof}
\leavevmode
\begin{enumerate}
\item \textbf{Specialize the relative cellular construction.}
This is Theorem~\ref{thm:relative-cellular-completion} with
$T=\RNU$ and $\tau=\tau_{\NU}$. The initial-stage assertion follows
from $D_\varnothing=P_\infty^\tau$.
\qedhere
\end{enumerate}
\end{proof}

\begin{proposition}[A relative presentation by simplex attachments]
\label{prop:relative-simplex-attachments}
For the same marked source $q:(D,W)\to(T,\tau)$ there is a filtered
presentation starting at $(D,W)$ in which each transition is presented
by finitely many attachments of simplices along their boundaries.
Its marked colimit is $(T,\tau)$ under $(D,W)$. Its envelope colimit,
its category of compatible sections, and its category of completed
extensions of a fixed section of $D$ agree with those of
Theorem~\ref{thm:relative-cellular-completion}.
\end{proposition}
\begin{proof}
\leavevmode
\begin{enumerate}
\item \textbf{Embed the source in a marked presentation of the target.}
Present $q$ by a map $s:S\to Q$ of quasicategories with the
corresponding edge markings. Let $J$ be the nerve of the groupoid
with two objects and a unique isomorphism between any ordered pair.
In simplicial sets form
\[
 M=(S\times J)\amalg_{S\times\{1\}}Q,\qquad
 \pi:M\longrightarrow Q,\qquad S_0=S\times\{0\}\subset M .
\]
The map $\pi$ is $s\operatorname{pr}_S$ on the cylinder and the
identity on $Q$. Both $S_0$ and $Q$ embed as disjoint simplicial
subsets. The endpoint inclusion $S\times\{1\}\to S\times J$ is a
monomorphism and a categorical equivalence, since $J\simeq *$.
The displayed pushout therefore presents a categorical pushout
\cite[Ex.~4.5.5.12]{Kerodon}. It follows that $Q\to M$ and
$\pi$ are inverse equivalences after categorical realization.
Give $M$ the marking generated by $W$ on $S_0$ and $\tau$ on $Q$.
Both $Q\to M$ and $\pi$ preserve these markings, since $q$ is
marked. Their inverse equivalences respect the markings by
invariance under equivalence of arrows. In particular the marked
realization of $\pi$ is an equivalence over $T$, whose restriction
to $S_0$ is $q$.

\item \textbf{Filter by finite relative cells.}
For each finite simplicial subset $A\subseteq M$, including the
empty subset, set $M_A=S_0\cup A$ and mark the edges generated by
$W$ and the edges of $A\cap Q$ which belong to $\tau$. All
markings here mean the generated markings after categorical
realization. The subsets form a filtered poset under union and
$M_\varnothing=S_0$. For $A\subseteq A'$ attach the finitely many
new nondegenerate simplices of $M_{A'}$ in order of dimension.
Each step is a simplicial pushout along
$\partial\Delta^n\hookrightarrow\Delta^n$, with boundary map
given by the actual faces in the preceding stage. The uniqueness of
the expression of a simplex as a degeneracy of a nondegenerate one
gives this pushout even when the boundary map identifies faces.
On the characteristic simplex mark the edges whose images are
prescribed edges of $S_0$ or of $A'\cap Q$, and the degenerate
edges. Give the boundary their restrictions. Every prescribed
boundary edge already has its mark at the preceding step, since
that step contains the edge itself and $S_0$ and $Q$ are disjoint.
The categorical pushout criterion just used,
followed by generation of markings, gives the corresponding marked
pushout over $T$.

For the colimit, regard $M_A$ as a degreewise discrete simplicial
space. Its filtered colimit is $M$ degreewise: for any simplex of
$M$, the poset of stages containing it is nonempty and filtered.
The complete-Segal reflector preserves this colimit and gives the
category presented by each simplicial set
\cite[Thm.~4.11]{JoyalTierney2007}. Every final marked generator
already belongs to $S_0$ or occurs in a finite stage of $Q$.
Thus the marked colimit is the realization of $M$, identified
by $\pi$ with $(T,\tau)$ under the prescribed source.

\item \textbf{Compute the extensions for a single cell.}
Write $\mathcal S_K(X)$ for the full category of lifts of a presented
marked $X\to T$ to $\int_TK$ carrying its marked edges to cocartesian
arrows. For a step $Y'=Y\amalg_{\partial\Delta^n}\Delta^n$ and a
section $b$ on $Y$, marked descent and pullback pasting give
\[
 \begin{aligned}
 \mathcal S_K(Y')&\simeq
 \mathcal S_K(Y)\times_{\mathcal S_K(\partial\Delta^n)}
                         \mathcal S_K(\Delta^n),\\
 \mathcal S_K(Y')\times_{\mathcal S_K(Y)}\{b\}
 &\simeq\mathcal S_K(\Delta^n)
       \times_{\mathcal S_K(\partial\Delta^n)}\{b|_{\partial\Delta^n}\}.
 \end{aligned}
\]
Here all shape maps to $T$ are the ones specified by $\pi$ and the
characteristic simplex. This includes $n=0$ with empty boundary.
Iteration gives the extension category at every finite transition.
Applying the envelope left adjoint yields the claimed envelope
colimit. Marked descent yields the section limit over the opposite
poset. Both this limit and \eqref{eq:relative-stage-internal-lim}
are identified with the same restriction problem over
$\mathcal S_K(D)$, since both marked colimit identifications
restrict to $q$. Taking the fiber at any section of $D$ proves
the last assertion.\qedhere
\end{enumerate}
\end{proof}

\section{Iterated theorization and higher operations}
\label{sec:higher-composition}

An operation category suggests the next step by the same reasoning
that first produced its kernels. A composition functor represents the
operations from an input string to a variable output. Keeping that
functor of outputs allows composition itself to be theorized.

We begin with the balanced composites already determined by a category
over $T$. We then construct independent multikernels on categories of
lower operations, their substitution product, and their coherent tree
composites. Corepresentability recovers composition functors; the image
of a universal operation determines each comparison map. This realizes
the composition-bimodule calculation underlying theorization
\cite[\S2.4]{Matsuoka2017}. A grading over the original doctrine
retains the base simplices throughout the construction. We also give the
complementary categorical iteration: complete category objects of
marked-cocartesian operation theories, which retain unrepresented
kernels in each external degree.

Write $\operatorname{Cat}^0(\mathcal A)=\mathcal A$ and
$\operatorname{Cat}^{r+1}(\mathcal A)
 =\operatorname{Cat}(\operatorname{Cat}^r(\mathcal A))$.
An $r$-fold complete category object is regarded as a diagram on
$(\Delta^{\op})^r$ satisfying the Segal and completeness comparisons in
each variable.  We write $\mathbf m=(m_1,\ldots,m_r)$ and use the
complete Segal nerve convention $N_qC=\Fun([q],C)^\core$.
We write $\delta^r A$ for the constant $r$-fold complete category
object on $A$ (and $\delta^0 A=A$).  In particular, a degreewise
construction on an $r$-fold object always acts in its external
multisimplicial degrees; it does not denote a cotensor in a slice.
All constructions are at a finite categorical stage and at the
relative universe levels fixed in Section~\ref{subsec:conventions}.

\subsection{Derived higher composition from normalized kernels}
\label{subsec:derived-higher-composition}

For $p:E\to T$ write $C_s=E_s$ and let $P_f$ be the normalized
kernel of~\eqref{thz:eq:kernel}. In particular, the construction
below does not require cocartesian representatives of the operations.

\begin{construction}[Bar composite of operation kernels]
\label{con:derived-bar-composite}
For composable $f_i:s_{i-1}\to s_i$ ($1\leq i\leq n$), define the
profunctor from $C_{s_0}$ to $C_{s_n}$ by
\begin{equation}\label{eq:derived-bar-composite}
 \mathsf B_{(f_1,\ldots,f_n)}(x,z)
 :=\int^{x_1\in C_{s_1}}\cdots
    \int^{x_{n-1}\in C_{s_{n-1}}}
    P_{f_1}(x,x_1)\times P_{f_2}(x_1,x_2)
    \times\cdots\times P_{f_n}(x_{n-1},z).
\end{equation}
For $n=1$ this is $P_{f_1}$; for $n=0$ it is the identity
profunctor $\Map_{C_{s_0}}$. The coends are homotopy coends in
spaces. For example, the two-factor composite
$P_g\odot P_f$ is computed by the geometric realization of the
simplicial two-sided bar object
\begin{equation}\label{eq:derived-bar-resolution}
 \operatorname{Bar}_k(P_f,C_t,P_g)(x,z)
 =\int_{(y_0,\ldots,y_k)\in(C_t^\core)^{k+1}}
 P_f(x,y_0)\times
 \prod_{i=1}^k\Map_{C_t}(y_{i-1},y_i)
 \times P_g(y_k,z).
\end{equation}
Its first and last faces use the respective profunctor actions,
its internal faces compose adjacent arrows, and its degeneracies
insert identities. Iterating this construction gives a multi-bar
presentation of~\eqref{eq:derived-bar-composite}.
\end{construction}

\begin{theorem}[Coherent derived composition and represented recovery]
\label{thm:derived-higher-composition}
For every bounded $p:E\to T$, composition in $E$ induces canonical
natural transformations of profunctors
\begin{equation}\label{eq:derived-bar-comparison}
 \mu_{f_n,\ldots,f_1}:
 \mathsf B_{(f_1,\ldots,f_n)}
 \longrightarrow P_{f_n\cdots f_1},\qquad n\geq0.
\end{equation}
They are functorial in functors over $T$, compatible with every
regrouping of the $f_i$, and retain all higher associative and unit
coherences. If $p$ has cocartesian lifts in each of the degrees
$f_1,\ldots,f_n$, then $\mu_{f_n,\ldots,f_1}$ is an equivalence.
In particular, for $p\in\Thz_\tau(T)$ and $f_i\in\tau$ all these
marked higher composition comparisons are equivalences, while the
unmarked composites remain the explicitly constructed bimodules
\eqref{eq:derived-bar-composite} and maps
\eqref{eq:derived-bar-comparison}.
\end{theorem}
\begin{proof}
\leavevmode
\begin{enumerate}
\item \textbf{Construct the relative tensor products.}
The categories of space-valued profunctors are presentable, and
cartesian product of spaces preserves colimits in each variable.
Consequently their relative tensor products exist and are computed by
the two-sided bar construction; this is the profunctor instance of
the relative tensor product and its bar formula in
\cite[\S4.4.2--4.4.3]{LurieHA}. The object
\eqref{eq:derived-bar-resolution} is obtained by evaluating this
bar construction on $x$ and $z$. Fubini for homotopy coends gives
canonical equivalences between all parenthesizations of the
multi-bar object; these equivalences are coherent because they are
induced by the same iterated colimit.
\item \textbf{Construct comparison with composition in $E$.}
For every intermediate string $(x_1,\ldots,x_{n-1})$,
composition of arrows in $E$ supplies a map from the product
of operation spaces in~\eqref{eq:derived-bar-composite}
to $P_{f_n\cdots f_1}(x,z)$. These maps respect the actions of
the $C_{s_i}$: composing with a vertical arrow at an intermediate
object on either adjacent side gives the same total composite.
They therefore descend along the multi-bar realization to
\eqref{eq:derived-bar-comparison}. The complete Segal diagram
$N(E)\to N(T)$ carries this construction on every string
simultaneously, so the comparisons commute with faces,
degeneracies, higher base simplices, and all iterated regroupings.
A functor $E\to E'$ over $T$ acts on this same string diagram,
proving functoriality.
\item \textbf{Compute the represented degrees.}
Let $u_i$ denote cocartesian transport over $f_i$. Cocartesian
factorization gives
$P_{f_i}(a,b)\simeq\Map_{C_{s_i}}((f_i)_!a,b)$.
Starting with the first factor of the iterated coend, co-Yoneda
eliminates the integration variable $x_1$; repeating for
$x_2,\ldots,x_{n-1}$ identifies the multi-bar composite with
\[
 \Map_{C_{s_n}}
  ((f_n)_!\cdots(f_1)_!x,z).
\]
Cocartesian arrows compose, and the composite cocartesian arrow
represents $P_{f_n\cdots f_1}(x,-)$ by the same formula.
Under these identifications $\mu_{f_n,\ldots,f_1}$ is induced by
the actual composite cocartesian lift and is therefore an
equivalence. Its compatibility with every regrouping has already
been established in Step~2.\qedhere
\end{enumerate}
\end{proof}

\subsection{Independent composition operations}
\label{subsec:independent-multikernels}

The bar construction starts with a string of operations and computes
its balanced composite. To theorize composition itself, we keep one
more kind of datum: the operations from a string to a proposed output.
Their substitution will again be balanced, but their spaces are now
part of the structure.

Fix a small set $X$ of lower objects and small categories
$\mathcal A(x,y)$ for $x,y\in X$. An object of $\mathcal A(x,y)$
is a potential operation from $x$ to $y$; a morphism compares two such
operations. This is a local construction at the specified lower
objects. All products, coends, and collections below lie in the
universe fixed in Section~\ref{subsec:conventions}.

\begin{definition}[Multicollections]\label{def:multicollections}
A multicollection $M$ on $\mathcal A$ consists, for every $n\geq0$
and string $\mathbf x=(x_0,\ldots,x_n)$, of a functor
\[
 M_{\mathbf x}:
 \left(\prod_{i=1}^n\mathcal A(x_{i-1},x_i)\right)^{\op}
 \times\mathcal A(x_0,x_n)\longrightarrow\Sps.
\]
We write its values as $M(a_1,\ldots,a_n;b)$, leaving the lower
objects implicit. For $n=0$ this is a functor
$\mathcal A(x_0,x_0)\to\Sps$. A map of multicollections is a family
of natural transformations. Their category is denoted
$\operatorname{Coll}(\mathcal A)$.
\end{definition}

Contravariance in the inputs lets us substitute a comparison into
an input. Covariance in the output lets us follow an operation by
a comparison. These two actions are exactly what is needed to
identify the two appearances of an intermediate operation.

\begin{construction}[Balanced substitution]\label{con:balanced-substitution}
For $M,N\in\operatorname{Coll}(\mathcal A)$ put
\begin{equation}\label{eq:multikernel-substitution}
\begin{split}
 (M\circ N)(a_1,\ldots,a_n;b)
 ={}&\coprod_{\substack{k\geq0,\ n_1+\cdots+n_k=n\\n_i\geq0}}
 \int^{c_1,\ldots,c_k}
 M(c_1,\ldots,c_k;b)\\[-2pt]
 &\hspace{1.8em}\times
 \prod_{j=1}^k
 N(a_{r_{j-1}+1},\ldots,a_{r_j};c_j),
 \qquad r_j=\sum_{i\leq j}n_i.
\end{split}
\end{equation}
Here $r_0=0$ and
$c_j\in\mathcal A(x_{r_{j-1}},x_{r_j})$.
An empty block contributes a nullary operation. The summand $k=0$
occurs only when $n=0$. Each integral is a homotopy coend, computed
by the two-sided bar construction.

Define the unit multicollection $I_{\mathcal A}$ by
\[
 I_{\mathcal A}(a;b)=\Map_{\mathcal A(x,y)}(a,b)
\]
in arity one and by the empty space in every other arity.
\end{construction}

\begin{proposition}[The substitution monoidal structure]
\label{prop:substitution-monoidal}
Formula~\eqref{eq:multikernel-substitution} and
$I_{\mathcal A}$ define a monoidal structure on
$\operatorname{Coll}(\mathcal A)$. A unital associative algebra
for this structure supplies coherent substitution along every finite
planar rooted tree, including trees with nullary vertices.
\end{proposition}
\begin{proof}
\leavevmode
\begin{enumerate}
\item \textbf{Identify the common three-level formula.}
Expand $(L\circ M)\circ N$ using
\eqref{eq:multikernel-substitution}. Its summands are three-level
planar trees, with $L$ at the root, $M$ at the middle vertices,
and $N$ at the bottom vertices. Every internal edge contributes
one coend variable. Expanding $L\circ(M\circ N)$ gives the same
trees and the same integrands. Fubini for homotopy coends and
distribution of finite products of spaces over colimits identify
the two expressions.
\item \textbf{Remove the unary identity vertices.}
In $I_{\mathcal A}\circ M$ only the unary root remains, and
co-Yoneda identifies its coend with $M$. In $M\circ I_{\mathcal A}$
only singleton input blocks remain. Applying co-Yoneda in each
input variable again gives $M$. These identifications also cover
the nullary part.
\item \textbf{Retain the entire coherence diagram.}
For any number of levels, use the single diagram indexed by the
corresponding levelled trees and their internal-edge variables.
Every reassociation is the canonical comparison between iterated
colimits of this diagram. The universal property of the colimit
supplies their higher compatibility, including the unit comparisons
from Step~2. This constructs the monoidal structure. An associative
algebra in it is, by definition, an algebra over the associative
$\infty$-operad; its multiplication and unit therefore evaluate
all levelled trees compatibly. Inserting unary identities allows
any finite planar rooted tree to be levelled, and Step~2 identifies
the results for different insertions.\qedhere
\end{enumerate}
\end{proof}

\begin{definition}[Normalized independent multikernels]
\label{def:independent-multikernels}
A system of independent composition multikernels on $\mathcal A$
is a unital associative algebra $M$ in
$(\operatorname{Coll}(\mathcal A),\circ,I_{\mathcal A})$
whose unary unit map
\[
 \Map_{\mathcal A(x,y)}(a,b)\longrightarrow M(a;b)
\]
is an equivalence. We use this equivalence to normalize its unary
part. The components of its multiplication are the grafting maps.
\end{definition}

Thus the unary comparisons are the morphisms already present in
$\mathcal A$. The higher-arity spaces are independent data, with
substitution supplied by the algebra structure. For example,
three consecutive lower objects determine the binary rows, and
a four-object string has the grafting map
\begin{equation}\label{eq:binary-multikernel-grafting}
 \int^{c\in\mathcal A(x_0,x_2)}
 M(a_1,a_2;c)\times M(c,a_3;b)
 \longrightarrow M(a_1,a_2,a_3;b).
\end{equation}
The alternative binary tree gives its second grafting map. Both
belong to the same associative algebra, which supplies their
compatibilities with further substitutions.

\begin{construction}[The multikernel of a tree]
\label{con:tree-multikernel}
Let $\mathsf t$ be a finite planar rooted tree with ordered leaves
$a_1,\ldots,a_n$ and root output $b$. Label its regions by the
lower objects, so that an edge spanning consecutive regions has
its prescribed category $\mathcal A(x_i,x_j)$. Define
\begin{equation}\label{eq:tree-multikernel}
 M_{\mathsf t}(\mathbf a;b)
 :=
 \int^{(c_e)_{e\in E_{\mathrm{int}}(\mathsf t)}}
 \prod_{v\in V(\mathsf t)}
 M\bigl((c_e)_{e\in\operatorname{in}(v)};c_{\operatorname{out}(v)}\bigr).
\end{equation}
The external edges have the fixed labels $\mathbf a,b$.
The edge with no vertex is assigned the unary mapping-space kernel.
Substitution gives
$\mu_{\mathsf t}:M_{\mathsf t}(\mathbf a;b)\to M(\mathbf a;b)$.
This formula computes every finite-height composite directly.
\end{construction}

\begin{proposition}[Tree substitution and change of operation categories]
\label{prop:tree-substitution}
The kernels~\eqref{eq:tree-multikernel} and their maps
$\mu_{\mathsf t}$ are compatible with contraction of any collection
of internal edges. Given functors
$H_{x,y}:\mathcal A(x,y)\to\mathcal B(x,y)$, a map of normalized
multikernel systems over $H$ consists of natural maps
\begin{equation}\label{eq:multikernel-map}
 M(\mathbf a;b)\longrightarrow N(H\mathbf a;Hb)
\end{equation}
compatible with units and grafting. Such maps act functorially on
every tree kernel.
\end{proposition}
\begin{proof}
\leavevmode
\begin{enumerate}
\item \textbf{Compute a contraction by its internal variables.}
Partition the vertices of $\mathsf t$ into the subtrees to be
contracted. Fubini rewrites~\eqref{eq:tree-multikernel} as the
coend of the corresponding subtree kernels. Apply their grafting
maps and then integrate over the remaining edges. The associative
algebra structure identifies this composite with the original
$\mu_{\mathsf t}$, coherently for all further contractions.
\item \textbf{Send every vertex and internal edge through the map.}
Apply~\eqref{eq:multikernel-map} at each vertex. Change each coend
variable along its functor $H_{x,y}$ using the induced map of
two-sided bars. This gives
$M_{\mathsf t}(\mathbf a;b)\to N_{\mathsf t}(H\mathbf a;Hb)$.
Compatibility at each vertex implies compatibility with contraction.
Maps of the same bar diagrams prove compatibility with composition
of $H$ and with parameter families of such maps.\qedhere
\end{enumerate}
\end{proof}

\begin{construction}[Freely generated binary multikernels]
\label{con:free-binary-multikernels}
Let $B(a_1,a_2;b)$ be any binary multicollection on $\mathcal A$.
For $n\geq2$, let $\operatorname{Bin}(n)$ be the set of finite planar
binary rooted trees with $n$ ordered leaves, and define
\[
 \mathsf F(B)(\mathbf a;b)
 =
 \coprod_{\mathsf t\in\operatorname{Bin}(n)}
 \int^{(c_e)_{e\in E_{\mathrm{int}}(\mathsf t)}}
 \prod_{v\in V(\mathsf t)}
 B(c_{v,1},c_{v,2};c_{\operatorname{out}(v)}).
\]
Put $\mathsf F(B)(a;b)=\Map_{\mathcal A(x,y)}(a,b)$ and give it
empty nullary rows. Grafting binary trees, followed by Fubini for
their internal-edge coends, supplies the substitution maps.
Co-Yoneda supplies substitution at unary vertices. The common
tree diagrams used in Proposition~\ref{prop:substitution-monoidal}
give all coherences. Thus $\mathsf F(B)$ is a normalized
multikernel system constructed from arbitrary binary data.

Its free universal property fixes the unary category and its unit.
Restriction along the two-leaf corolla identifies maps
$\mathsf F(B)\to M$ over the identity with maps of binary collections
$B\to M_2$. To construct the inverse, label the vertices of each tree
by the given map, pass to the coend along its internal edges, and
apply $\mu_{\mathsf t}^M$. The unary component is the unit of $M$.
Compatibility with a grafting is compatibility of $\mu^M$ with the
corresponding tree contraction; for a chain of graftings use the same
chain of contractions in $M$. Restriction of this extension to the
corolla is the original map. Conversely the summand of each tree in
$\mathsf F(B)$ is itself the grafting of its generating corollas:
the map from their coend to that summand is its identity by Fubini.
Naturality of an algebra morphism for this grafting therefore identifies
it on the summand with the described extension. On a chain of
contractions the same identification is the coherence of that algebra
morphism on the chain, so these identifications form a homotopy of
algebra maps. The construction applies to space-parameterized maps
throughout, proving the mapping-space equivalence.
\end{construction}

For example, take one lower object, its operation category
$\mathcal A=\one$, and the constant one-point binary kernel $B$.
Then $\mathsf F(B)_n$ is the discrete space of planar binary trees
with $n$ leaves. Its ternary row has two points, corresponding to
the two binary bracketings. That row has no corepresentative in
$\one$, whose unique representable functor is the point. This gives
an explicit system of independently unrepresented higher operations.

\subsection{Recovery from universal composition operations}
\label{subsec:multikernel-recovery}

An output represents a row when every operation in that row factors
uniquely through its universal operation. To recover associative
composition, we also need to understand what happens when universal
operations are substituted into one another. The tree formula lets
us test this without making separate coherence choices.

\begin{definition}[Represented rows and substitution stability]
\label{def:multikernel-representability}
A row $M(\mathbf a;-)$ is corepresented by $m(\mathbf a)$ if
there is an equivalence, natural in $b$,
\begin{equation}\label{eq:multikernel-corepresentation}
 M(\mathbf a;b)\simeq
 \Map_{\mathcal A(x_0,x_n)}(m(\mathbf a),b).
\end{equation}
Its universal operation
$u_{\mathbf a}\in M(\mathbf a;m(\mathbf a))$ corresponds to the
identity. A complete corepresenting record includes this operation
and its universal property.

The universal operations are \emph{stable under tree substitution}
if the grafted operation obtained from any finite tree of universal
operations again corepresents the resulting row.
\end{definition}

\begin{theorem}[Represented recovery for independent multikernels]
\label{thm:independent-represented-recovery}
Let $M$ be a normalized multikernel system with every row
corepresentable, including the nullary rows.
\begin{enumerate}
\item The representatives assemble into functors
\[
 m_{\mathbf x}:
 \prod_{i=1}^n\mathcal A(x_{i-1},x_i)
 \longrightarrow\mathcal A(x_0,x_n),
 \qquad m_{(x,y)}\simeq\id_{\mathcal A(x,y)}.
\]
For each tree $\mathsf t$, write $m_{\mathsf t}$ for their iterated
application along that tree. Grafting recovers coherent comparisons
\begin{equation}\label{eq:oplax-composition-comparison}
 \alpha_{\mathsf t}:m_{\mathbf x}(\mathbf a)
 \longrightarrow m_{\mathsf t}(\mathbf a).
\end{equation}
Thus the represented rows recover an oplax composition structure.
\item The following conditions are equivalent: all grafting maps
$\mu_{\mathsf t}$ are equivalences; all comparisons
$\alpha_{\mathsf t}$ are equivalences; and universal operations
are stable under tree substitution. Under these conditions the
functors $m_{\mathbf x}$ give unital associative composition up to
coherent equivalence.
\item If $N$ is another system with represented rows, maps
$M\to N$ over $H$ correspond to coherent lax comparisons
\begin{equation}\label{eq:multikernel-lax-comparison}
 m^N_{\mathbf x}(H\mathbf a)
 \longrightarrow Hm^M_{\mathbf x}(\mathbf a).
\end{equation}
The correspondence retains the full categories of transformations.
A map preserves the universal operations in a specified family
of rows exactly when the corresponding comparisons are equivalences.
\end{enumerate}
\end{theorem}
\begin{proof}
\leavevmode
\begin{enumerate}
\item \textbf{Assemble the representing objects.}
Regard the rows as a functor from the opposite of the input
category to
$\Fun(\mathcal A(x_0,x_n),\Sps)$.
Covariant Yoneda
$\mathcal A(x_0,x_n)^{\op}\to
\Fun(\mathcal A(x_0,x_n),\Sps)$ is fully faithful.
Corepresentability places the row functor in its essential image,
so it lifts uniquely up to a contractible space of complete lifts.
Taking opposites gives $m_{\mathbf x}$ with all its naturality.
Unary normalization identifies $m_{(x,y)}$ with the identity.
Nullary rows supply the unit objects.
\item \textbf{Compute a represented tree.}
Starting with the bottom vertices of $\mathsf t$, substitute
\eqref{eq:multikernel-corepresentation} into
\eqref{eq:tree-multikernel}. Co-Yoneda eliminates each internal
edge variable. Repeating upwards gives a natural equivalence
\begin{equation}\label{eq:represented-tree-kernel}
 M_{\mathsf t}(\mathbf a;b)
 \simeq\Map(m_{\mathsf t}(\mathbf a),b).
\end{equation}
Under this equivalence, $\mu_{\mathsf t}$ is a map
\[
 \Map(m_{\mathsf t}(\mathbf a),-)
 \longrightarrow\Map(m_{\mathbf x}(\mathbf a),-).
\]
Covariant Yoneda gives precisely
\eqref{eq:oplax-composition-comparison}. Its direction follows from
precomposition. Applying this fully faithful functor to the entire
contraction diagram transfers all the algebra coherences at once.
\item \textbf{Recognize universal substituted operations.}
The identity of $m_{\mathsf t}(\mathbf a)$ represents the tree
of universal operations in~\eqref{eq:represented-tree-kernel}.
Its image under $\mu_{\mathsf t}$ is the operation corresponding
to $\alpha_{\mathsf t}$. This operation is universal exactly when
$\alpha_{\mathsf t}$ is an equivalence, which by Yoneda is exactly
when $\mu_{\mathsf t}$ is an equivalence. This proves all three
conditions equivalent. In that case the coherent comparisons from
Step~2 can be inverted, giving associative and unital composition.
\item \textbf{Recover maps by evaluating universal operations.}
Evaluate~\eqref{eq:multikernel-map} at
$b=m^M_{\mathbf x}(\mathbf a)$ and at $u_{\mathbf a}$.
Corepresentability in $N$ turns its image into
\eqref{eq:multikernel-lax-comparison}. Conversely, such a comparison
sends an operation represented by
$v:m^M_{\mathbf x}(\mathbf a)\to b$ to the composite
\[
 m^N_{\mathbf x}(H\mathbf a)
 \longrightarrow Hm^M_{\mathbf x}(\mathbf a)
 \xrightarrow{H(v)}Hb.
\]
Yoneda identifies the two procedures as inverse. Applied to each
tree, it identifies compatibility with grafting with the lax
composition axioms, including units.
\item \textbf{Recover transformations and strongness.}
A transformation between maps over $H$ and $G$ consists of natural
transformations $\eta_{x,y}:H_{x,y}\to G_{x,y}$ whose output action
on each operation agrees coherently with the input actions.
For represented rows, Step~4 identifies this diagram with
\[
 \eta_{m^M(\mathbf a)}\,\alpha_H
 \simeq
 \alpha_G\,m^N(\eta_{\mathbf a}).
\]
Use the same Yoneda equivalence on compatible diagrams indexed
by $[q]$, for every $q$. Naturality makes these equivalences
commute with all simplicial operators; they therefore identify the
complete nerves of the transformation categories. Finally, the
image of $u_{\mathbf a}$ is universal exactly when its representing
arrow~\eqref{eq:multikernel-lax-comparison} is an equivalence.\qedhere
\end{enumerate}
\end{proof}

\paragraph{Representability does not imply substitution stability.}
Use one lower object and let its operation category be
$\mathcal A=(\mathbb N,\leq)$. Put $c_1=0$, $c_n=1$ for $n\ne1$,
and $m_n(a_1,\ldots,a_n)=\sum_i a_i+c_n$ including $n=0$.
Set $M(\mathbf a;b)=\mathcal A(m_n(\mathbf a),b)$. For a partition
$n=n_1+\cdots+n_k$ the grafting exists since
\[
 c_n\leq c_k+\sum_{j=1}^k c_{n_j}.
\]
Indeed the right side vanishes only for $k=n_1=1$. This includes
empty blocks and nullary roots. Whenever the input operations exist,
their inequalities compose with this inequality to give the output
operation. This is balanced and functorial in the intermediate
variables; all unit, associativity and higher diagrams commute since
each output space is subterminal. The unary row is the hom functor.
Every row is represented, but at three zero inputs the flat
representative is $1$ and the binary-tree representative is $2$.
At output $1$ the flat operation exists and the tree operation space
is empty. Thus even representation in every arity leaves a genuine
substitution condition. The representative comparison is $1\to2$,
in precisely the direction of \eqref{eq:oplax-composition-comparison}.

The independent algebra itself also cannot be recovered from the unary
category: on the sole unary category $*$, the free nonsymmetric operad
on a binary point has two ternary operations, while the positive
associative operad has one. Both have the same empty nullary, terminal
unary and terminal binary spaces.

\subsection{Coherent degrees over the original doctrine}
\label{subsec:graded-multikernels}

We now put the base degrees back into the construction. Choose
$t_x\in T$ for every $x\in X$ and set
$\mathcal B(x,y)=\Map_T(t_x,t_y)$, regarded as a category which is
a space. The complete Segal nerve of $T$ supplies its coherent
composition functors. Equivalently, restrict the category enriched
in spaces presented by $T$ to the specified objects. Write
$m^T(\mathbf f)$ for this coherent composition.

\begin{construction}[The degree multikernel and its fibers]
\label{con:degree-multikernel}
Define the represented multikernel
\[
 \mathcal P_T(\mathbf f;g)
 :=\Map_{\mathcal B(x_0,x_n)}(m^T(\mathbf f),g).
\]
Composition in $T$ and co-Yoneda give its normalized associative
multikernel structure. A $T$-graded system consists of functors
$d_{x,y}:\mathcal A(x,y)\to\mathcal B(x,y)$ and a map of
normalized multikernel systems $M\to\mathcal P_T$ over $d$.
For a coherent degree operation
$\sigma\in\mathcal P_T(\mathbf f;g)$, take the homotopy fiber
\begin{equation}\label{eq:graded-independent-multikernel}
 M_\sigma(\mathbf a;b)
 :=\fib_\sigma\bigl(M(\mathbf a;b)
       \longrightarrow\mathcal P_T(\mathbf f;g)\bigr),
 \qquad d\mathbf a=\mathbf f,\quad db=g.
\end{equation}
Thus $M_\sigma$ is contravariant on the input fibers
$\mathcal A_{f_i}$ and covariant on $\mathcal A_g$.
\end{construction}

The map of systems is additional information when the operations are
independently specified. For example, take one lower object, unary
category $*$, and the terminal nonsymmetric operad. In the ordinary
category of finite sets put $t=\{0,1\}$ and send the unique unary
object to its transposition $w:t\to t$. This gives the specified
functor into the discrete space of maps $t\to t$. A binary operation
would require an identification $w^2=w$, which is impossible. Thus
the functor on degrees alone does not construct a grading.

The degree includes the entire compositional identification
$m^T(\mathbf f)\simeq g$. As $\sigma$ varies, the fibers in
\eqref{eq:graded-independent-multikernel} vary functorially over
its space. This records the higher base homotopies together with
the operations they grade.

\begin{proposition}[Graded substitution and recovery]
\label{prop:graded-independent-recovery}
Tree substitution in a $T$-graded system carries a family of vertex
degrees to its composite degree in $\mathcal P_T$.
The constructions of Propositions~\ref{prop:substitution-monoidal}
and~\ref{prop:tree-substitution} therefore give coherent graded
substitution at every finite tree height.
If the rows of $M$ are represented and their universal operations
are stable under substitution, Theorem~\ref{thm:independent-represented-recovery}
recovers associative composition over $T$.
\end{proposition}
\begin{proof}
\leavevmode
\begin{enumerate}
\item \textbf{Compose the degrees in the same tree.}
The map $M\to\mathcal P_T$ commutes with grafting. Its restriction
to fixed vertex-degree fibers therefore lands in the fiber over
the grafted degree. The functors $d_{x,y}$ carry each internal
edge to its base degree; the induced maps of bar diagrams preserve
the balancing relations. Tree contractions commute with these
maps by Proposition~\ref{prop:tree-substitution}. Taking homotopy
fibers retains their full coherent parameter diagrams.
\item \textbf{Read off the degree of a universal operation.}
The image of $u_{\mathbf a}$ is an element of
\[
 \mathcal P_T(d\mathbf a;d\,m(\mathbf a))
 =
 \Map_{\mathcal B(x_0,x_n)}
       (m^T(d\mathbf a),d\,m(\mathbf a)).
\]
It identifies the degree of the representing operation with the
composite input degree. These identifications are coherent because
the grading is a multikernel map. Together with the invertible
comparisons supplied by
Theorem~\ref{thm:independent-represented-recovery}, they give the
claimed composition over $T$.\qedhere
\end{enumerate}
\end{proof}

Taking $T=\RNU$ supplies the degree system for the independent
composition operations over the NU doctrine. The lower marking
$\tau_{\NU}$ remains the marking used for transport and dependent
base change in Section~\ref{sec:theorization} and in the iterated
microcosm of Section~\ref{sec:stratified-microcosm}.
At the level of the new rows, preservation of a specified universal
operation is expressed by the equivalence condition in
Theorem~\ref{thm:independent-represented-recovery}(3).
These are explicit conditions on the operations and their maps.

\subsection{Unrepresented categorical strata and intrinsic limits}
\label{subsec:unrepresented-categorical-strata}

The multikernel construction theorizes composition locally. A
complementary iteration forms complete category objects of theorized
structures, retaining unrepresented operations in every external degree.
To construct these objects, we first establish limits of
marked-cocartesian objects without requiring finite cocartesian
exactness. Exactness is then imposed degreewise.

\begin{construction}[Marked-cocartesian ambient category]
\label{con:marked-cocart-ambient}
Let $\mathscr C_\tau(T)$ be the $\infty$-category whose objects are
bounded functors $p:E\to T$ admitting globally $p$-cocartesian lifts
of every $\tau$-marked base arrow and whose morphisms are functors
over $T$ preserving these marked cocartesian arrows. It is obtained
as the corresponding (generally non-full) subcategory of $(\Cat)_{/T}$.
The finite cocartesian-exact objects form the subcategory underlying
$\Thz_\tau(T)$.
\end{construction}

\begin{theorem}[Limits of marked-cocartesian objects]
\label{thm:marked-cocart-limits}
The category $\mathscr C_\tau(T)$ admits bounded small limits, created
by the underlying functors in $(\Cat)_{/T}$. For an arrow of such a
limit whose degree is marked, cocartesianness is equivalent to
cocartesianness of all its projections.
\end{theorem}
\begin{proof}
\leavevmode
\begin{enumerate}
\item \textbf{Test slice limits by lifts with specified base.}
For $\sigma:D\to T$ and $p:E\to T$ put
\[
 \mathcal D_\sigma(p)=\Fun(D,E)\times_{\Fun(D,T)}\{\sigma\}.
\]
For each $n$ currying gives a natural equivalence
\[
 N_n\mathcal D_\sigma(p)\simeq
 \Map_{\Cat/T}\bigl((D\times[n],\sigma\operatorname{pr}_D),(E,p)\bigr).
\]
Indeed, both sides are the fiber of the mapping-space map for
$D\times[n]\to E\to T$ at $\sigma\operatorname{pr}_D$.
Each level therefore preserves all limits in the slice. For
$p=\lim_{i\in I}p_i$ in $\Cat/T$ we obtain
$\mathcal D_\sigma(p)\simeq\lim_i\mathcal D_\sigma(p_i)$, naturally
in $\sigma$. Taking $\sigma$ to be a vertex gives
$E_t\simeq\lim_i(E_i)_t$. For $\sigma=g:t\to r$ take the fiber of
the two endpoint evaluations at $(y,z)\in E_t\times E_r$.
This fiber is a space: it consists of arrows $y\to z$ with their
identification of degree with $g$, and endpoint-fixed transformations
are equivalences. More explicitly, on cores it is the iterated fiber
of the degree-one complete nerve over the endpoints and $g$, which
is exactly the defining normalized kernel. Commuting these limits gives
\[
 P_g^E(y,z)\simeq\lim_i P_g^{E_i}(y_i,z_i).
 \tag*{$(*)$}
\]
These are equivalences of entire families, since restriction to
endpoints and the degree maps are natural.

\item \textbf{Assemble the marked lifts without a cleavage.}
For marked $f:s\to t$ let $U_i(f)\subseteq\mathcal D_f(p_i)$ be
the full subcategory of globally $p_i$-cocartesian lifts. Source
evaluation $U_i(f)\to(E_i)_s$ is essentially surjective by the
definition of $\mathscr C_\tau(T)$. For $u:x\to y$ and $u':x'\to y'$
in $U_i(f)$ its mapping space is
\[
 \Map_{(E_i)_s}(x,x')\times_{P_f^{E_i}(x,y')}
                     \Map_{(E_i)_t}(y,y').
\]
The two maps are composition with $u'$ and precomposition with $u$:
a point of the pullback is precisely a vertical square over the
identity square of $f$, including its commutativity homotopy.
Precomposition with $u$ is an equivalence by its cocartesian
property in degree $\id_t$. Thus source evaluation is fully faithful.
These evaluation maps give a natural equivalence of $I$-diagrams,
because all transition maps preserve the indicated lifts.

Consequently $U(f):=\lim_iU_i(f)\to E_s$ is an equivalence.
Its map into $\mathcal D_f(p)$ is fully faithful with replete image
the lifts whose projections belong to $U_i(f)$ for every $i$:
full subcategories restrict objects and inherit all mapping spaces,
so compatible families of such objects and their coherences factor
through the limit of the full subcategories. In particular every
source admits a lift in $U(f)$.

\item \textbf{Prove global universality and detect it.}
Let $u:x\to y$ lie in $U(f)$ and take any $z\in E_r$ and
$g:t\to r$. The comparison
$P_g^E(y,z)\to P_{gf}^E(x,z)$ is, by $(*)$, the limit of the
comparisons induced by $u_i$. This identification respects
precomposition: restriction of a lifted composable pair to its
long edge commutes with every projection. Each component is an
equivalence by global cocartesianness of $u_i$. Varying $g$ over
$\Map_T(t,r)$ proves that \eqref{thz:eq:cocart} is an equivalence,
by the fiber criterion for maps of spaces over a fixed space
\cite[Prop.~5.1.2.1, dual]{Kerodon}. Conversely any globally
cocartesian lift with the same source and degree differs from this
one by a vertical equivalence, by Lemma~\ref{thz:lem:cocart-calculus}.
It belongs to the same replete subcategory $U(f)$. This proves
existence and componentwise detection for every marked $f$,
including the empty limit, where $p=\id_T$.

\item \textbf{Restrict the universal cone.}
For every $p'$ in $\mathscr C_\tau(T)$, the mapping space into $p$
in the ambient slice is the limit of the mapping spaces into $p_i$.
A map in this space preserves the marked cocartesian arrows precisely
when all its projections do, by the preceding step. The predicate is
invariant under equivalence of maps. Taking its union of components
on both sides gives the mapping-space universal property in
$\mathscr C_\tau(T)$. The projections themselves preserve the marked
lifts by componentwise detection. This proves creation of the limit.
\qedhere
\end{enumerate}
\end{proof}

\begin{corollary}[Limit-preserving intrinsic internalization]
\label{cor:intrinsic-internalization-limits}
The functor
\[
 \Int_\tau:\mathscr C_\tau(T)\longrightarrow\Cat,
 \qquad p\longmapsto\Fun_{/T}^{\tau\text{-}\coc}(T,E),
\]
preserves bounded small limits. Its restriction to finite
cocartesian-exact objects is the intrinsic internalization of
Definition~\ref{def:general-internalization}.
\end{corollary}
\begin{proof}
\leavevmode
\begin{enumerate}
\item \textbf{Compute sections and their marked predicate in the limit.}
The section category is $\mathcal D_{\id_T}(p)$.
Its complete nerve is represented levelwise in the slice, as in the
first step of Theorem~\ref{thm:marked-cocart-limits}, so it preserves
all slice limits. The predicate that a section
carry a marked arrow to a cocartesian one is detected on each
projection by Theorem~\ref{thm:marked-cocart-limits}. Taking the full
subcategories on that predicate commutes with the resulting limit.
Thus $\Int_\tau(\lim_i p_i)\simeq\lim_i\Int_\tau(p_i)$ as
full $\infty$-categories.
\qedhere
\end{enumerate}
\end{proof}

\begin{construction}[Higher categorical theorized objects]
\label{con:higher-unrepresented-strata}
For every $r\geq0$, define
\begin{equation}\label{eq:higher-unrepresented-strata}
 \mathsf{Thz}^{\langle r\rangle}_\tau(T)
 \subseteq\operatorname{Cat}^r(\mathscr C_\tau(T))
\end{equation}
to be the full subcategory on those complete $r$-fold category objects
$X$ for which each external multidegree
$X_{\mathbf m}=(p_{\mathbf m}:E_{\mathbf m}\to T)$
is finitely cocartesian exact. At $r=0$ this is the underlying
$\infty$-category of $\Thz_\tau(T)$.
Each multidegree retains its normalized, potentially unrepresented
operation kernels. The external face, degeneracy, Segal, and
completeness structure belongs to the complete category object in
$\mathscr C_\tau(T)$, with the fixed base marking $\tau$.
\end{construction}

\begin{theorem}[Higher intrinsic internalization and represented recovery]
\label{thm:higher-unrepresented-internalization}
Pointwise intrinsic internalization constructs, for every $r\geq0$,
\begin{equation}\label{eq:higher-unrepresented-internalization}
 \Int_\tau^{\langle r\rangle}:
 \mathsf{Thz}^{\langle r\rangle}_\tau(T)
 \longrightarrow\operatorname{Cat}^r(\Cat),
 \qquad
 \bigl(\Int_\tau^{\langle r\rangle}X\bigr)_{\mathbf m}
 =\Int_\tau(p_{\mathbf m}).
\end{equation}
Its external Segal, completeness, and mixed coherences are canonical.
On the degreewise represented objects with strong external structure
maps, and with strong morphisms, there is a natural equivalence
\[
 \mathsf{Thz}^{\langle r\rangle,\mathrm{rep,strong}}_\tau(T)
 \simeq\operatorname{Cat}^r(\Lex(T,\Cat)).
\]
In particular for $(T,\tau)=(\RNU,\tau_{\NU})$ the represented
subcategory is $\operatorname{Cat}^{r+1}(\TT)$.
Moreover the enlargement is genuine: the constant $r$-fold object on
any nonrepresented object of $\Thz_\tau(T)$ belongs to
$\mathsf{Thz}^{\langle r\rangle}_\tau(T)$ and remains nonrepresented.
\end{theorem}
\begin{proof}
\leavevmode
\begin{enumerate}
\item \textbf{Apply internalization to each categorical direction.}
Corollary~\ref{cor:intrinsic-internalization-limits} carries every
external Segal pullback and every completeness limit comparison to
an equivalence in $\Cat$. Applying the one functor $\Int_\tau$
degreewise preserves all simplicial operators and their mixed
identities, proving the first assertion. 
\item \textbf{Identify the represented categorical strata.}
On the strong represented
subcategory, Theorem~\ref{thz:thm:recovery} identifies the degreewise
functors and their maps with left-exact categorical diagrams. Its
straightening equivalence preserves and reflects the external
complete-Segal conditions. Indeed, with the sharp marking
Theorem~\ref{thm:marked-cocart-limits} computes limits of
cocartesian fibrations and strong maps in $\Cat/T$. The inclusion
of left-exact diagrams among all diagrams creates limits, since
limits commute with finite limits. Including the sharp marked
objects into $\mathscr C_\tau(T)$ thus computes the same limiting
cones, by the same theorem for $\tau$. Applying this to the
external structure maps proves the displayed equivalence. The NU specialization is
Theorem~\ref{thm:nu-categorical}. 
\item \textbf{Retain unrepresented examples in every stratum.}
Finally, a constant multisimplicial
object has identity Segal and completeness comparisons and retains the
underlying nonrepresented kernel in each degree.
\qedhere
\end{enumerate}
\end{proof}

\begin{proposition}[Higher unrepresented operation kernels]
\label{prop:higher-unrepresented-kernels}
For $X\in\mathsf{Thz}^{\langle r\rangle}_\tau(T)$, unstraightening
is not required to construct its higher operation categories: the
degreewise total categories form
$\mathbb E(X)\in\operatorname{Cat}^r(\Cat/T)$, and for any
$f:s\to t$ the pullback
\[
 \mathsf{Ar}^{(r)}_f(X)=
 \operatorname{Ar}_{\mathrm{lev}}(\mathbb E(X))
 \times_{\delta^r\Fun([1],T)}\delta^r\{f\}
\]
is a complete $r$-fold category object. Its fiber in external degree
$\mathbf m$ at $(x,y)$ is
\[
 P_f^{\mathbf m}(x,y)
 =\fib_f\bigl(\Map_{E_{\mathbf m}}(x,y)
 \longrightarrow\Map_T(s,t)\bigr).
\]
The complete string diagrams induce composition and all its external
coherences. When $f\in\tau$, the subcategory of globally cocartesian
arrows over $f$ is equivalent by source evaluation to the fiber
category at $s$, compatibly with every external simplicial operator.
\end{proposition}
\begin{proof}
\leavevmode
\begin{enumerate}
\item \textbf{Form the external arrow categories and their fibers.}
The forgetful functor $\mathscr C_\tau(T)\to\Cat/T$
preserves limits by Theorem~\ref{thm:marked-cocart-limits}; hence it
carries the external Segal and completeness comparisons of $X$ to
those of $\mathbb E(X)$. The fixed-degree arrow category is $\mathcal D_f(p_{\mathbf m})$.
The first step of Theorem~\ref{thm:marked-cocart-limits} shows directly
that this construction preserves slice limits.
In degree $\mathbf m$ the resulting category is
$\Fun([1],E_{\mathbf m})\times_{\Fun([1],T)}\{f\}$, whose
fiber at $(x,y)$ is the displayed normalized kernel. 
\item \textbf{Retain composition and marked evaluation.}
Functoriality
of the whole string diagram supplies composition and all mixed
simplicial compatibilities. For marked $f$, the universal lift
property, contractibility of lifts, and marked preservation by every
external operator give the source-evaluation equivalence.
\qedhere
\end{enumerate}
\end{proof}

\begin{remark}[Categorical and profunctor composition]
\label{rem:higher-bimodule-boundary}
For every multidegree of a higher unrepresented theorized object, the
normalized operation spaces have explicit iterated composition
profunctors and comparison maps, as constructed in
Theorem~\ref{thm:derived-higher-composition}. The external Segal and
completeness directions of
\eqref{eq:higher-unrepresented-strata} organize these constructions
without replacing the structural marking $\tau$ by a new one.
\end{remark}

\subsection{Retaining the operation data during realization}

The preceding construction keeps the operation categories, their
degrees, and their multikernels together. The next two calculations
explain this choice of data: they compare what is retained by a
total category and by the represented envelope.

\begin{lemma}[A semilattice obstruction to equalizers]
\label{lem:semilattice-total-category-obstruction}
Let $T$ be small and finitely complete, and let $F:T\to\Sps$ be
left exact and discrete-valued, equipped naturally with meet-semilattice
structures with top. Suppose $a\ne1$ in some $F(t)$.
Let $BM$ denote the one-object category with endomorphism monoid
$(M,\wedge,1)$, and define $K(t)=BF(t)$. Then $K$ is left exact,
but $\int_TK$ has no equalizer of the vertical endomorphisms $a$ and
$1$ of its object over $t$.
\end{lemma}
\begin{proof}
\leavevmode
\begin{enumerate}
\item \textbf{Verify left exactness of the categorical model.}
Only $1$ is invertible in a meet-semilattice monoid. Thus the complete
nerve satisfies $N_nK(t)=F(t)^n$ as discrete spaces, with $N_0K(t)=*$.
Every nerve level is left exact, so $K$ is left exact.


\item \textbf{Write the equalizer universal property in coordinates.}
Write $E=\int_TK$. There is one object over each base object, and
\[
 \Map_E(s,t)\simeq\Map_T(s,t)\times F(t),
 \qquad
 (g,c)\circ(f,b)=(gf,c\wedge F(g)(b)).
\]
The chosen object and the arrows $(f,1)$ form a cocartesian section of
$E\to T$. Suppose an equalizer of $a$ and $1$ existed, with arrow
$h=(f,b):s\to t$ and its equalizing homotopy. Discreteness of $F(t)$
gives $a\wedge b=b$. Put
\[
 L=\Eq_T(\id_t,\id_t).
\]
The base arrow and the base image of the homotopy give $v:s\to L$.
Restricting the equalizer universal property along the cocartesian
section, for every $z$ we obtain
\begin{equation}\label{eq:semilattice-equalizer-test}
 \Map_T(z,s)\times F(s)
 \xrightarrow{\simeq}
 \Map_T(z,L)\times\mathord\downarrow a,
 \qquad
 \mathord\downarrow a=\{d\in F(t):a\wedge d=d\}.
\end{equation}
The displayed map is the product of $\Map_T(z,v)$ and
\[
 \varphi:F(s)\to\mathord\downarrow a,
 \qquad c\longmapsto b\wedge F(f)(c).
\]
Indeed, the homotopy equalizer of the identity actions on the base
mapping space is its free-loop space, represented by $L$; the discrete
semilattice factor gives exactly $\mathord\downarrow a$.


\item \textbf{Derive the contradiction from the two factors.}
At $z=s$, both mapping-space factors in
\eqref{eq:semilattice-equalizer-test} are inhabited, by $\id_s$ and
$v$. The product map being an equivalence forces $\varphi$ to be a
bijection. For every $z$ the same equivalence then forces
$\Map_T(z,v)$ to be an equivalence, so Yoneda gives $s\simeq L$.
Since $F$ is discrete-valued and left exact, the projection $L\to t$
induces a bijection $F(L)\to F(t)$. Hence $F(f)$ is a bijection.
Injectivity of $\varphi$ would imply injectivity of meet with $b$ on
$F(t)$. But $b\wedge1=b=b\wedge a$ and $a\ne1$, a contradiction.
\qedhere
\end{enumerate}
\end{proof}

\begin{proposition}[Failure of total-category closure at the NU base]
\label{prop:NU-total-category-obstruction}
There exists a represented object of
$\Thz_{\tau_{\NU}}(\RNU)$ whose total category is not finitely
complete. It has cocartesian lifts of every base arrow.
\end{proposition}
\begin{proof}
\leavevmode
\begin{enumerate}
\item \textbf{Construct a semilattice object in type theories.}
Let $A=(\{0<1\})_{\minrep}$ be the two-element poset with its minimal
theory structure. Meet $A\times A\to A$ preserves finite limits, and
its unit selects the terminal object $1$. These are theory morphisms
and make $A$ a meet-semilattice object in $\TT$.


\item \textbf{Obtain its discrete NU model.}
Transport this structure through
$\Lex(\RNU,\Sps)\simeq\TT$ to obtain $F:\RNU\to\Sps$.
Each $F(t)$ is discrete: the covariant representable $h_t$ corresponds
to a theory $A_t$, and
\[
 F(t)\simeq\Map_{\TT}(A_t,A).
\]
A functor category with poset target has no nonidentity natural
equivalences; the subspace of theory morphisms is therefore
$0$-truncated. The transported meet and unit make these sets natural
meet-semilattices. The model $F$ is not terminal, since $A$ is not the
terminal theory. Thus some $F(t)$ contains an element different from $1$.


\item \textbf{Form the represented operation theory.}
Apply Lemma~\ref{lem:semilattice-total-category-obstruction}.
The functor $K=BF$ is left exact, so its cocartesian unstraightening is
finitely cocartesian exact by Lemma~\ref{thz:lem:exact}. Being
represented, it satisfies every marked-lift condition, in particular the
NU one. Its total category lacks the indicated equalizer.
\qedhere
\end{enumerate}
\end{proof}

The construction of Section~\ref{subsec:graded-multikernels} keeps the
finite-limit doctrine as the degree base and the operation categories
as separate inputs. In this example the former retains its finite
limits while the latter retain the semilattice operations.

\subsection{The resulting operation data}
\label{subsec:comparison-higher-theories}

The constructions now give explicit descriptions of three kinds of
composition data. For an operation category over $T$, the bar formula
of Theorem~\ref{thm:derived-higher-composition} computes balanced
strings and their comparison with composition. For a family of
categories of lower operations, Definition~\ref{def:independent-multikernels}
supplies independent higher operation spaces, and
Construction~\ref{con:tree-multikernel} computes their substitution
along any finite planar tree. Construction~\ref{con:free-binary-multikernels}
produces concrete unrepresented systems from arbitrary binary kernels.
Finally, Construction~\ref{con:degree-multikernel} places these operations
over the coherent simplices of the original doctrine.

The recovery statements can be read directly from the universal
operations. Theorem~\ref{thm:independent-represented-recovery} identifies
the representing composition functors, determines the direction of
their comparisons, and recognizes invertibility by stability under
substitution. Its map calculation recovers lax comparisons and their
full transformation categories. The categorical NU strata retain
their intrinsic internalization and represented recovery from
Theorem~\ref{thm:higher-unrepresented-internalization} and
Corollary~\ref{cor:stratified-represented-recovery}.

For the completed NU source, the corresponding global construction is
the relative cellular resolution of
Corollary~\ref{cor:NU-relative-completion}. Its envelope colimit
computes the walking classifier, while its marked-section limit
computes internal theories and every extension category. The local
operation formulas and the global classifiers thus each retain the
source data used in their universal properties.

\section{The iterated NU microcosm}
\label{sec:stratified-microcosm}

The walking NU theory represents an internal theory together with its
ordinary morphisms. Once those morphisms are present, one can ask for
families of them, their compositions, and further categorical comparisons.
The representing source should account for these higher degrees as well.

The key point is already available: the internal-object functor is
represented and hence preserves limits. It can therefore be applied to
complete category objects, and then applied in further categorical
directions. Compact syntax identifies the arities of these semantics;
free complete-category evaluation constructs the higher comparison
objects. The direct marked-operation construction verifies that they
recover the original NU operations in every degree.

The preceding section constructed independent higher composition
operations and, separately, categorical strata of operation theories
with potentially unrepresented kernels. The latter admit intrinsic
internalization in every external degree. We now identify the represented
NU semantics and construct their higher universal comparison objects.
Throughout, Segal and completeness conditions govern the ambient
categorical directions, while the original NU marking governs structural
strongness of an internal object.

We use the iterated categorical conventions fixed in Section~\ref{sec:higher-composition}.

\subsection{Iterated NU semantics}
\label{subsec:stratified-semantics}

The first step is to identify the ambient semantic category at each
stage. Finite limits and complete-category conditions can both be tested
pointwise, so adding a categorical direction commutes with forming models
of the NU doctrine. This interchange supplies the semantic induction.

\begin{theorem}[Iterated categorical NU classification]
\label{thm:stratified-NU-semantics}
For every $r\geq0$ there is a canonical equivalence
\begin{equation}\label{eq:stratified-NU-semantics}
 \operatorname{Cat}^r(\TT)
 \simeq
 \Lex\bigl(\RNU,\operatorname{Cat}^r(\Sps)\bigr).
\end{equation}
It is natural on ordinary morphisms and higher transformations, and
commutes with adding a complete categorical direction and with the
successive discrete-category embeddings.
\end{theorem}
\begin{proof}
\leavevmode
\begin{enumerate}
\item \textbf{Interchange left exactness and complete category objects.}
For a small finitely complete $T$ and a complete $\infty$-category
$\mathcal A$, the full subcategory $\Lex(T,\mathcal A)$ is closed
under limits, which are computed pointwise.  A simplicial diagram
in it is therefore a bifunctor $T\times\Delta^{\op}\to\mathcal A$
left exact in the first variable.  Its Segal comparisons and its
walking-isomorphism completeness comparison are limits in
$\mathcal A$, tested pointwise on $T$.  Interchanging the two
variables consequently constructs
\[
 \operatorname{Cat}(\Lex(T,\mathcal A))
 \simeq\Lex(T,\operatorname{Cat}(\mathcal A)).
\]
The functors in both directions retain the same bifunctor and all its
natural transformations.  The construction is compatible with further
simplicial variables.
\item \textbf{Iterate the NU equivalence.}
The case $r=0$ is
$\Th(\RNU)\simeq\TT$ from
Theorem~\ref{thm:nu-categorical}.  Apply the interchange of Step~1
inductively.  The constant-simplicial-object inclusions are pointwise
and are carried to the corresponding inclusions by this same
interchange.  This gives all the asserted naturality and compatibility.
\qedhere
\end{enumerate}
\end{proof}

\subsection{Compact semantics and its finite-limit doctrine}
\label{subsec:compact-semantic-syntax}

To describe a categorical semantic stage by a small doctrine, it is useful
to ask which observations determine its objects. Compact objects provide
such observations, and finite colimits among them become finite-limit
relations in the opposite category. The following construction records
this passage before it is applied to the successive NU stages.

Let \(\mathcal A_T\) be a compactly generated presentable \(\infty\)-category
and let \(\mathcal A_T^\omega\) be its full subcategory of compact objects.
Compact generation is a source-specific property and must be proved or cited
where this construction is invoked.

\begin{lemma}[Finite-colimit closure of compact objects]\label{old:c8s7:lemma1}
\label{rn:compact-finite-colimits}
The essentially small subcategory \(\mathcal A_T^\omega\) is closed under
finite colimits.
\end{lemma}

\begin{proof}
\leavevmode
\begin{enumerate}
\item \textbf{Commute finite limits with filtered colimits.} For a finite diagram of compact objects $(c_i)$ and a filtered diagram
$(X_j)$, the equivalences
\[
 \Map(\colim_i c_i,\colim_jX_j)
 \simeq\lim_i\colim_j\Map(c_i,X_j)
 \simeq\colim_j\Map(\colim_i c_i,X_j)
\]
prove compactness. The middle interchange is commutation of filtered
colimits and finite limits in spaces.\qedhere
\end{enumerate}
\end{proof}

\begin{construction}[Compact semantic syntax]\label{old:c8s7:construction1}
\label{rn:compact-syntax}
Define $\widehat T:=((\mathcal A_T^\omega)^{\op})_{\minrep}.$ Lemma~\ref{rn:compact-finite-colimits} makes \(\widehat T\) finitely complete,
and the minimal representable-map structure makes it an \(\infty\)-type theory.
\end{construction}

\begin{proposition}[Semantic recognition from compact generation]\label{old:c8s7:proposition1}
\label{rn:semantic-recognition}
There are canonical equivalences
\[
 \Th(\widehat T)
 \simeq
 \Ind(\mathcal A_T^\omega)
 \simeq
 \mathcal A_T
\]
and $\Lex(\widehat T,\Cat) \simeq \operatorname{Cat}(\mathcal A_T).$ \end{proposition}

\begin{proof}
\leavevmode
\begin{enumerate}
\item \textbf{Apply compact generation degreewise.} Apply Theorem~\ref{thm:Th-filtered-colimit-generation} to
\(((\mathcal A_T^\omega)^{\op})_{\minrep}\) and use the compact-generation
equivalence. Applying it degreewise to complete Segal objects gives the categorical statement, since finite limits and the Segal and completeness tests are computed pointwise.\qedhere
\end{enumerate}
\end{proof}

\subsection{Compact doctrines and categorical arities}
\label{subsec:stratified-compact-doctrines}

The next calculation both supplies the new doctrines and specifies their
categorical arities.  It uses the enriched left-exact reflector already
constructed in Lemma~\ref{thz:lem:lex-reflector}.

\begin{lemma}[Compact categorical evaluation generators]
\label{lem:stratified-compact-generators}
For a small finitely complete $T$, put $\mathcal L_T=\Lex(T,\Cat)$.
For $t\in T$ and $m\geq0$ set
\[
 G^T_{t,m}:=L_{\mathrm{lex},T}(h_t\times[m]),
 \qquad h_t=T(t,-),
\]
where the pointwise product with $[m]$ is formed in
$\Fun(T,\Cat)$.  The objects $G^T_{t,m}$ are compact and jointly detect equivalences.
Let $\mathcal G_T$ be their idempotent-complete finite-colimit closure
in $\mathcal L_T$. Then $\mathcal G_T=\mathcal L_T^\omega$ and
$\mathcal L_T\simeq\Ind(\mathcal G_T)$; in particular, every object
is a filtered colimit of objects of this explicitly constructed closure.
Every compact object of $\mathcal L_T$ is a retract of a finite colimit
of the $G^T_{t,m}$.  Their evaluation formula is
\begin{equation}\label{eq:stratified-compact-evaluation}
 \Map_{\mathcal L_T}(G^T_{t,m},K)
 \simeq\Map_{\Cat}([m],K(t))=N_mK(t).
\end{equation}
\end{lemma}
\begin{proof}
\leavevmode
\begin{enumerate}
\item \textbf{Construct the representing objects.}
By enriched reflection and copowered Yoneda,
\[
 \Map_{\mathcal L_T}(L_{\mathrm{lex},T}(h_t\times[m]),K)
 \simeq\Map_{\Fun(T,\Cat)}(h_t\times[m],K)
 \simeq\Map_{\Cat}([m],K(t)).
\]
These equivalences are natural in $t,m,K$.
\item \textbf{Verify compactness.}
Filtered colimits in $\Cat$ commute with finite limits and are
computed degreewise under the complete Segal nerve.  Consequently,
filtered colimits of left-exact diagrams $T\to\Cat$ are left exact
and are computed pointwise in $\mathcal L_T$.  For every finite
ordinal $[m]$, the evaluation $N_m=\Map_{\Cat}([m],-)$ preserves
filtered colimits: on complete Segal nerves it is evaluation at the
fixed simplicial degree $m$.  Thus, for a filtered diagram
$(K_i)$ in $\mathcal L_T$,
\begin{align*}
 \Map(G^T_{t,m},\colim_iK_i)
 &\simeq\Map_{\Cat}([m],\colim_iK_i(t))\\
 &\simeq\colim_i\Map_{\Cat}([m],K_i(t))
 \simeq\colim_i\Map(G^T_{t,m},K_i).
\end{align*}
This proves compactness.
\item \textbf{Detect equivalences and identify the compact closure.}
A map $K\to L$ induces equivalences on maps out of every
$G^T_{t,m}$ precisely when $N_mK(t)\to N_mL(t)$ is an equivalence
for every $t,m$.  Full faithfulness of the complete Segal nerve then
makes $K\to L$ an equivalence.  The reflector makes $\mathcal L_T$
presentable, so the compact-generation and Ind-recognition results
of~\cite[\S5.3.5]{Lurie2009} identify $\mathcal L_T$ with the
Ind-completion of the idempotent-complete finite-colimit closure of
these compact generators.  This proves both generation and the
assertion about all compact objects.\qedhere
\end{enumerate}
\end{proof}

\begin{proposition}[Compact doctrines at every stage]
\label{prop:stratified-compact-doctrines}
Every $\operatorname{Cat}^r(\TT)$ is compactly generated and
presentable.  Put $T_0=\RNU$ and, for $r\geq1$, construct
\begin{equation}\label{eq:stratified-compact-doctrines}
 T_r:=\bigl(\operatorname{Cat}^r(\TT)^\omega\bigr)^{\op}_{\minrep},
\end{equation}
using a small equivalent category of compact objects at the fixed
relative size.  Then $T_r$ is an $\infty$-type theory and there are
natural equivalences
\begin{equation}\label{eq:stratified-doctrine-recognition}
 \Th(T_r)\simeq\operatorname{Cat}^r(\TT),
 \qquad
 \Lex(T_r,\Cat)\simeq\operatorname{Cat}^{r+1}(\TT).
\end{equation}
For $r\geq1$ the representable arrows of $T_r$ are exactly its
equivalences, as constructed in \eqref{eq:stratified-compact-doctrines};
$T_0$ retains the original NU representable-map structure.
\end{proposition}
\begin{proof}
\leavevmode
\begin{enumerate}
\item \textbf{Start with filtered generation of the NU theory.}
Theorems~\ref{thm:nu-categorical} and
\ref{thm:Th-filtered-colimit-generation} give
\[
 \TT\simeq\Th(\RNU)\simeq\Ind(\RNU^{\op}).
\]
Thus the stage $r=0$ is compactly generated and presentable.
\item \textbf{Construct the successive compact doctrines.}
Suppose $\Th(T_{r-1})\simeq\operatorname{Cat}^{r-1}(\TT)$.
Theorem~\ref{thm:nu-categorical} gives
\[
 \operatorname{Cat}^r(\TT)
 \simeq\operatorname{Cat}(\Th(T_{r-1}))
 \simeq\Lex(T_{r-1},\Cat).
\]
Lemma~\ref{lem:stratified-compact-generators} proves compact generation.
Its compact objects are closed under finite colimits by
Lemma~\ref{rn:compact-finite-colimits}.  Hence their opposite is
finitely complete, and its minimal representable-map structure makes
$T_r$ an $\infty$-type theory by Definition~\ref{def:infty-type-theory}.
Compact Ind-recognition and Theorem~\ref{thm:Th-filtered-colimit-generation}
now give
\[
 \Th(T_r)
 \simeq\Ind\bigl(\operatorname{Cat}^r(\TT)^\omega\bigr)
 \simeq\operatorname{Cat}^r(\TT).
\]
Apply Theorem~\ref{thm:nu-categorical} again to get the category-valued
equivalence in \eqref{eq:stratified-doctrine-recognition}.  All
identifications are induced by universal functors and so are natural
in the complete categorical directions.\qedhere
\end{enumerate}
\end{proof}

\begin{proposition}[Categorical arities and their relations]
\label{prop:stratified-categorical-arities}
For $r\geq1$, transport the compact generators
$G^{T_{r-1}}_{t,m}$ through the preceding equivalences.  Their
opposites construct a coherent functor
\[
 a_r:T_{r-1}\times\Delta^{\op}\longrightarrow T_r,
 \qquad (t,[m])\longmapsto(G^{T_{r-1}}_{t,m})^{\op}.
\]
It is left exact in the $T_{r-1}$-variable, and its image generates
$T_r$ under finite limits and retracts.  With $G_{t,m}$ abbreviating
$G^{T_{r-1}}_{t,m}$, its canonical categorical relations include
\begin{align}
 G_{t,m}&\simeq
 \underbrace{G_{t,1}\amalg_{G_{t,0}}\cdots
 \amalg_{G_{t,0}}G_{t,1}}_{m\text{ copies}}
 &&(m\geq2),\label{eq:stratified-arity-Segal}\\
 \operatorname*{colim}_{\sigma:\Delta[k]\to\mathbf E}G_{t,k}
 &\simeq G_{t,0},\label{eq:stratified-arity-complete}
\end{align}
where $\mathbf E$ is the simplicial nerve of the walking-isomorphism
groupoid.  Taking opposites gives the corresponding existing limit
comparisons in $T_r$.
\end{proposition}
\begin{proof}
\leavevmode
\begin{enumerate}
\item \textbf{Assemble the arities and verify their variance.}
Covariant Yoneda $t\mapsto h_t$ is contravariant in $t$, while
$[m]\mapsto[m]$ is covariant in $[m]$.  Reflection and taking
opposites give the displayed functor with both variances correct.
For any finite limiting cone $\ell\to t_j$, left exactness of
$K:T_{r-1}\to\Cat$ gives
\begin{align*}
 \Map(G_{\ell,m},K)
 &\simeq N_mK(\ell)
 \simeq\lim_jN_mK(t_j)
 \simeq\lim_j\Map(G_{t_j,m},K).
\end{align*}
Yoneda therefore gives
$G_{\ell,m}\simeq\colim_{j\in J^{\op}}G_{t_j,m}$,
so $a_r(-,[m])$ is left exact after taking opposites.
\item \textbf{Construct the categorical relations.}
For fixed $t$, the functor
$A\mapsto L_{\mathrm{lex},T_{r-1}}(h_t\times A)$ is the left
adjoint of evaluation at $t$ in $\mathcal L_{T_{r-1}}$;
this is the categorical form of
\eqref{eq:stratified-compact-evaluation}.  It preserves colimits.
The spine colimit
$[m]\simeq[1]\amalg_{[0]}\cdots\amalg_{[0]}[1]$ in $\Cat$
gives \eqref{eq:stratified-arity-Segal}.  Simplicial density gives
$\colim_{\sigma:\Delta[k]\to\mathbf E}[k]\simeq\mathbf E$ in
$\Cat$, and $\mathbf E\simeq[0]$; applying the same left adjoint
gives \eqref{eq:stratified-arity-complete}.  The colimit maps belong
to a single cosimplicial functor, so their face, degeneracy, and
higher coherences agree.
\item \textbf{Identify the generated doctrine.}
By Lemma~\ref{lem:stratified-compact-generators}, every compact object
of $\operatorname{Cat}^r(\TT)$ is a retract of a finite colimit
of the $G_{t,m}$.  Taking opposites proves the finite-limit-and-retract
generation of $T_r$.\qedhere
\end{enumerate}
\end{proof}

The successive categorical arities also yield an explicit image of the
original NU type and operation arities in every $T_r$.

\begin{construction}[Iterated NU categorical arities]
\label{con:stratified-NU-arities}
Set $b_0=\id_{\RNU}$ and define recursively
\begin{equation}\label{eq:stratified-native-arities}
 b_r(t,m_1,\ldots,m_r)
 :=a_r\bigl(b_{r-1}(t,m_1,\ldots,m_{r-1}),[m_r]\bigr).
\end{equation}
Thus the functors $a_r$ construct
\[
 b_r:\RNU\times(\Delta^{\op})^r\longrightarrow T_r.
\]
For each $\mathbf m$, the restriction $b_r(-,\mathbf m)$ is left
exact.  The original NU structural arrows are retained as the marked
arrows of the $\RNU$ source; no arrow in a new minimal doctrine is
declared representable by this construction.
\end{construction}

\begin{theorem}[Iterated NU arity corepresentation]
\label{thm:stratified-native-arities}
Let $r\geq1$ and $X\in\operatorname{Cat}^r(\TT)$.
Write $\widehat X:T_r\to\Sps$ for its theory under
\eqref{eq:stratified-doctrine-recognition}, and let
$X(t)_{\mathbf m}$ denote its evaluation under
\eqref{eq:stratified-NU-semantics}.  There is a natural equivalence
\begin{equation}\label{eq:stratified-native-evaluation}
 \widehat X\bigl(b_r(t,\mathbf m)\bigr)
 \simeq X(t)_{\mathbf m}.
\end{equation}
The image of $b_r$ generates $T_r$ under finite limits and retracts.
For each categorical coordinate, the functor $b_r$ satisfies the
co-Segal and co-completeness comparisons, with limits understood in
$T_r$ and all their mixed simplicial coherences.  For fixed $\mathbf m$,
it carries the finite NU operation and polynomial-arity diagrams to
their left-exact images in $T_r$.
\end{theorem}
\begin{proof}
\leavevmode
\begin{enumerate}
\item \textbf{Corepresent the first categorical direction.}
For $r=1$ the compact-generation equivalence identifies the model
$\widehat X$ at $G_{t,m}^{\op}$ with the mapping space from the
compact object $G_{t,m}$ to $X$.  Thus
\[
 \widehat X(b_1(t,m))
 \simeq\Map_{\operatorname{Cat}(\TT)}(G_{t,m},X)
 \simeq N_mX(t).
\]
This is \eqref{eq:stratified-native-evaluation} for $r=1$.
\item \textbf{Iterate the evaluation equivalence.}
Under
$\operatorname{Cat}^r(\TT)\simeq\Lex(T_{r-1},\Cat)$,
write $X^\mathrm{lex}$ for the categorical diagram corresponding to
$X$.  Its $m_r$th complete-nerve level corresponds, by
$\Lex(T_{r-1},\Sps)\simeq\operatorname{Cat}^{r-1}(\TT)$,
to the $m_r$th outer simplicial level $X_{m_r}$.  Consequently
\begin{align*}
 \widehat X(b_r(t,\mathbf n,m_r))
 &\simeq\Map_{\Cat}([m_r],X^\mathrm{lex}(b_{r-1}(t,\mathbf n)))\\
 &\simeq\widehat{X_{m_r}}(b_{r-1}(t,\mathbf n))\\
 &\simeq X(t)_{\mathbf n,m_r},
\end{align*}
where the last step is the inductive equivalence.  These are
identifications of the complete nerve diagrams and therefore commute
with every simplex operator, not only with their vertices.
\item \textbf{Prove generation.}
For $r=1$ it is
Proposition~\ref{prop:stratified-categorical-arities}.  Inductively,
$b_{r-1}$ generates $T_{r-1}$ under finite limits and retracts,
while $a_r(-,[m])$ preserves both.  Every $a_r(t,[m])$ therefore
belongs to the finite-limit-and-retract closure of the objects
$a_r(b_{r-1}(u,\mathbf n),[m])=b_r(u,\mathbf n,m)$.
Since the $a_r(t,[m])$ generate $T_r$, the latter family does also.
\item \textbf{Construct all categorical relations by Yoneda.}
In any chosen categorical coordinate, the Segal comparison for
$X(t)$ gives, naturally in $X$,
\[
 \widehat X(b_r(t,\ldots,m,\ldots))
 \simeq
 \widehat X(b_r(t,\ldots,1,\ldots))
 \times_{\widehat X(b_r(t,\ldots,0,\ldots))}\cdots
 \times_{\widehat X(b_r(t,\ldots,0,\ldots))}
 \widehat X(b_r(t,\ldots,1,\ldots)).
\]
The compact representables in $\operatorname{Cat}^r(\TT)$ therefore
satisfy the corresponding colimit comparison.  Taking opposites gives
its finite-limit comparison in $T_r$.  For completeness, first take the
colimit of the compact representing objects over the simplex category
of $\mathbf E$ in the presentable category
$\operatorname{Cat}^r(\TT)$.  Mapping this colimit into every $X$
identifies it, by \eqref{eq:stratified-native-evaluation} and
completeness of $X(t)$, with the compact representing object at
coordinate $0$.  Yoneda identifies the colimit with that object.
Taking opposites consequently constructs the required particular
limit in $T_r$, without asserting arbitrary infinite limits there.
All comparisons are natural in the complete multisimplicial diagram,
which supplies their mixed coherence.
\item \textbf{Transport the NU operation diagrams.}
Each $b_r(-,\mathbf m)$ is a composite of the left-exact functors
$a_i(-,[m_i])$.  It therefore carries the finite pullback and
product diagrams of Section~\ref{subsec:nu-doctrine}, the polynomial
base-change squares, and their specified relation cells to coherent
finite-limit diagrams.  This statement concerns the underlying
finite-limit categories: the selected representable arrows and their
dependent-product semantics remain the ones in the NU source and are
not reassigned in the minimal target theory.\qedhere
\end{enumerate}
\end{proof}

\subsection{Iteration of the internal NU classification}
\label{subsec:stratified-internalization}

The arities now describe each semantic stage. To identify its internal
theories, apply the walking-source functor degreewise. Its preservation
of limits ensures that the resulting diagrams still satisfy the Segal
and completeness conditions; the operation calculation then identifies
what those degrees and comparisons mean.

Write $\CNU=\Lex(\RNU,\Cat)\simeq\operatorname{Cat}(\TT)$.
The walking-theory theorem and its limit corollary give
\begin{equation}\label{eq:stratified-internalization-base}
 \Int^{\NU}_\infty(K)
 \simeq\Nat_{\CNU}(\Omega_{\NU},K),
 \qquad
 \Int^{\NU}_\infty:\CNU\longrightarrow\Cat
 \quad\text{preserves bounded small limits}.
\end{equation}
Here $\Nat$ is the full $\infty$-category of ordinary natural
transformations and modifications; the marked lax comparisons are
encoded in $\Omega_{\NU}$.

\begin{theorem}[Stratified NU internalization]
\label{thm:stratified-internalization}
For every $r\geq0$, pointwise application of
\eqref{eq:stratified-internalization-base} constructs a functor
\begin{equation}\label{eq:stratified-internalization}
 \operatorname{Cat}^r(\Int^{\NU}_\infty):
 \operatorname{Cat}^{r+1}(\TT)
 \longrightarrow\operatorname{Cat}^{r+1}(\Sps).
\end{equation}
For $K\in\operatorname{Cat}^r(\CNU)$ its value has
\begin{equation}\label{eq:stratified-internalization-degree}
 \bigl(\operatorname{Cat}^r(\Int^{\NU}_\infty)(K)\bigr)_{\mathbf m}
 \simeq\Int^{\NU}_\infty(K_{\mathbf m}).
\end{equation}
All Segal compositions, completeness comparisons, and mixed higher
coherences are induced by the complete categorical structure of $K$.
\end{theorem}
\begin{proof}
\leavevmode
\begin{enumerate}
\item \textbf{Preserve complete categorical structure.}
Corollary~\ref{thz:cor:limits} proves that
$\Int^{\NU}_\infty$ preserves the limits defining Segal and
walking-isomorphism completeness comparisons.  Applying it to the
entire $r$-simplicial diagram therefore produces a diagram satisfying
both comparisons in each coordinate.  The resulting object lies in
$\operatorname{Cat}^r(\Cat)
 =\operatorname{Cat}^{r+1}(\Sps)$.
\item \textbf{Retain all morphisms and higher coherence.}
A map of $r$-fold category objects is a natural map of the complete
multisimplicial diagrams.  Applying the single functor
$\Int^{\NU}_\infty$ objectwise is functorial on those maps and their
higher simplices.  It commutes with every face, degeneracy, and mixed
simplicial identity.  This gives \eqref{eq:stratified-internalization}
and \eqref{eq:stratified-internalization-degree}, compatibly with
iteration.\qedhere
\end{enumerate}
\end{proof}

We next construct the marked NU operation description of this object
without using its representing higher comparison objects.

\begin{lemma}[Categorical realization of iterated NU arities]
\label{lem:stratified-categorical-arity-realization}
Let $K\in\operatorname{Cat}^r(\CNU)$ and put
$A_r=\operatorname{Cat}^r(\TT)$. In compact-doctrine recognition
use the internal complete-nerve coordinate of $K_{\mathbf m}(t)$
as the outer coordinate of $\operatorname{Cat}(A_r)$.
Denote the resulting left-exact model by
$\widetilde K:T_r\to\Cat$. Then
\begin{equation}\label{eq:stratified-categorical-arity-realization}
 \widetilde K\,b_r\simeq
 \bigl((t,\mathbf m)\longmapsto K_{\mathbf m}(t)\bigr)
\end{equation}
naturally as category-valued diagrams. This construction carries the
NU operation squares and all external operators. At $r=0$ it is the
identification $T_0=\RNU$, $b_0=\id$ and $\widetilde K=K$.
\end{lemma}
\begin{proof}
\leavevmode
\begin{enumerate}
\item \textbf{Move the distinguished nerve coordinate.}
Unpack $K$ as the diagram of spaces
\[
 Z(t;\mathbf m,q)=N_q(K_{\mathbf m}(t)).
\]
For fixed $q$ the entries are left exact in $t$ and complete Segal
in each of the $r$ variables $\mathbf m$: the nerve level $N_q$
preserves limits. By Theorem~\ref{thm:stratified-NU-semantics}
these entries define an object $K^\sharp_q\in A_r$ whose
$\mathbf m$-level at $t$ is $Z(t;\mathbf m,q)$. Restriction in
$q$ gives $K^\sharp:\Delta^{\op}\to A_r$. Its Segal and
completeness comparisons are equivalences since evaluation at
$(t,\mathbf m)$ preserves limits and jointly detects equivalences,
and there it is the complete nerve of $K_{\mathbf m}(t)$.

This is a coordinate permutation on multisimplicial diagrams:
regroup the slots as $(q;\mathbf m)$ and keep the value
$Z(t;\mathbf m,q)$. The inverse puts the $q$-slot back inside.
The defining finite-limit condition in $t$ and the coordinatewise
Segal and completeness conditions are invariant under this
regrouping. It thus gives an equivalence
$\operatorname{Cat}^r(\CNU)\to\operatorname{Cat}(A_r)$ on the
full diagram categories with these properties, including their
morphisms. For the iterated category-object convention compare
\cite[Def.~4.1 and Rem.~4.3]{HaugsengSpans}.

\item \textbf{Evaluate on the compact arities.}
For $r\geq1$ define $\widetilde K$ by applying
\eqref{eq:stratified-doctrine-recognition} to $K^\sharp$.
At $c^{\op}\in T_r$ its $q$th nerve level is
$\Map_{A_r}(c,K^\sharp_q)$ by compact recognition, applied degreewise.
Write $b_r(t,\mathbf m)=c_{t,\mathbf m}^{\op}$. Then
Theorem~\ref{thm:stratified-native-arities} gives
\[
 \begin{aligned}
 N_q\widetilde K(b_r(t,\mathbf m))
 &\simeq \Map_{A_r}(c_{t,\mathbf m},K^\sharp_q)\\
 &\simeq (K^\sharp_q(t))_{\mathbf m}
  = N_q(K_{\mathbf m}(t)).
 \end{aligned}
\]
The equivalences are induced by evaluation and are natural in
$t,\mathbf m,q$ and $K$. Full faithfulness of the complete nerve
therefore gives \eqref{eq:stratified-categorical-arity-realization}.
In particular it identifies both arrows and the specified cells of
every NU square, and all restrictions of those cells in each external
direction. The NU mark is tested on its original $t$-coordinate.

\item \textbf{Check the two directions on a represented NU model.}
For $r=1$ take the constant external category object on the model
$L=BF$ of Proposition~\ref{prop:NU-total-category-obstruction}.
At $t$ with $F(t)$ nontrivial the one-object category $L(t)$ has
only the identity invertible, so $N_0L(t)=*$ and $N_1L(t)=F(t)$.
Consequently $Z(t;1,0)=*$ and $Z(t;0,1)=F(t)$.
The coordinate permutation used above puts the endomorphisms in
the internal nerve direction, as required in
\eqref{eq:stratified-categorical-arity-realization}.
\qedhere
\end{enumerate}
\end{proof}

\begin{corollary}[Stratified recovery of represented theorization]
\label{cor:stratified-represented-recovery}
Let $T$ be small and finitely complete and let $\tau$ be a Berman
marking.  Write $\Thz_\tau(T)^{\mathrm{represented,strong}}$ for
the underlying $\infty$-category of represented theorized objects
and strong morphisms.  For every $r\geq0$, degreewise cocartesian
unstraightening gives a canonical equivalence
\begin{equation}\label{eq:stratified-represented-recovery}
 \operatorname{Cat}^r\!\bigl(
   \Thz_\tau(T)^{\mathrm{represented,strong}}\bigr)
 \simeq
 \operatorname{Cat}^r\!\bigl(\Lex(T,\Cat)\bigr).
\end{equation}
The equivalence preserves the external face, degeneracy, Segal,
completeness, and mixed coherence diagrams.  In particular,
\begin{equation}\label{eq:stratified-NU-represented-recovery}
 \operatorname{Cat}^r\!\bigl(
   \Thz_{\tau_{\NU}}(\RNU)^{\mathrm{represented,strong}}\bigr)
 \simeq \operatorname{Cat}^{r+1}(\TT).
\end{equation}
\end{corollary}
\begin{proof}
\leavevmode
\begin{enumerate}
\item \textbf{Iterate represented recovery.}
Theorem~\ref{thz:thm:recovery} identifies the underlying
$\infty$-categories of represented, strong theorized objects and
left-exact categorical diagrams by mutually inverse straightening and
unstraightening.  Apply $\operatorname{Cat}^r$ to this equivalence.
At every multidegree the construction is the same cocartesian
unstraightening; since it is a functor, the entire simplicial
structure and all higher coherence are transported together.
\item \textbf{Specialize to the NU doctrine.}
Take $T=\RNU$ and $\tau=\tau_{\NU}$ and compose with
$\Lex(\RNU,\Cat)\simeq\operatorname{Cat}(\TT)$ from
Theorem~\ref{thm:nu-categorical}.  This gives
\eqref{eq:stratified-NU-represented-recovery} without changing the
selected representable arrows of $\RNU$.\qedhere
\end{enumerate}
\end{proof}

\subsection{Normalized operations and native stratified records}

The higher categorical operation records can be calculated directly in
each external degree. Their complete string diagrams retain the base
degrees, the structural cocartesian arrows, and all comparison cells.
The following constructions identify these native records with the
internal-object semantics of the stratified NU doctrine.

\begin{lemma}[Higher normalized operation arrows]
\label{lem:stratified-normalized-operations}
Let $T$ be small and finitely complete and let
$K\in\operatorname{Cat}^r(\Lex(T,\Cat))$.  There is a canonically
constructed $r$-fold complete category object
\[
 \mathbb E_T(K)\in\operatorname{Cat}^r(\Cat/T)
 \simeq
 \operatorname{Cat}^r(\Cat)/\delta^r T,
\]
whose value at $\mathbf m$ is the cocartesian unstraightening
$p_{\mathbf m}:E_{\mathbf m}=\int_TK_{\mathbf m}\to T$.
For an $r$-fold category object $A$ in $\Cat$, write
$\operatorname{Ar}_{\mathrm{lev}}(A)$ for the degreewise arrow
object with
$\operatorname{Ar}_{\mathrm{lev}}(A)_{\mathbf m}
 =\Fun([1],A_{\mathbf m})$.
For every arrow $f:s\to t$ in $T$, form the pullback \emph{in}
$\operatorname{Cat}^r(\Cat)$
\begin{equation}\label{eq:stratified-normalized-arrow-category}
 \mathsf{Ar}^{(r)}_f(K):=
 \operatorname{Ar}_{\mathrm{lev}}(\mathbb E_T(K))
 \times_{\delta^r\Fun([1],T)}\delta^r\{f\}.
\end{equation}
Here the map to the base applies $\Fun([1],-)$ degreewise to
$p_{\mathbf m}$, and $\delta^r\{f\}$ is the constant
$r$-fold category object on the point $f\in\Fun([1],T)$.
In particular, \eqref{eq:stratified-normalized-arrow-category} is
\emph{not} the cotensor of $\mathbb E_T(K)$ in the slice $\Cat/T$:
it retains arrows of the specified, possibly nonidentity, degree $f$.
The object $\mathsf{Ar}^{(r)}_f(K)$ has source and target maps to
$K(s)$ and $K(t)$, where $K(s)_{\mathbf m}=K_{\mathbf m}(s)$.
Its fiber in degree $\mathbf m$ at
$x\in K_{\mathbf m}(s)$ and $y\in K_{\mathbf m}(t)$ is
the normalized operation space
\[
 P^{\mathbf m}_f(x,y)
 \simeq\fib_f\bigl(\Map_{E_{\mathbf m}}(x,y)
                  \longrightarrow\Map_T(s,t)\bigr).
\]
The subcategory of cocartesian arrows over $f$ is naturally equivalent
through source evaluation to $K(s)$ as an $r$-fold complete category
object.  The complete string diagrams of the $E_{\mathbf m}$ induce
composition, units, and all higher coherence of these operation arrows,
compatibly with every external categorical operator.  The same holds
for the displayed-operation comparison and pasting maps of
Proposition~\ref{thz:prop:displayed} when its vertical arrows are marked.
\end{lemma}
\begin{proof}
\leavevmode
\begin{enumerate}
\item \textbf{Unstraighten the complete categorical diagram.}
The degreewise objects and maps are precisely the iterated represented
objects and strong morphisms identified in
Corollary~\ref{cor:stratified-represented-recovery}.
The unstraightening equivalence takes limits of functor diagrams over
$T$ to limits of cocartesian fibrations and cocartesian-preserving maps
over $T$.  Theorem~\ref{thm:marked-cocart-limits} with the sharp marking
shows that these limits are created in $\Cat/T$.  Consequently the degreewise
unstraightenings $E_{\mathbf m}$ satisfy the external Segal and
completeness comparisons in $\Cat/T$.  This constructs
$\mathbb E_T(K)$ with all its external simplicial operators.
\item \textbf{Take degreewise arrow fibers and identify their kernels.}
For fixed $f$ the relative lift category
$\mathcal D_f(p_{\mathbf m})$ preserves limits of its slice
argument by the nerve calculation in
Theorem~\ref{thm:marked-cocart-limits}. Applied degreewise it
therefore preserves each external Segal and completeness comparison.
Its defining pullback is exactly
\eqref{eq:stratified-normalized-arrow-category}.  Source and target evaluation land over
$s$ and $t$ and hence define maps to $K(s)$ and $K(t)$.
In degree $\mathbf m$ the pullback is explicitly
\[
 \Fun([1],E_{\mathbf m})
 \times_{\Fun([1],T)}\{f\}.
\]
Its homotopy fiber over $(x,y)$ is the space of arrows $x\to y$
whose degree is identified with $f$, namely the displayed normalized
kernel.  This calculation uses the ordinary total arrow category,
not the category of vertical arrows in the slice.
\item \textbf{Recover universal structural transport.}
Every $p_{\mathbf m}$ is cocartesian.  The source functor from its
category of cocartesian lifts over $f$ is an equivalence: every $x$
has a lift, its space of lifts is contractible, and cocartesian
factorization identifies maps between lifts with maps between their
sources.  This is natural in $K_{\mathbf m}$ and therefore is an
equivalence of complete external diagrams.
\item \textbf{Construct composition and displayed pasting.}
For each $n$, the complete string category
$\Fun([n],E_{\mathbf m})$ simultaneously records composable operation
arrows, their composites, and their higher simplices.  The categorical
spine equivalences and the simplicial operators construct balanced
composition and its entire associativity and unit coherence, exactly
as in \eqref{thz:eq:kernel}.  Their formation is functorial in
$E_{\mathbf m}$, so it commutes with the external categorical
operators.  Cocartesian factorization and the mapping-space pullbacks
of Proposition~\ref{thz:prop:displayed} are equally functorial; its
comparison, conjugation, and pasting constructions therefore assemble
throughout the same external diagrams.\qedhere
\end{enumerate}
\end{proof}

\begin{lemma}[Limits of local marked operation records]
\label{lem:stratified-local-operation-limits}
For a marked simplex $\sigma:(\Delta^n,W_\sigma)\to(T,\tau)$,
let $\mathscr R_\sigma(F)$ be the full category of lifts of the
specified simplex to $\int_TF$, with the edges in $W_\sigma$
required to be cocartesian.  Thus
\[
 \mathscr R_\sigma(F)\subseteq
 \Fun([n],\textstyle\int_TF)
 \times_{\Fun([n],T)}\{\sigma\}.
\]
The functor $\mathscr R_\sigma:\Lex(T,\Cat)\to\Cat$ preserves
bounded small limits.  Consequently applying it in every external
simplicial degree to $K\in\operatorname{Cat}^r(\Lex(T,\Cat))$
constructs a complete $r$-fold category object
$\mathscr R^{(r)}_\sigma(K)$ in $\Cat$.  Its objects and arrows
are explicitly the vertex assignments and composable normalized
operation arrows of Lemma~\ref{lem:stratified-normalized-operations},
with the stipulated marked transport and all their coherent cells.
\end{lemma}
\begin{proof}
\leavevmode
\begin{enumerate}
\item \textbf{Commute the unmarked lift category with limits.}
Let $F=\lim_iF_i$ in $\Lex(T,\Cat)$.  Limits are pointwise,
and unstraightening yields
$\int_TF\simeq\lim_{i,\Cat/T}\int_TF_i$.
The unmarked lift category is $\mathcal D_\sigma$ of this slice
object. The first step of Theorem~\ref{thm:marked-cocart-limits}
proves its limit preservation by representing every complete nerve
level directly in $\Cat/T$.
\item \textbf{Commute the structural locus with limits.}
An arrow over $u:s\to t$ in a cocartesian unstraightening is
cocartesian precisely when its vertical comparison
$F(u)x\to y$ is an equivalence in $F(t)$.  In the pointwise
limit $F(t)=\lim_iF_i(t)$, an arrow is an equivalence if and only
if all its projections are equivalences.  Thus the condition on
every edge of $W_\sigma$ is preserved and detected by the limit.
Taking the full subcategories on this object predicate gives
$\mathscr R_\sigma(F)\simeq\lim_i\mathscr R_\sigma(F_i)$,
including every morphism and higher cell.
\item \textbf{Recover the external category object.}
Every external Segal and completeness comparison of $K$ is a limit
comparison in $\Lex(T,\Cat)$.  Apply the limit-preserving functor
$\mathscr R_\sigma$ to each comparison.  The resulting
multisimplicial diagram is therefore complete Segal in every external
direction.  The lift-category construction uses the complete string
functors $\Fun([n],E_{\mathbf m})$, so its terms are precisely the
normalized operation strings and their stated higher coherence.
\qedhere
\end{enumerate}
\end{proof}

\begin{construction}[Native stratified NU operation records]
\label{con:stratified-native-records}
Choose a quasicategorical presentation
$\phi:Q\xrightarrow{\simeq}\RNU$ as in
Lemma~\ref{thz:lem:marked-resolution}.  Let
$\Sigma_Q=\Delta/Q$ be its category of simplices and give each
simplex $\sigma:\Delta^n\to Q$ the pullback of the positive
marking $\tau_{\NU}$.  For
$K\in\operatorname{Cat}^r(\CNU)$, form the complete
$r$-fold category object of local normalized operation strings
\[
 \mathscr R^{(r)}_{\sigma}(K)
 :=\mathscr R^{(r)}_{\phi\sigma}(K),
\]
using the preceding lemma and the operation arrow categories of
\eqref{eq:stratified-normalized-arrow-category}, applied to the
semantic arity diagram with the internal nerve
coordinate placed outermost, $\widetilde K b_r\simeq((t,\mathbf m)\mapsto K_{\mathbf m}(t))$
of Lemma~\ref{lem:stratified-categorical-arity-realization}.  Restriction of
strings along simplex operators constructs a functor
$\Sigma_Q^{\op}\to\operatorname{Cat}^r(\Cat)$.  Define the
\emph{native stratified NU operation category} by the limit
\begin{equation}\label{eq:stratified-native-sections}
 \mathsf{Int}^{\NU,(r)}_{\mathrm{nat}}(K)
 :=\lim_{\sigma\in\Sigma_Q^{\op}}
          \mathscr R^{(r)}_{\sigma}(K)
 \quad\in\operatorname{Cat}^r(\Cat).
\end{equation}
An object is thus a coherent system of NU vertex objects, operations
over every native NU arrow, and all the higher composition simplices.
For structural arrows its operation is cocartesian.  The seven NU
operation squares, polynomial arities, and comparison cubes occur as
the corresponding named faces of this complete operation diagram.
The defining limit uses the categorical NU arities in $T_r$ and
every simplex of the completed $\RNU$, not just its finite
named-cell shadow.  No higher classifier or
levelwise application of $\Int^{\NU}_\infty$ enters the definition.
\end{construction}

\begin{theorem}[Native stratified NU comparison]
\label{thm:stratified-native-microcosm}
The operation construction
\eqref{eq:stratified-native-sections} is independent, up to its
canonical comparison equivalence, of the quasicategorical presentation
$Q$ and is natural in categorical NU macrocosms and ordinary maps.
For every $r\geq0$ and $K\in\operatorname{Cat}^r(\CNU)$ there are
natural equivalences of full iterated category objects
\begin{equation}\label{eq:stratified-native-microcosm}
 \mathsf{Int}^{\NU,(r)}_{\mathrm{nat}}(K)
 \simeq
 \operatorname{Cat}^r(\Int^{\NU}_\infty)(K)
 \quad\text{in }\operatorname{Cat}^r(\Cat).
\end{equation}
They identify the native normalized NU operations, their structural
cocartesian comparisons, polynomial dependent arities, substitution,
and all mixed categorical composition cells.  In particular, in every
external degree $\mathbf m$ and ordinary morphism degree $q$,
\begin{equation}\label{eq:stratified-native-cotensor}
 N_q\bigl(\mathsf{Int}^{\NU,(r)}_{\mathrm{nat}}(K)_{\mathbf m}\bigr)
 \simeq
 \Map_{\CNU}(\Omega_{\NU},K_{\mathbf m}^{[q]}),
 \qquad
 K_{\mathbf m}^{[q]}(t)=\Fun([q],K_{\mathbf m}(t)).
\end{equation}
The comparison is natural on complete $q$-simplices, not only on
objects or invertible morphisms.
\end{theorem}
\begin{proof}
\leavevmode
\begin{enumerate}
\item \textbf{Descend the constructed operation strings.}
Lemma~\ref{thz:lem:marked-resolution} constructs the marked colimit
\[
 (\RNU,\tau_{\NU})
 \simeq\operatorname*{colim}_{\sigma\in\Sigma_Q}
    (\Delta^{n_\sigma},W_\sigma)
\]
in marked categories over $\RNU$.  Lemma~\ref{lem:stratified-categorical-arity-realization} identifies
the native categorical diagram $\widetilde K b_r$ with the
complete multisimplicial NU diagram of $K$.  In degree
$\mathbf m$, the local entries of
\eqref{eq:stratified-native-sections} are therefore exactly the
full categories of lifts of these marked simplices into the
cocartesian unstraightening of $K_{\mathbf m}$.  Marked descent,
Proposition~\ref{prop:marked-descent}, therefore gives a natural
equivalence of full categories
\[
 \bigl(\mathsf{Int}^{\NU,(r)}_{\mathrm{nat}}(K)\bigr)_{\mathbf m}
 \simeq
 \Fun_{/\RNU}^{\tau_{\NU}\text{-}\coc}
 \bigl(\RNU,\textstyle\int_{\RNU}K_{\mathbf m}\bigr).
\]
This comparison is induced by restriction and gluing of entire
coherent string functors.  It includes arbitrary modifications and
all higher cells.  It also shows that different presentations $Q$
construct equivalent objects, since each maps by the same universal
marked-colimit equivalence to the completed source.
\item \textbf{Identify the classifier and external categorical coherence.}
Theorems~\ref{prop:marked-sections} and
\ref{thz:thm:walking} identify the right-hand side, naturally
in $K_{\mathbf m}$, with
$\Nat_{\CNU}(\Omega_{\NU},K_{\mathbf m})$.
Every map in this chain is natural under the external simplicial
operators of $K$, and both sides satisfy the external complete-Segal
relations by Lemma~\ref{lem:stratified-local-operation-limits} and
Theorem~\ref{thm:stratified-internalization}.  The degreewise
comparisons therefore assemble to the asserted equivalence in
$\operatorname{Cat}^r(\Cat)$.
\item \textbf{Identify the full morphism direction.}
Apply the complete nerve in an arbitrary degree $q$.  The
complete-nerve comparison lemma, Lemma~\ref{lem:complete-nerve-comparison},
gives
\[
 N_q\Nat_{\CNU}(\Omega_{\NU},K_{\mathbf m})
 \simeq
 \Map_{\CNU}(\Omega_{\NU},K_{\mathbf m}^{[q]}).
\]
Pointwise cotensoring by $[q]$ preserves finite limits, so the
target lies in $\CNU$.  Naturality in all $[q]$ and all external
simplicial operators proves \eqref{eq:stratified-native-cotensor},
including the noninvertible comparison arrows.
\item \textbf{Compare NU operation terms and dependent substitution.}
For every named NU arrow $u$, its operation is the corresponding
arrow in the complete string diagram of Lemma~\ref{lem:stratified-normalized-operations}.
Under the descent equivalence it is the section arrow over $u$;
under the comma-classifier equivalence it is the image of the
universal comma arrow
$(d,a(u))\to(d',\id)$.
Lemma~\ref{lem:term} identifies that image with the original
operation comparison, preserving its source and target boundaries.
Its marked strongness is equivalently cocartesianness, as checked in
Theorem~\ref{prop:marked-sections}.
For a named NU polynomial pullback or comparison cube, the local
string record contains its complete coherently commuting diagram.
Left exactness of $K_{\mathbf m}$ identifies its polynomial
input fibers; Proposition~\ref{prop:polynomial-arity} identifies
them with the comprehension fibers.  Proposition~\ref{thz:prop:displayed}
constructs the operation-kernel comparisons and their horizontal,
vertical, and substitution pastings.  Those constructions are natural
in maps of categorical NU models, so Lemma~\ref{lem:stratified-normalized-operations}
assembles them in every external direction.  Since all comparisons
are maps of the complete string diagrams rather than choices of
individual binary composites, their mixed higher coherences are
preserved.\qedhere
\end{enumerate}
\end{proof}

\subsection{Universal higher comparison objects}
\label{subsec:stratified-higher-classifiers}

A classifier for objects should also classify a chain of morphisms
between them. Tensoring the walking source by a finite ordinal gives
that chain, after lex reflection. Free complete-category evaluation then
places it in any external categorical degree. These two adjunctions
produce the full family of higher comparison objects.

For $q\geq0$, define in $\CNU$
\begin{equation}\label{eq:stratified-copower}
 \Omega_{\NU}\langle q\rangle
 :=L_{\mathrm{lex},\RNU}
 \bigl(t\longmapsto\Omega_{\NU}(t)\times[q]\bigr).
\end{equation}
The product on the right is pointwise in $\Fun(\RNU,\Cat)$, before
left-exact reflection.  Enriched reflection and cartesian closure of
$\Cat$ give
\begin{equation}\label{eq:stratified-copower-UP}
 \Map_{\CNU}(\Omega_{\NU}\langle q\rangle,K)
 \simeq
 \Map_{\Cat}([q],\Int^{\NU}_\infty(K))
 =N_q\Int^{\NU}_\infty(K).
\end{equation}
In particular, $\Omega_{\NU}\langle0\rangle\simeq\Omega_{\NU}$.
The construction is cosimplicial in $q$ and represents
noninvertible morphism chains, not merely invertible comparisons.

\begin{proposition}[Co-Segal relations of the NU comparison arities]
\label{prop:stratified-NU-copower-relations}
The cosimplicial family $q\mapsto\Omega_{\NU}\langle q\rangle$
satisfies
\begin{align}
 \Omega_{\NU}\langle q\rangle
 &\simeq
 \underbrace{\Omega_{\NU}\langle1\rangle
 \amalg_{\Omega_{\NU}}\cdots
 \amalg_{\Omega_{\NU}}\Omega_{\NU}\langle1\rangle}
 _{q\text{ copies}}
 &&(q\geq2),\label{eq:stratified-NU-copower-Segal}\\
 \operatorname*{colim}_{\sigma:\Delta[k]\to\mathbf E}
 \Omega_{\NU}\langle k\rangle
 &\simeq\Omega_{\NU}.
 \label{eq:stratified-NU-copower-complete}
\end{align}
The comparisons are natural as maps of the complete cosimplicial
arity diagram.
\end{proposition}
\begin{proof}
\leavevmode
\begin{enumerate}
\item \textbf{Construct the spine colimits.}
For every $K\in\CNU$, mapping the proposed spine coproduct into
$K$ gives the iterated fiber product of
$\Map_{\Cat}([1],\Int^{\NU}_\infty(K))$ over
$\Map_{\Cat}([0],\Int^{\NU}_\infty(K))$.
The complete Segal nerve identifies this with
$\Map_{\Cat}([q],\Int^{\NU}_\infty(K))$.
Equation~\eqref{eq:stratified-copower-UP} and Yoneda give
\eqref{eq:stratified-NU-copower-Segal}.
\item \textbf{Construct co-completeness.}
Simplicial density in $\Cat$ gives
$\colim_{\sigma:\Delta[k]\to\mathbf E}[k]\simeq\mathbf E$,
and the walking-isomorphism groupoid satisfies
$\mathbf E\simeq[0]$.  Mapping the proposed colimit into any $K$
and using \eqref{eq:stratified-copower-UP} therefore gives
\begin{align*}
 \lim_{\sigma:\Delta[k]\to\mathbf E}
 \Map_{\Cat}([k],\Int^{\NU}_\infty(K))
 &\simeq\Map_{\Cat}(\mathbf E,\Int^{\NU}_\infty(K))\\
 &\simeq\Map_{\Cat}([0],\Int^{\NU}_\infty(K))\\
 &\simeq\Map_{\CNU}(\Omega_{\NU},K).
\end{align*}
Yoneda proves \eqref{eq:stratified-NU-copower-complete}.  Both
comparisons arise from the same cosimplicial arity functor, so their
higher simplicial compatibility is inherited rather than separately
specified.\qedhere
\end{enumerate}
\end{proof}

\begin{lemma}[Free complete-category evaluation]
\label{lem:stratified-free-evaluation}
For $r\geq1$ and $\mathbf m=(m_1,\ldots,m_r)$,
evaluation on $\operatorname{Cat}^r(\CNU)$ has a left adjoint
\[
 F^{(r)}_{\mathbf m}:\CNU
 \rightleftarrows\operatorname{Cat}^r(\CNU)
 :\operatorname{ev}_{\mathbf m}.
\]
Let $\Delta[\mathbf m]=\prod_i\Delta[m_i]$ be the external
multisimplicial representable, and write $\odot$ for copower by spaces.
If $L_{\mathrm{CSS}}^{(r)}$ is the simultaneous complete-Segal
reflection of $r$-simplicial diagrams in $\CNU$, then
\begin{equation}\label{eq:stratified-free-evaluation-construction}
 F^{(r)}_{\mathbf m}A
 :=L_{\mathrm{CSS}}^{(r)}(\Delta[\mathbf m]\odot A),
\end{equation}
and
\begin{equation}\label{eq:stratified-free-evaluation-UP}
 \Map_{\operatorname{Cat}^r(\CNU)}
 (F^{(r)}_{\mathbf m}A,K)
 \simeq\Map_{\CNU}(A,K_{\mathbf m}).
\end{equation}
The left adjoints form a coherent multicosimplicial family in
$\mathbf m$.  For $r=0$ set
$F^{(0)}_{()}=\id_{\CNU}$.
\end{lemma}
\begin{proof}
\leavevmode
\begin{enumerate}
\item \textbf{Construct complete-Segal reflection from universal maps.}
Lemma~\ref{lem:stratified-compact-generators} makes $\CNU$
presentable and supplies a small family of compact generators
$G^{\RNU}_{t,n}$.  In
$\Fun((\Delta^{\op})^r,\CNU)$ form the images of each spine
inclusion and the walking-isomorphism completeness map in one
simplicial coordinate, externally multiplied by representables in the
other coordinates and copowered by every member of this generating
family.  Their local objects are exactly the diagrams whose Segal and
completeness comparisons in each direction are equivalences: mapping
from the generators detects equivalences in $\CNU$.  Localization at
this small family exists by the presentable-localization theorem
\cite[\S5.5.4]{Lurie2009}.  It constructs the functor
$L_{\mathrm{CSS}}^{(r)}$ without selecting Segal compositions or
completeness witnesses.
\item \textbf{Compose copowered Yoneda with reflection.}
The unrestricted diagram-category adjunction gives
\[
 \Map_{\Fun((\Delta^{\op})^r,\CNU)}
 (\Delta[\mathbf m]\odot A,K)
 \simeq\Map_{\CNU}(A,K_{\mathbf m}).
\]
Restrict the target to complete-Segal objects and apply the
localization adjunction.  This constructs
\eqref{eq:stratified-free-evaluation-UP}.
\item \textbf{Construct multicosimplicial coherence by mates.}
An operator $\theta:[m]\to[n]$ in any simplicial coordinate gives
a natural transformation of evaluation functors
$\operatorname{ev}_n\to\operatorname{ev}_m$.
Its left mate gives $F_m^{(r)}\to F_n^{(r)}$.
The mate correspondence is functorial on natural transformations and
their higher simplices.  Applied in all coordinates, it constructs
the multicosimplicial family and all mixed interchange coherence.
\qedhere
\end{enumerate}
\end{proof}

\begin{theorem}[Higher classifying objects and all their relations]
\label{thm:stratified-higher-classifiers}
For $r\geq0$, $\mathbf m\in\Delta^r$, and $q\geq0$, construct
\begin{equation}\label{eq:stratified-higher-classifier}
 \Omega^{(r)}_{\mathbf m,q}
 :=F^{(r)}_{\mathbf m}(\Omega_{\NU}\langle q\rangle)
 \in\operatorname{Cat}^{r+1}(\TT).
\end{equation}
There are natural equivalences
\begin{equation}\label{eq:stratified-higher-classifier-UP}
 \Map_{\operatorname{Cat}^{r+1}(\TT)}
 (\Omega^{(r)}_{\mathbf m,q},K)
 \simeq
 N_q\!\left(\mathsf{Int}^{\NU,(r)}_{\mathrm{nat}}(K)_{\mathbf m}\right).
\end{equation}
The objects assemble into a coherent $(r+1)$-cosimplicial functor
\begin{equation}\label{eq:stratified-higher-cosimplicial}
 \boldsymbol\Omega^{(r)}_{\NU}:
 \Delta^{r+1}\longrightarrow\operatorname{Cat}^{r+1}(\TT).
\end{equation}
In each of its $r+1$ categorical coordinates it satisfies the
co-Segal and co-completeness relations.  With other indices fixed,
for $m\geq2$ these include
\begin{align}
 \Omega^{(r)}_{(\ldots,m,\ldots),q}
 &\simeq
 \underbrace{\Omega^{(r)}_{(\ldots,1,\ldots),q}
 \amalg_{\Omega^{(r)}_{(\ldots,0,\ldots),q}}\cdots
 \amalg_{\Omega^{(r)}_{(\ldots,0,\ldots),q}}
 \Omega^{(r)}_{(\ldots,1,\ldots),q}}_{m\text{ copies}},
 \label{eq:stratified-higher-coSegal}\\
 \operatorname*{colim}_{\sigma:\Delta[k]\to\mathbf E}
 \Omega^{(r)}_{(\ldots,k,\ldots),q}
 &\simeq\Omega^{(r)}_{(\ldots,0,\ldots),q}.
 \label{eq:stratified-higher-coComplete}
\end{align}
The same two relations hold in the $q$-direction.  In particular,
these classifiers retain all ordinary morphisms of internal NU
theories and the higher categorical directions with their mixed
coherences.
\end{theorem}
\begin{proof}
\leavevmode
\begin{enumerate}
\item \textbf{Represent every degree and every morphism chain.}
Apply \eqref{eq:stratified-free-evaluation-UP} to
$A=\Omega_{\NU}\langle q\rangle$, then
\eqref{eq:stratified-copower-UP} and
Theorem~\ref{thm:stratified-native-microcosm}:
\begin{align*}
 \Map(\Omega^{(r)}_{\mathbf m,q},K)
 &\simeq\Map_{\CNU}(\Omega_{\NU}\langle q\rangle,K_{\mathbf m})\\
 &\simeq N_q\Int^{\NU}_\infty(K_{\mathbf m})\\
 &\simeq N_q\mathsf{Int}^{\NU,(r)}_{\mathrm{nat}}(K)_{\mathbf m}.
\end{align*}
This is \eqref{eq:stratified-higher-classifier-UP} with all
parameter naturality.
\item \textbf{Prove co-Segal in the external categorical directions.}
Mapping the proposed spine coproduct into any $K$ gives the
iterated fiber product of the represented complete-nerve levels.
Theorem~\ref{thm:stratified-native-microcosm} says that the external
multisimplicial category of native records is Segal in each coordinate.
Applying the $q$th complete-nerve level preserves these limits.
Thus the resulting mapping-space comparison is an equivalence for
every $K$, and Yoneda gives
\eqref{eq:stratified-higher-coSegal}.
\item \textbf{Prove external co-completeness.}
Mapping the left side of
\eqref{eq:stratified-higher-coComplete} into any $K$ gives the limit
of the $q$th complete-nerve levels of the native record categories
over the simplex category of $\mathbf E$.  Completeness of the
external categorical direction, followed by the limit-preserving
functor $N_q$, identifies this with the represented degree at $0$.
Yoneda gives \eqref{eq:stratified-higher-coComplete}.
\item \textbf{Prove the relations in the internal-morphism direction.}
The left adjoint $F^{(r)}_{\mathbf m}$ preserves colimits.
Applying it to
\eqref{eq:stratified-NU-copower-Segal} and
\eqref{eq:stratified-NU-copower-complete} gives respectively
\[
 \Omega^{(r)}_{\mathbf m,q}
 \simeq\underbrace{
 \Omega^{(r)}_{\mathbf m,1}
 \amalg_{\Omega^{(r)}_{\mathbf m,0}}\cdots
 \amalg_{\Omega^{(r)}_{\mathbf m,0}}
 \Omega^{(r)}_{\mathbf m,1}}_{q\text{ copies}}
\]
for $q\geq2$, and
\[
 \operatorname*{colim}_{\sigma:\Delta[k]\to\mathbf E}
 \Omega^{(r)}_{\mathbf m,k}
 \simeq\Omega^{(r)}_{\mathbf m,0}.
\]
Thus every one of the $r+1$ directions has its relations.
\item \textbf{Assemble all higher coherence.}
The cosimplicial copower family and the multicosimplicial left
adjoints are constructed functorially.  Their composite is therefore
\eqref{eq:stratified-higher-cosimplicial}.  Every relation above is
the image of a canonical map of cosimplicial diagrams, or is detected
naturally by Yoneda.  Hence their face, degeneracy, and mixed
coherence relations agree on all higher simplices.\qedhere
\end{enumerate}
\end{proof}

\begin{corollary}[Explicit reflected-comma formula]
\label{cor:stratified-higher-explicit}
For the NU marked comma diagram
$R_{\id_{\RNU},\tau_{\NU}}(t)
 =((\RNU)_{/t})[W_t^{-1}]$,
the higher classifying objects have the construction
\begin{equation}\label{eq:stratified-higher-explicit}
 \Omega^{(r)}_{\mathbf m,q}
 \simeq
 L_{\mathrm{CSS}}^{(r)}\!\left(
 \Delta[\mathbf m]\odot
 L_{\mathrm{lex},\RNU}
 \bigl(t\mapsto R_{\id_{\RNU},\tau_{\NU}}(t)\times[q]\bigr)
 \right).
\end{equation}
No arity, comparison term, or higher coherence is adjoined outside the
comma localization and the two displayed reflections.
\end{corollary}
\begin{proof}
\leavevmode
\begin{enumerate}
\item \textbf{Commute the reflected copower with the first reflection.}
Theorem~\ref{thz:thm:walking} gives
$\Omega_{\NU}=L_{\mathrm{lex},\RNU}R_{\id,\tau_{\NU}}$.
For every left-exact $K$, enriched reflection and the tensor--cotensor
adjunction identify both
$L_{\mathrm{lex}}((L_{\mathrm{lex}}R)\times[q])$
and $L_{\mathrm{lex}}(R\times[q])$ as representing
$\Map_{\Cat}([q],\Nat(R,K))$.  Yoneda therefore identifies them.
Substitute this equivalence into
\eqref{eq:stratified-free-evaluation-construction} and
\eqref{eq:stratified-higher-classifier}.\qedhere
\end{enumerate}
\end{proof}

\subsection{Structural strongness and the walking strata}
\label{subsec:stratified-strongness}

The categorical Segal and completeness comparisons are relations among
the representing arities, not additional marked-strongness conditions
on a section.  The following computation distinguishes these two
constructions even on a complete categorical diagram.

\begin{example}[Categorical spine arrows are not structural strongness]
\label{ex:stratified-spine-strongness}
Let $C:\Delta^{\op}\to\Cat$ be the constant diagram on the
ordinary category $\mathbf{FinSet}$ of finite sets (with a small
skeleton understood).  It is a complete
category object in $\Cat$: every Segal comparison and the
walking-isomorphism completeness comparison is an identity, up to the
canonical equivalences for a constant diagram.

An unmarked lax cone on $C$ is exactly a functor
$X:\Delta^{\op}\to\mathbf{FinSet}$, with the comparison along
$u:[m]\to[n]$ given by $X(u):X_m\to X_n$.  Take the representable
simplicial finite set
\[
 X_m=\operatorname{Hom}_{\Delta}([m],[1]).
\]
It gives an actual coherent unmarked cone, with all face, degeneracy,
and composition comparisons.  Yet $|X_2|=4$ and $|X_1|=3$.
Each of the two spine restrictions $X_2\to X_1$ is therefore
noninvertible in $\mathbf{FinSet}$.  Requiring those arrows to be
marked-strong excludes this valid unmarked cone, although $C$ already
satisfies every complete-Segal relation.  Thus complete-Segal
relations constrain the ambient categorical diagram, whereas marking
simplicial operators imposes additional invertibility on a selected
section.  In the constructions above, categorical constraints are
carried by arity relations and complete-Segal reflection, while NU
structural strongness remains exactly the original $\tau_{\NU}$
condition in \eqref{eq:stratified-native-sections}.
\end{example}

\begin{theorem}[Walking strata and the complete higher classifier]
\label{thm:stratified-walking-strata}
For every $r\geq0$, the zero-degree free categorical evaluation
of the walking NU theory is the iterated discrete object
\[
 \mathbb W_r
 :=F^{(r)}_{(0,\ldots,0)}(\Omega_{\NU})
 \simeq\delta^r\Omega_{\NU}
 \quad\text{in }\operatorname{Cat}^r(\CNU).
\]
It represents the space of zero-degree internal NU theories:
\begin{equation}\label{eq:stratified-walking-zero}
 \Map(\mathbb W_r,K)
 \simeq N_0\mathsf{Int}^{\NU,(r)}_{\mathrm{nat}}(K)_{(0,\ldots,0)}.
\end{equation}
The complete system of higher degrees and their ordinary
noninvertible comparison arrows is represented by the single coherent
multicosimplicial functor
$\boldsymbol\Omega^{(r)}_{\NU}$ of
\eqref{eq:stratified-higher-cosimplicial}, and specializes to the
native NU records through
\eqref{eq:stratified-higher-classifier-UP}.
\end{theorem}
\begin{proof}
\leavevmode
\begin{enumerate}
\item \textbf{Identify the free zero-degree object.}
The representable $\Delta[0]$ is terminal as a simplicial space, so
$\Delta[(0,\ldots,0)]\odot\Omega_{\NU}$ is the constant
$r$-simplicial diagram on $\Omega_{\NU}$.  This diagram already
satisfies every Segal and completeness comparison, hence its
complete-Segal reflection is itself.  This constructs
$\mathbb W_r\simeq\delta^r\Omega_{\NU}$.
\item \textbf{Compute its classification property and the higher strata.}
Free evaluation and the walking-theory equivalence give
\begin{align*}
 \Map(\mathbb W_r,K)
 &\simeq\Map_{\CNU}(\Omega_{\NU},K_{(0,\ldots,0)})\\
 &\simeq N_0\Int^{\NU}_\infty(K_{(0,\ldots,0)})\\
 &\simeq N_0\mathsf{Int}^{\NU,(r)}_{\mathrm{nat}}(K)_{(0,\ldots,0)}.
\end{align*}
Theorem~\ref{thm:stratified-higher-classifiers} proves the
classification and all co-Segal and co-completeness relations for
$\boldsymbol\Omega^{(r)}_{\NU}$; the native equivalence is
Theorem~\ref{thm:stratified-native-microcosm}.  Together these
identify every higher categorical degree and its comparison terms.
\qedhere
\end{enumerate}
\end{proof}

The first categorical direction is the ordinary morphism direction of
internal NU theories, represented by
$\Omega_{\NU}\langle q\rangle$.  The additional directions are free
complete-category evaluations of that same constructed source.  Their universal relations and direct marked-section semantics have now
been constructed. Section~\ref{subsec:stratified-relative} below
supplies their relative realization over the indexing base introduced
in Section~\ref{sec:relative}. These are successive complete
categorical directions of the represented semantics. The higher
relative tensor products and the substitution product for independent
composition multikernels were constructed in
Section~\ref{sec:higher-composition}.

\subsection{Stratified relative classification}
\label{subsec:stratified-relative}

The higher classifiers constructed above and the relative realization
of Section~\ref{sec:relative} address independent parts of the same
problem: one retains categorical degrees, and the other retains the
indexing base. Since the realization preserves limits, they combine
degreewise to classify higher marked cones over an arbitrary doctrine.

The construction also applies to any marked source of the main
relative classification theorem, without passing to the completed
NU source or changing the marking.

\begin{theorem}[Universal relative stratified classifier]
\label{thm:stratified-relative-classifier}
Let $a:D^W\to T$ be a bounded marked source over a small finitely
complete $T$, and put
$\mathcal L_T=\Lex(T,\Cat)$ and
$\Omega_{a,W}=\Env_T(a,W)$.
For $K\in\operatorname{Cat}^r(\mathcal L_T)$ define
\[
 \mathsf S^{(r)}_{a,W}(K)_{\mathbf m}
 :=\Fun^{W\text{-}\coc}_{/T}
 \left(D,\int_TK_{\mathbf m}\right).
\]
These entries form an $r$-fold complete category object in $\Cat$.
Writing
\[
 \Omega_{a,W}\langle q\rangle
 =L_{\mathrm{lex},T}
 \bigl(t\mapsto\Omega_{a,W}(t)\times[q]\bigr),
 \qquad
 \Omega^{(r)}_{a,W;\mathbf m,q}
 =F^{(r)}_{\mathbf m,T}(\Omega_{a,W}\langle q\rangle),
\]
where $F^{(r)}_{\mathbf m,T}$ is the free complete-category evaluation
constructed as in Lemma~\ref{lem:stratified-free-evaluation} with
$\CNU$ replaced by $\mathcal L_T$, let
$\J_T^{(r)}:=\operatorname{Cat}^r(\J_T)$.
Theorem~\ref{thm:J}, together with the iterated slice
comparison established below, identifies its target with
$\operatorname{Cat}^{r+1}(\TT)/\delta^r\B_T$.
There are natural equivalences of mapping spaces
\begin{equation}\label{eq:stratified-relative-classifier}
\begin{aligned}
 N_q\mathsf S^{(r)}_{a,W}(K)_{\mathbf m}
 &\simeq
 \Map_{\operatorname{Cat}^r(\mathcal L_T)}
 (\Omega^{(r)}_{a,W;\mathbf m,q},K)\\
 &\simeq
 \Map_{\operatorname{Cat}^{r+1}(\TT)/\delta^r\B_T}
 \bigl(\J_T^{(r)}\Omega^{(r)}_{a,W;\mathbf m,q},
       \J_T^{(r)}K\bigr).
\end{aligned}
\end{equation}
These equivalences are natural in every external simplicial index
$\mathbf m$ and in the complete internal-morphism index $q$, so
they classify the entire mixed higher comparison diagram, including
noninvertible morphisms.  The native classifiers satisfy the co-Segal and co-completeness
relations in all $r+1$ directions. The construction following the
proof gives classifiers with these colimit relations in the entire
relative NU slice and a natural comparison on realized targets.  In particular,
the induced relative NU realization
\begin{equation}\label{eq:stratified-relative-NU-embedding}
 \J_T^{(r)}:
 \operatorname{Cat}^r(\mathcal L_T)
 \longrightarrow
 \operatorname{Cat}^{r+1}(\TT)/\delta^r\B_T
\end{equation}
is fully faithful and retains the same specified relative base
$\delta^r\B_T$ throughout the iteration.
\end{theorem}
\begin{proof}
\leavevmode
\begin{enumerate}
\item \textbf{Construct the native relative record and its higher degrees.}
Theorems~\ref{thz:thm:envelope} and
\ref{prop:marked-sections} give, naturally in $K_{\mathbf m}$,
\[
 \mathsf S^{(r)}_{a,W}(K)_{\mathbf m}
 \simeq\Nat_{\mathcal L_T}(\Omega_{a,W},K_{\mathbf m}).
\]
The enriched hom preserves limits in the target, by the same
complete-nerve argument as Corollary~\ref{thz:cor:limits}.  Thus the
entries assemble into an $r$-fold complete category object.  Applying
pointwise cotensors, enriched reflection, and free complete-category
evaluation gives
\begin{align*}
 N_q\mathsf S^{(r)}_{a,W}(K)_{\mathbf m}
 &\simeq\Map_{\mathcal L_T}
 (\Omega_{a,W}\langle q\rangle,K_{\mathbf m})\\
 &\simeq\Map_{\operatorname{Cat}^r(\mathcal L_T)}
 (F^{(r)}_{\mathbf m,T}\Omega_{a,W}\langle q\rangle,K).
\end{align*}
This proves the first equivalence in
\eqref{eq:stratified-relative-classifier}.  The proof of
Theorem~\ref{thm:stratified-higher-classifiers}, applied to these
adjunctions, gives all stated categorical relations.
\item \textbf{Iterate the relative element realization.}
Theorem~\ref{thm:J} constructs an enriched fully faithful,
limit-preserving functor
$\J_T:\mathcal L_T\to\CNU/\B_T$.
Applying it degreewise to complete category objects gives the
fully faithful functor $\J_T^{(r)}$ of
\eqref{eq:stratified-relative-NU-embedding}, initially with target
$\operatorname{Cat}^r(\CNU/\B_T)$.  In particular, for
$A,K\in\operatorname{Cat}^r(\mathcal L_T)$ it induces an equivalence
\[
 \Map_{\operatorname{Cat}^r(\mathcal L_T)}(A,K)
 \xrightarrow{\simeq}
 \Map_{\operatorname{Cat}^r(\CNU/\B_T)}
       (\J_T^{(r)}A,\J_T^{(r)}K).
\]
This supplies the second equivalence of
\eqref{eq:stratified-relative-classifier}, compatibly with every
multicosimplicial operator.
\item \textbf{Identify the iterated slice target.}
The forgetful functor from a slice creates limits indexed by a
weakly contractible category. Indeed, the structure maps of the
diagram induce a map from its underlying limit to the limit of the
constant diagram at the base, and the latter is the base itself
for such an indexing category. This verifies the slice universal
property on mapping spaces.

The indexing categories for the Segal pullbacks are weakly
contractible. The completeness comparison uses the walking
isomorphism; its nerve, and hence its simplex category, is
contractible. Thus both the Segal and completeness conditions on
a simplicial object in $\CNU/\B_T$ can be tested in $\CNU$.
The compatible maps to $\B_T$ form a map of simplicial objects to
the constant object $\delta\B_T$. Consequently a complete category
object in $\CNU/\B_T$ is equivalently a complete category object
in $\CNU$ over $\delta\B_T$. Iterating gives
\[
 \operatorname{Cat}^r(\CNU/\B_T)
 \simeq\operatorname{Cat}^{r+1}(\TT)/\delta^r\B_T.
\]
Degreewise straightening, the inverse on the essential image of
$\J_T$, reconstructs the relative categorical model and all
simplicial operators, proving full faithfulness and the specified
base compatibility in \eqref{eq:stratified-relative-NU-embedding}.
\qedhere
\end{enumerate}
\end{proof}

For an arbitrary bounded marked diagram $C:D\to\Cat$, the free
finite-limit extension $\overline C:QD\to\Cat$ of
Section~\ref{subsec:free-completion} supplies the relative doctrine
$T=QD$ and native source $i_D:D^W\to QD$.  Its iterated discrete
categorification $\delta^r\overline C$ is a canonical input to
Theorem~\ref{thm:stratified-relative-classifier} at every finite
stage, and the same theorem applies to every $r$-fold categorical
model over $QD$.  As in Theorem~\ref{thm:main}, the source marking
remains $W$; the free completion introduces no additional strong
comparison directions.

\paragraph{Higher classifiers for arbitrary targets in the relative slice.}
Put $\mathcal B=\CNU/\B_T$ and $Z=\J_T\Omega_{a,W}$. For $Z\to\B_T$
write $UZ$ for its underlying object in $\CNU$. In $\mathcal B$ form
\[
 Z\langle q\rangle_{\mathcal B}
 =\left(L_{\mathrm{lex},\RNU}(t\mapsto UZ(t)\times[q])
              \longrightarrow UZ\longrightarrow\B_T\right).
\]
Here $U$ denotes the forgetful functor of the slice, and the first map
is adjoint to projection. Its mapping space into
$Y\to\B_T$ is the homotopy fiber of the ambient mapping-space map
at the composite to $\B_T$. By the copower adjunction this is
$N_q\Homc_{\mathcal B}(Z,Y)$: in $\CNU$ the relevant chains of
transformations project to the constant chain at $UZ\to\B_T$ with
a specified identification. Equivalently the slice cotensor is
$\B_T\times_{\B_T^{[q]}} (UY)^{[q]}$ with the constant-chain map
from $\B_T$.

The slice is presentable. Apply the construction of
Lemma~\ref{lem:stratified-free-evaluation} there, using any small
generating family to detect its complete-Segal local objects. Set
\[
 D_{\mathbf m,q}=F^{(r)}_{\mathbf m,\mathcal B}
       (Z\langle q\rangle_{\mathcal B}).
\]
For every $Y\in\operatorname{Cat}^r(\mathcal B)$ the two adjunctions give
\[
 \Map(D_{\mathbf m,q},Y)\simeq N_q\Homc_{\mathcal B}(Z,Y_{\mathbf m}).
\]
The right side is complete Segal in $q$, and in each external
coordinate because the enriched hom in the target preserves limits
(test with the slice copowers at every nerve level). Applying Yoneda
as in Theorem~\ref{thm:stratified-higher-classifiers} gives all the
co-Segal pushouts and walking-isomorphism colimits in the entire slice.
All operators are induced by the same adjunctions. On realized targets
$\J_T^{(r)}K$ the right side agrees with
\eqref{eq:stratified-relative-classifier} by enriched full faithfulness.
Evaluating this agreement at the identity of the native classifier
gives the natural comparison
$D_{\mathbf m,q}\to\J_T^{(r)}\Omega^{(r)}_{a,W;\mathbf m,q}$.
Precomposition is an equivalence on mapping spaces into every
realized target.

\emph{The two copowers can differ.} Take $r=0$, $T=\mathrm{FinSet}$,
and the empty marked source. Its native envelope is initial in $\LT$.
It is $h_1=\Map_T(1,-)$ regarded as category-valued: by copowered
Yoneda maps into $K$ form $K(1)^\core\simeq *$. All native ordinal
copowers of it are also initial by their mapping property (the full
hom into $K$ is $K(1)\simeq *$). Its realization $Z$ is the discrete
category object on $V_{\minrep}\to T_{\minrep}$,
where $V=T_{1/}$ is the category of pointed finite sets. Test in
the slice against the projection $Y=\B_T\times W\to\B_T$, where
$W_n=\Fun([n],V)_{\minrep}$. These levels form a complete category
object of $\TT$: limits of minimal theories are formed on their
underlying categories, the spine glues composable strings, and
evaluation at a vertex of the walking isomorphism is an equivalence
of functor categories. In degree $n$, the enriched hom from $Z$ to $Y$
therefore has nerve
\[
 \Map_{\TT}(V_{\minrep},\Fun([n],V)_{\minrep})
       \simeq \Fun([n],\Lex(V,V))^\core.
\]
Indeed the product slice removes the base coordinate, maps from a
discrete category object evaluate at level zero, and preservation of
finite limits into a functor category is pointwise. The displayed
formula applied to the ordinal cotensor gives the asserted degree.
The natural transformation $\id_V\to\mathrm{const}_*$ is not invertible
(on a pointed two-element set its component is not an isomorphism).
For $q=1$ precomposition with $D_{(),1}\to Z$ is the identity-arrow
degeneracy on this complete nerve, which misses that component.
Thus the comparison of ambient and realized native copowers actually
fails to be an equivalence in this example.

\begin{thebibliography}{CHack22}
\raggedright
% The twisted-arrow end is not required for the graded hom proof;
% full hom categories are computed by the slice tensor--cotensor formula.
\bibitem[Ber24]{Berman2020}
J.~D. Berman, \textbf{On lax limits in $\infty$-categories},
Proceedings of the American Mathematical Society 152 (2024), no.~12,
5055--5066. doi:10.1090/proc/16968. arXiv:2006.10851.

\bibitem[Blo24]{Blom2024}
T.~Blom, \textbf{On the straightening of every functor},
\href{https://arxiv.org/abs/2408.16539}{arXiv:2408.16539}.

\bibitem[Hau18]{HaugsengSpans}
R.~Haugseng, \textbf{Iterated spans and classical topological field theories},
Mathematische Zeitschrift 289 (2018), 1427--1488.
doi:10.1007/s00209-017-2005-x.

\bibitem[HHLN23]{HHLN2023}
R.~Haugseng, F.~Hebestreit, S.~Linskens, and J.~Nuiten,
\textbf{Lax monoidal adjunctions, two-variable fibrations and the calculus
of mates}, Proceedings of the London Mathematical Society
127 (2023), 889--957. doi:10.1112/plms.12548.

\bibitem[JT07]{JoyalTierney2007}
A.~Joyal and M.~Tierney, \textbf{Quasi-categories vs Segal spaces},
Contemporary Mathematics 431 (2007), 277--326.
doi:10.1090/conm/431/08278.

\bibitem[Ker]{Kerodon}
J.~Lurie, \textbf{Kerodon}, online reference, \url{https://kerodon.net}
(accessed 11 October 2026).

\bibitem[Kos21]{Kositsyn2021}
R.~Kositsyn, \textbf{Completeness for monads and theories},
\href{https://arxiv.org/abs/2104.00367}{arXiv:2104.00367v2} (2021).

\bibitem[Lur09]{Lurie2009}
J.~Lurie, \textbf{Higher Topos Theory}, Annals of Mathematics Studies
170, Princeton University Press (2009).

\bibitem[LurHA]{LurieHA}
J.~Lurie, \textbf{Higher Algebra}, manuscript (2017).
\url{https://www.math.ias.edu/~lurie/papers/HA.pdf}.

\bibitem[Mat17]{Matsuoka2017}
T.~Matsuoka, \textbf{Higher theories of algebraic structures},
arXiv:1601.00301v3 (2017).

\bibitem[NU25]{NguyenUemura2025}
H.~K. Nguyen and T.~Uemura, \textbf{$\infty$-type theories},
Higher Structures 9(1) (2025), 179--226.
doi:10.21136/HS.2025.04.

\end{thebibliography}
\end{document}
